\quellseite{65}
ist nun Kains Opfer vom Korn auf dem Felde gewesen, und so zugegangen, wie Moses schreibt im dritten Buche [Cap. 2, 1. ff.], daß man sollte versengte und gedörrte Aehren nehmen und zerstößen, und auf den Altar legen, und Gott verbrennen, denn er sagt viel von dem Mehlopfer, gebacken, geröstet, gestoßen, gedörrt etc. Das alles heißt auf hebräisch Mincha, Speisopfer, das man essen mag, als Brod. Solches hat der Kain geopfert von dem Seinen, wie es seine Nahrung gab.

14. Desgleichen opferte auch Abel von seiner Nahrung, nämlich von\footnote{„von“ fehlt in der Erlanger.} den Schafen, und sonderlich von den Erstlingen oder Frühlingen, und von dem Fetten. Davon steht auch im dritten Buch [Mosis, Cap. 3, 9. 10.], wie man das Fette davon nehmen soll, den Schwanz mit dem Rücken, und das Fett, das das Eingeweide bedeckt, sammt allem Fetten, das inwendig ist, und die Nieren mit dem Fette, das daran ist an den Lenden, und das Netz um die Leber an den Nieren dazu, und soll das Netz um die Nieren herab scheiden, und soll es auf den Altar legen und anzünden. So hat hier Abel auch gethan, und ist die Weise von ihm hergenommen.

15. Nun fragt man hier: Welch Opfer unter diesen das beste gewesen sei? Die Juden sagen, daß Kain darum verworfen sei, daß er zu gering Opfer gethan hat; wollen die Person verdammen und rechtfertigen nach den Werken, wie man pflegt. Aber laßt uns auf den Text sehen, [so] werden wir finden, wie er wider die Werke auf den Glauben bringt.

16. Aufs erste soll man nicht dafür halten, daß Kain ein geringer Opfer gethan habe denn Abel, denn es sind die besten Früchte gewesen. Und wenn man die Opfer in Mose ansieht, sind die Speisopfer fast die edelsten unter allen Opfern, bedeuten aber das lautere Evangelium. Derhalben gibt’s [der Text] nicht, daß er um des geringsten Opfers willen sollte verworfen sein, denn das Werk des Opfers halben gilt gleich vor Gott, es sei gering oder köstlich.

17. Zum andern ist auch nicht zu halten, daß Abel Gott gefallen habe eben des Opfers halben, sondern um des Glaubens willen; wie Kain nicht seines Werkes, sondern seines Unglaubens halben verdammt ist. Das schließt sich gewaltig aus dem Text, da er sagt: „Gott hielt sich zu Abel und zu seinem Opfer; aber zu Kain und seinem Opfer hielt er sich nicht.“ Denn Moses will so sagen, daß Gott zum ersten auf die Person oder den Mann sieht, und hernach auf das Werk, das er thut; und nicht wiederum. Darum muß der Mensch zuvor fromm sein und Gott gefallen, ehe er ein gut Werk thun kann; sonst hätte er so müssen reden: Gott hielt sich freundlich zum Opfer, und darnach zu Abel etc. Er will niemand ansehen um seines Werks willen, sondern allein um der Person willen: ist die gut, so gefällt ihm das Werk; ist sie böse, so ist das Werk auch nicht gut. Denn sie beide, nach dem äußerlichen Ansehen, einerlei Werk thun; [es opfert] jeglicher von dem Seinen, und ist je eines so gut als das andere, ja, des Kains Opfer ist herrlicher und scheinbarer. Derhalben müßte auch Gott, wo er nach den Werken richtete, entweder Kains Opfer auch für gut ansehen, oder Abels auch mit verdammen, daß also einerlei Urtheil über beide einerlei Werk gestellt würde.

18. Aber siehe, was er für ein Richter ist. Das Werk, das am scheinbarlichsten, hübschesten und besten ist, verdammt er, das andere lobt er. Deß hätte sich weder Kain noch Abel, Adam noch Eva versehen. Kain meinte, er hätte es so gewiß, daß ihm nicht könnte fehlen, daß sein Opfer Gott gefallen würde, dazu viel mehr gefallen denn des Bruders. Dachte also: Ich bin je der erste Sohn, mir gebührt vor Gott das Priesterthum, dazu das Regiment, bringe auch das beste Opfer; darum wird Gott je mein Opfer ansehen, und ihm besser gefallen lassen denn Abels, meines Bruders. Wiederum dachte Abel also: Ich bin der Geringste, er ist der Beste, hat viel Vortheil vor mir, hat auch das beste Opfer gebracht, darum muß Gott mein Opfer aus lauter Gnaden ansehen. So richtet nun Gott auf dieser Seite, sieht die geringe Person und Opfer an, die größte Person und beste Opfer verdammt er.

19. So haben wir nun ein gewisses, göttliches Urtheil, daß man nicht nach Werken richten soll, und nur sich hüten vor hübschen, gleißenden Werken; denn je hübscher, je fährlicher es ist. Es verdammt nicht die groben Stücke, sondern gemeiniglich die höchsten Werke. Darum ist der Text wohl und stark zu fassen, daß wir uns nicht mit etwas Anderem lassen die Augen blenden. Wenn dieser Text nicht stößt, so weiß ich nicht,
\quellseite{66}
was stoßen soll. So siehe nun und urtheile, ob es genug sei, daß unsere Werktheiligen sagen: Willst du GOtt einen großen Dienst thun, so werde ein\footnote{„ein“ fehlt in der Erlanger.} Pfaff, Mönch oder Nonne, stifte und halte viel Messen und Gottesdienst; meinen, es sei allerdings genug, wenn es nur den Namen hat, und scheint als Gottesdienst und gut Werk. Wenn sie denn nur hören, so platzen und plumpen sie darauf, als wären sie blind, toll und thöricht. Hier ist auch ein Gottesdienst, so herrlich, als man immer thun kann; noch will er [GOtt] sein nicht. Darum muß man weiter fahren, und endlich nach der Version richten, daß das Gottesdienst sei, so aus dem Glauben kommt. Die Version muß zuvor fromm sein und GOtt gefallen, darnach wird der Dienst auch angenehm.

20. Wodurch ist nun Abel fromm worden? Ohne Zweifel durch den Glauben, wie wir; denn er auf keine andere Weise hat mögen rechtfertig werden, denn alle anderen Leute. So schreibt auch von ihm die Epistel an die Hebräer [Cap. 11, 4.]: „Durch den Glauben hat Abel GOtt ein besser Opfer gethan, denn Kain, durch welchen er Zeugniß überkommen hat, daß er gerecht sei, als GOtt zeugete über seine Gabe; und durch denselben redet er noch, wiewohl er gestorben ist“, das ist, dies Exempel redet und predigt noch, und wird bekannt in aller Welt, daß wir alle fromm werden allein durch den Glauben.

21. Darum siehe, wie sein Herz gestanden sei. Siehe, der allmächtige GOtt hat das Wort und Zusagung Adam gegeben, und uns verkündigen lassen, daß ein Same kommen solle vom Weibe, welcher alles Unglück zerträte, das der Teufel gestiftet hat. Nun bin ich auch unter denselben Jammer, habe es mit mir von Vater und Mutter gebracht, und bin natürlich ein Kind des Zorns und des Todes [Eph. 2, 3.]. Weil ich aber sehe, daß mir GOtt solche Gnade erzeigt, und läßt mir durch sein Wort anbieten, daß er mir von allem Unglück wolle helfen durch den Samen, der kommen soll, so will ich mich darauf erwägen mit ganzer Zuversicht, ihn loben und danken. Darauf kommt er, und bringt das Opfer. Da ist sein Herz in GOtt fröhlich durch das Erkenntniß, steht auf seinem Werk; opfert aber darum, daß er sich äußerlich will erzeigen, daß er GOttes Gnade von Herzen dankt. Siehe, solch gläubig Herz will GOtt haben. Darum thut er sich freundlich zu ihm und seinem Opfer, daß wir solches erkennen lernen, wie es ihm allein um den Glauben zu thun sei.

22. Wie aber das zugegangen ist, daß sich GOtt zu seinem Opfer gehalten hat, drückt Moses nicht aus. Man sagt aber, daß zu der alten Väter Zeit das Feuer vom Himmel kommen sei, und habe das Opfer angezündet, das ein gewiß Zeichen sein solle, daß ihm GOtt solch Opfer ließ wohlgefallen. Wie man auch hernach liest im Buch der Richter [Cap. 6, 21.] von Gideons Opfer. Darum muß man die Opfer der Väter dafür halten, daß sie GOtt eingesetzt habe zum Zeichen, bei welchem man wüßte, daß GOtt gnädig wäre. Wie wir jetzt bei und neben dem Wort ein Zeichen haben, die Taufe und das Sacrament, daß wir beide durch Wort und Zeichen gewiß seien GOttes Gnade, so haben sie auch Wort und ihr Zeichen, daß das Opfer vom Himmel angezündet worden, gehabt, daß ihr Glaube stark und gewiß würde.

23. Dagegen siehe nun auch, wie Kains Gemüth gestanden ist, nämlich also, wie ich gesagt habe: Ich bin der erste Sohn, das liebe Kind, mir gebührt das Priesterthum, der Vater steht auf meiner Seite, [ich] bin mehr und besser denn Abel; so will ich nun GOtt ein Opfer thun. Also hat er ohne Zweifel müssen denken. Denn wenn er den Glauben gehabt hätte und mit sich gebracht, hätte ihn GOtt nicht verworfen. Weil er ihn denn verworfen hat, ist gewiß, daß kein Glaube da gewesen sei, ohne welchen unmöglich ist, GOtt zu gefallen [Hebr. 11, 6.]. Denn der Glaube steht nicht also auf sich und den Werken. Darum urtheilt GOtt recht, verdammt beide, Person mit der Vermessenheit, und das Werk. So ist nun das das erste Exempel der Schrift, darinne wir lernen sollen, was vor GOtt gilt und ihm gefällt, und wie er richtet, nicht nach dem Werk, sondern nach dem Glauben.

\begin{quote}
V. 5—7. Da ergrimmte Kain sehr, und seine Gebärde verstellte sich. Da sprach der HErr zu Kain: Warum ergrimmest du? und warum verstellest sich deine Gebärde? Ist es nicht also? wenn du fromm wärest, so wäre es ein Opfer; wenn du aber nicht fromm bist, so wird die Sünde liegen zur Offenbarung.\footnote{Randglosse: Zur Offenbarung. Das ist, ob die Sünde zu dieser Zeit stille liegt und ungestraft bleibt, so
\quellseite{67}
währet es doch nicht länger, denn bis an Gottes Gericht, da es alles muß offenbar und vergolten werden. Darum hilft kein Heucheln noch Bergen der Sünde.}
\end{quote}
24. Hier findet sich nun der Unglaube Kains mit seinen Früchten. Da sich Gott zu seines Bruders Opfer so freundlich hielt, zu seinem aber nicht, wird er zornig, schlägt den Kopf nieder, sieht sauer, und verdroß ihn übel. Denn Adam und Eva gemeint hatten, er sollte der rechte Mann sein und Gott gefallen. Da ihm das fehlt, bricht es heraus, was er im Herzen hat, und wird erbittert wider Gott und seinen Bruder. Denn wo nicht ein rechter Glaube ist, da sieht gewiß das Herz Gott mit scheelen Augen an, und denkt: Ich wollte, daß Gott nicht Gott wäre. Also thun eigentlich alle hoffärtigen Heiligen, können nicht leiden, daß man ihre Werke verwirft; schlechtes alle ihrem Vater nach. Wiederum, kann Gott auch nicht leiden, daß sie sich empor werfen und andere Leute verdammen; ihr Zorn ist zu gering gegen Gott. Wiewohl aber Abel darüber muß sterben, wird es doch vielfältig gerochen.

25. Als er nun so zürnt, murrt und schnurrt, spricht Gott zu ihm: Warum zürnest du? wärest du fromm, so gülte das Opfer. Als sollte er sagen: Du sollst es umkehren, und zuvor fromm sein, darnach wäre das Opfer auch gut. Darum, wenn du nicht fromm, das ist, ohne Glauben bist, so ist Sünde alles, was an dir ist, und was du thust. Darum wird die Sünde, spricht er, sich vor die Thür lagern, oder, zur Offenbarung liegen. Diese Worte sind etwas finster; ich halte aber, daß dies die Meinung sei: Wenn du nicht fromm bist, so thust du nichts denn Sünde, darauf wird auch die Strafe folgen, nicht lange außen bleiben, wird einmal kommen, und kann können so lange harren, bis du die Thür aufthust. So trifft er eben sein Herz mit diesen Worten: Wenn du ungläubig bist, so steckst du voll Sünde; aber verziehest du es, so wird es herausbrechen; bricht es heraus, so wird es ungerochen nicht bleiben. Dies ist ein einfältiger Verstand dieses Stückes, streckt sich aber auch weiter in das Gewissen.

\begin{quote}
Aber laß sie sich vor dir bücken, und sei du ihr Herr.
\end{quote}

26. Du hast Sünde, spricht er, und steckst voll Unglaubens, bist ergrimmt; aber siehe dich vor, zwinge und halte die Sünde unter dich, daß sie sich vor dir bücken müsse, und du ihr Herr seiest. Und braucht eben die Worte, die er zuvor [1 Mos. 3, 16.] zum Weibe redet: „Du sollst dich vor deinem Manne bücken, und er soll dein Herr sein.“ Will so sagen: Wie das Weib dem Manne unterworfen sein soll, und sich regieren lassen, so soll auch die Sünde dir zu Füßen liegen. Ob du gleich Sünde hast, so sei doch ihr Herr; laß sie nicht über dich herrschen, sondern sei du ihr Herr, und dämpfe sie, daß du nicht thust, was dich gelüstet [Röm. 6, 12.].

27. Aus diesem Spruch haben etliche wollen bewähren die Lehre vom freien Willen; es hat aber keinen Grund. Denn der Spruch ist nicht mehr, denn eine Lehre oder Gesetz, wie man thun soll; als, wenn man zu einem sagt, wenn du recht wolltest handeln, solltest du so oder so thun. Ist allerdings wie ein ander Gebot in [den] zehn Geboten, welche alle das fordern, daß man die Sünde und böse Lust nicht lasse herrschen, sondern dämpfe und unterwerfe. Wie man nun aus den andern Geboten nicht kann den freien Willen beweisen, so kann man auch aus diesen Worten nichts schließen; sintemal alle Gebote eben dazu gegeben sind und dienen, daß wir sehen, wie der freie Wille und unser Vermögen nichts ist. Denn wir kein Gebot aus unsern Kräften können halten, wie auch Kain dieses nicht gehalten hat.

28. Da ist aber noch keine Uebelthat beschrieben, sondern die hohe Sünde, so da steckt in den Werktheiligen, welche, ob sie gleich viel gute Werke thun, doch inwendig im Herzen daneben voll Zorn wüthen und toben.

29. Also ist uns bei diesen zween Brüdern vorgestellt ein Exempel beide falscher und rechtschaffener Heiligen. Kain ist ein Bild, ja, ein Vater aller Heuchler, die freundlich Gott dienen mit schönen Werken, aber mit falschem Herzen; Abel aber der Vornehmste aller, so da wandeln in rechtschaffenem Glauben, und Gott von Herzen dienen. Dieselben zween Haufen aber mögen nimmer mit einander eins sein; ist nicht zu hoffen noch zu denken, daß man Gottes Wort so handele, daß es beide Theile annehmen. Die Werktheiligen müssen die andern Gläubigen verfolgen, wie der Kain. Wenn man ihre Werke verwirft, heben sie an zu toben und zürnen wider Gott und den Nächsten. Wie man jetzt sieht, wie wüthig sie sind wider das Evangelium und seine Prediger.
\quellseite{68}

30. Diese Worte aber, die hier zu Kain gesagt worden, sind, als ich halte, durch Adam geredet, denn er ist der oberste Vater gewesen eine lange Zeit hernach, bis ins neunhundertsechste Jahr, daß er auch den Patriarchen Lamech, den Vater Noah, erlebt hat, im achten Gliede nach ihm. Darum hat Gott durch ihn gepredigt, und hier den Kain gestraft; wie er jetzt durch uns thut. Ich halte nicht, daß eine Stimme vom Himmel herab kommen sei, sondern ist, gleichwie Christus Matth. 19, 5. sagt, Gott habe [1 Mos. 2, 24.] gesagt: „Ein Mensch wird Vater und Mutter verlassen, und an seinem Weibe hangen“; so doch, wie wir oben [Cap. 2, § 37] gehört haben, Adam dasselbe geredet hat. Daß Christus so will sagen, daß Gott solches durch Adams Mund gesagt habe. Wie auch die ganze Schrift sagt, Gott habe es geredet, was die Propheten sagen. Darum wird Adam neben Kain also gezüchtigt haben, da er sich so stellte, den Kopf niederschlug, und läunisch ward auf den Bruder. Was thut nun Kain dazu, da er gestraft war? Spricht der Text:

\begin{quote}
V. 8. Da redete\footnote{Randglosse: Scham halben mußte er sich äußerlich stellen und reden mit seinem Bruder, weil er gestraft ward, ob er wohl im Herzen ihn zu tödten gedachte. Also ist Kain aller Heuchler und falschen Heiligen Vater.} Kain mit seinem Bruder Abel.
\end{quote}

31. Das ist, Scham halben, daß er so öffentlich überwunden und gestraft ward, mußte er sich äußerlich freundlich gegen ihn stellen, als wollte er nimmer zürnen; richtete den Kopf wieder auf, aber aus falschem Herzen. Daß fast der Spruch daher gezogen ist im 28. Psalm, V. 3.: „Die da freundlich reden mit ihrem Nächsten, und haben Böses im Herzen.“ Dabei der Heilige Geist sein abgemalt hat aller falschen Heiligen Herz, Muth und Sinn, die alle diese Kainische Art an ihnen haben. Was er nun gedenkt im Herzen, mit seinem Bruder zu handeln, beweist er bald hernach, wie folgt:

\begin{quote}
V. 8—12. Und es begab sich, da sie auf dem Felde waren, erhob sich Kain wider seinen Bruder Abel, und schlug ihn todt. Da sprach der HErr zu Kain: Wo ist dein Bruder Abel? Er aber sprach: Ich weiß nicht; soll ich meines Bruders Hüter sein? Er aber sprach: Was hast du gethan? Die Stimme deines Bruders Bluts schreiet zu mir von der Erde. Und nun verflucht seist du auf der Erde, die ihr Maul hat aufgethan, und deines Bruders Blut von deinen Händen empfangen. Wenn du den Acker bauen wirst, soll er dir fort kein Vermögen nicht geben, unstät und flüchtig sollst du sein auf Erden.
\end{quote}

32. Diese Worte alle nach einander hat Gott freilich durch Adam geredet. Siehe, da sind bereits im Anfang der Welt so tröstliche Worte geredet, die da Hülfe, Trost und Beistand zusagen allen, die um Gottes willen leiden. Wie viel seiner Sprüche sind in den Propheten hin und her aus diesem Text gezogen. Abel ging auf rechtem Wege des Glaubens, also, daß er das Leben darüber mußte lassen von seinem nächsten Freunde. Was thut aber Gott? Sobald er todt ist, da Kain meint, er habe gewonnen, sei nun allein das liebe Kind, Abels sei gar vergessen: da lebt er stärker und mächtiger, denn vor nie. Da nimmt sich die hohe Majestät seiner selbst an, kann nicht leiden noch schweigen, daß die Seinen unterdrückt werden. Und ob sie gleich ein wenig unterdrückt werden, so kommen sie bald herrlicher und stärker hervor, denn da sie lebten. Da bringt es Abel jetzt hin, daß seinem Mörder die Welt zu enge wird ohne alles Schwert, und Schläge, daß er an seinem Ort sicher sein kann. So ist schon abgemalt in dem Anfang der Welt, was für ein Lohn bei Gott sei deren, die wohlthun, und wiederum der Uebelthäter; jenen zum Trost, diesen aber zum Schrecken.

33. So wisse nun, daß dies Exempel nicht umsonst geschrieben ist. Zum ersten, daß es so gehen wird und muß, daß wir um des Evangelii willen müssen gewarten Feindschaft, nicht allein fremder Leute, sondern unserer nächsten Freunde und Brüder. Das Christus auch im Evangelio zuvor gesagt hat: „Ich bin kommen, daß ich ein Feuer anzünde auf Erden, was wollte ich lieber, denn daß es wäre schon angezündet? Es werden fünf in einem Hause spännig sein, drei wider zwei, und zwei wider drei. Es wird sich der Vater setzen wider den Sohn, und die Mutter wider die Tochter“ etc. [Luc. 12, 49. 52. 53.] Solches ist verkündigt und geschrieben, erstlich an diesem Ort im Anfang der Welt, und viel mehr hernach beweiset mit mannigfaltigen Exempeln. Darum müssen wir Christen uns getrost solches erwägen, und singen wie die Braut im Hohenliede Salomonis [Cap. 1, 6.]: „Meiner Mutter Kinder zürnen mit mir.“ Man hat mich zur
\quellseite{69}
Hüterin der Weinberge gesetzt; aber meinen Weinberg, den ich hatte, habe ich nicht behütet.“ Als sollte er sagen: Sie wollten, daß ich sollte äußerliche Satzungen hüten und halten, so wollte ich frei sein; sie wollten mich zwingen, so predigte ich dawider, darum muß ich mich leiden.

34. Dies ist ein Stück, daß wir uns deß müssen erwägen, und fröhlich dazu sein, wenn, die uns am nächsten sind, uns feind werden, die uns schützen und handhaben sollten. Denn es ist ein gut Zeichen, daß das Evangelium recht sei, wie wir jetzt auch sehen, wo man es verachtet und verfolgt; wiewohl es denen, die es verhöhnen, nicht gut ist. Das gewisse Zeichen ist, daß es im Herzen gefühlt wird, daß man sein gewiß wird; darnach schlägt jenes dazu. Sonst, wo es nicht wäre, sollte man wohl müde werden, das Evangelium zu predigen. Darum soll man es gern sehen und frisch predigen, wenn man es nicht gern hört.

35. Das andere Stück, das wir hier sehen sollen, ist: So bald der Mensch stirbt um Gottes willen, daß er so bald gepreiset wird, und lebt in Gott. Da ist schon\footnote{Erlanger: schön.} angezeigt die Auferstehung der Todten, daß Abel wieder hervor kommen muß und leben, herrlicher denn zuvor. Denn Gott kann und will sein nicht vergessen. Daher geht der Spruch Christi im Matthäo [Cap. 22, 32.]: „Gott ist Abrahams, Isaaks und Jakobs Gott; Gott aber ist nicht ein Gott der Todten, sondern der Lebendigen.“ Wie er nun Abrahams Gott ist, so ist er auch Abels Gott. Weil nun Gott nicht ein todtes, sondern ein lebendiges Volk haben muß, so muß er auch leben. Denn er sich hier sein annimmt als seines Dieners, daß man sieht, wie er nicht allein lebt, sondern Gott auch für ihn streitet. Er will das Blut nicht ungerochen lassen, es muß wieder hervor. Es ist aber noch nicht gar erfüllt, geht noch immer im Schwange bis an jüngsten Tag. Denn daß jetzt Abel schläft, ist eine kurze Zeit, daß, wenn die Auferstehung sein wird, wird er und wir alle meinen, wir seien erst jetzt gestorben. Es hat wohl jetzt einen Schein in dem Leben, daß sich’s nacheinander zeucht, ein, zwei tausend Jahr; aber dort ist es alles als eine Stunde. So liegt nun Abel, dennoch muß er hervor, und geht die Stimme noch, und wird nicht aufhören zu rufen, bis an jüngsten Tag.

36. Zum dritten ist hier auch das schreckliche Urtheil vorgehalten, und der Lohn, so den Gottlosen und den Verfolgern des Glaubens gegeben wird, sie damit zu schrecken, daß sie abstehen. Denn siehe, wie zugleich hier Gott mit dem Kain fährt und redet. Zum ersten fragt er: „Wo ist Abel, dein Bruder?“ Er aber gibt eine freche und stolze Antwort: „Ich weiß nicht; soll ich meines Bruders Hüter sein?“ Ist das nicht eine große Verachtung göttlicher Majestät? So zeucht immer eine Sünde die andere und größere, bis man gar hin hinter kommt. So verstockt ist er: hat seinen Bruder erwürgt, will es nicht allein leugnen, sondern dazu trotzen und pochen wider Gott; meint, er sei nicht schuldig seines Bruders zu hüten. Wenn ein Fünklein Liebe in seinem Herzen wäre, sollte er sich seines Bruders wie sein selbst annehmen und hüten; so ist er so tief in Blindheit gefallen, daß er meint, er thue wohl daran, daß er so antwortet.

37. Das währet aber nicht länger, denn bis der Reuel kommt, und Gott spricht: „Was hast du gethan?“ da sieht er, was er gethan hat, und schreiet, als ihm die Bosheit aufgedeckt und unter Augen gestellt ist; wie der Prophet [Ps. 50, 21.] spricht: „Ich will dich strafen, und mich dir unter Augen stellen“; also auch von Babylon [Jes. 47, 3.]: „Ich will dich aufdecken, daß alle Welt deine Schande sehe.“ Da ist nun kein Trost, sondern eitel Schrecken und Zagen. So deckt er hier den Kain auf, daß er fühlt, daß er verdammt ist, und in die Hölle soll gestoßen werden. Eben wie Judas, da er [Matth. 27, 4.] schrie und sprach: „O, ich habe gesündiget, und das unschuldige Blut verrathen.“ Wenn die Plage kommt, so ist es aus, ist keine Rettung, sondern die Hölle gar da, daß man nimmer heraus kommen kann, und muß verzweifeln, wo das Evangelium nicht heraus hilft, und Gott Trost gibt, denen geholfen soll werden.

38. Hier aber folgt kein Trost noch süße Worte. Siehst du, spricht er, „deines Bruders Blut schreiet zu mir“. Darinne bleibt er stecken, dazu wird er auch verdammt, und folgt die Strafe auch leiblich, als er spricht: „Verflucht seist du auf Erden“ etc. Ob diese Strafe noch währet über alle Gottlosen, wollen wir hier nicht ausfechten. Aber das will der Text haben, daß Kain gestraft wird um Gut, Leib und Seele. Die Seele ist dem Teufel gegeben, dazu soll ihm die Nahrung sauer und schwer werden; wie auch
\quellseite{70}
geschieht; auch allen Ungläubigen, wenn man es ansieht, daß ihnen ihre Arbeit und Nahrung viel saurer wird denn denen, so da glauben. Da ist viel Mühe und Arbeit, und folgt wenig Frucht; wiederum thun diese leichtere Arbeit, und folgt viel Frucht. Dort ist die Arbeit und Sorge größer, denn es trägt, daß die Kost und Mühe bezahlt; jene haben wohl Arbeit, haben aber nicht Sorge noch Angst.

39. Ueber das spricht der HErr: „Unstät und flüchtig sollt du sein auf Erden“; da ist angezeigt sein blöde und verzagt Gewissen. So soll es gehen allen, die ungläubig sind, und solche böse Stücke auf sich haben, daß sie nimmer sicher sind; wie auch im 3. Buch Mosis [Cap. 26, 36.] steht: Du wirst erschrecken vor einem rauschenden Blatt am Baum; so blöde und verzagt ist es, sonderlich wenn es etwas auf ihn weiß; wie auch an vielen Orten mehr geschrieben ist; als abermal Moses [5 Mos. 28, 67.] sagt: „Des Abends wirst du sagen: Ach, wäre es Morgen! und des Morgens: Ach, wäre es Abend!“ und wird dein Leben vor dir pampeln und hangen, daß du nimmer sicher wirst sein. Item, in [den] Sprüchen [Cap. 28, 1.] sagt Salomo also: „Der Gottlose fleucht, wenn ihn niemand jagt“; fleucht vor seinem eigenen Stern und Gedanken; „aber der Rechtfertige steht wie ein muthiger Löwe.“ Solche Sprüche sind alle hieher gezogen. Denn er hier auch sagt, der Kain soll nimmer sicher sein, nicht vor dem, der ihm nachjage, sondern vor seinem eigenen Gewissen. Dies ist die Straße, der niemand entlaufen kann. Was thut nun Kain? Folgt:

\begin{quote}
B. 13. 14. Kain aber sprach zu dem HErrn: Meine Misthat ist größer, denn daß sie mir vergeben werden möge. Siehe, du treibest mich heute vom Erdboden, und werde mich vor deinem Angesicht verbergen, und muß unstät und flüchtig sein auf Erden. So wird mir’s gehen, daß mich todtschlage, wer mich findet.
\end{quote}

40. Da ist er am Ende, hat ein Urtheil, das bestätigt er selbst; da ist kein Glaube noch Trost. Das Evangelium ist ihm genommen, und beraubt der Erkenntniß GOttes; sieht nichts denn den greulichen Ernst GOttes und seine Sünde, darum muß er verzweifeln und untergehen. So müssen nun auch sagen und fühlen alle Verdammten. Und siehe doch, was er thut. Denn solch verzagt Herz sieht hin und her, weiß nicht, wo es bleiben soll, [es] wird ihm die Welt zu enge, daß er sagen muß: „Siehe, du treibst mich vom Erdboden.“ Das sind eitel verzweifelte Worte. Wie sollte er ihn vom Erdboden stoßen, war doch die Welt so weit? aber das Gewissen macht ihm wohl tausend Welten zu enge. Item, wie sollte er sich vor GOttes Augen verbergen? Wer kann GOtt entlaufen? Es ist alles des Gewissens Schuld, daß ist in solcher Angst, und wollte gern aus der Welt laufen und vor GOttes Angesicht fliehen, wenn es könnte. Das ist die rechte, höchste Angst der bösen Gewissen, das wird auch eigentlich die höllische Pein sein, daß die Verdammten wollen fliehen und sich verbergen, daß sie GOtt nicht sehe, und nicht können.

41. Item, über das fürchtet er sich noch mehr: „Ich werde unstät und flüchtig sein. So wird mir’s gehen, daß mich todtschlage, wer mich findet.“ Da spricht ihm sein Gewissen selbst ein Urtheil, daß man ihn würgen werde, wie er seinen Bruder erwürgt hat. Und ist wahrlich eine greuliche, jämmerliche Plage, davon unser wenig wissen. O, wie ist der Mord so hart gerochen und bezahlt! Noch haben wir immer Sorge, GOtt habe uns verlassen. Sollten wir doch gern zehnmal sterben, weil wir sehen, wie greulich er die Mörder straft, und uns so reichlich unser Leiden vergilt; und uns nicht zu rächen begehren, sondern GOtt heimgeben, ja, gern für unsere Feinde bitten, so hoch wir können, daß sie GOtt je nicht lasse fallen in solche Angst. So hoch sollte uns ihr Jammer erbarmen; wie auch Abel ohne Zweifel gethan hätte, wo er gelebt hätte, und gern noch einmal gestorben wäre, seinen Bruder zu erretten.

42. Als nun Kain in solcher Angst ist, hebt GOtt die Plage ein wenig auf, doch nicht gar, und verzeucht sie. Denn, hätte er sich so sollen sehen lassen, wie Kain fürchtete, hätte er müssen so bald zu Pulver werden. Darum gibt er ihm eine leibliche Rettung, gibt ihm aber kein Evangelium und Sicherheit, und spricht:

\begin{quote}
B. 15. Aber der HErr sprach zu ihm: Es soll also sein, daß wer Kain todtschlägt, das soll siebenfältig gerochen werden. Und der HErr machte ein Zeichen an Kain, daß ihn niemand erschlüge, wer ihn auch fände.
\end{quote}

43. In diesen Worten hebt GOtt die zeitliche Strafe auf, und macht ihn sicher, daß er nicht soll ermürgt werden, setzt dazu ein Gebot darauf.
\quellseite{71}
Ist aber alles geschehen durch Adam, wie oben [§ 30] gesagt ist. Und hat GOtt ohne Zweifel den Todtschlag darum so hart verboten, daß es nicht so tief einrisse unter den Leuten, wie es doch gethan hat. Hier streiten nun die Lehrer darüber, wann und wie diese Strafe erfüllt sei. Antwort kürzlich: Es liegt keine Macht daran, ob man es nicht wisse; sie meinen wohl, Lamech, der den Kain erschlug [B. 23. 24.], sei also gestraft, daß sieben von seinem Geschlechte erwürgt seien. [Es] ist aber da kein Grund aus der Schrift; (G)Ott wird die Strafe wohl gefunden haben. [Es] ist genug, daß wir die Ursache wissen, warum es GOtt verboten hat.

44. Was ist aber das Zeichen, das GOtt auf ihn legt? Die jüdischen Meister schreiben also: es sei ein solches gewesen, daß er immerdar gegangen ist zitternd und bebend. [Es] ist wohl glaublich, ich weiß es aber nicht, denn es der Text nicht gibt; wiewohl sie es daher ziehen, daß GOtt gesagt hat: „Er soll unstät und flüchtig sein.“ Darum ist es möglich, daß GOtt ein Zeichen an seinem Leibe gesetzt habe, daß, wie er innerlich gebeht und gesagt hat, auch also außen gegangen ist, wenn ihn die Leute sähen, daß sie sagen sollten: Siehe, das ist Kain, der seinen Bruder erschlagen hat, den soll niemand todtschlagen. Also hat er doch müssen die Schande tragen an Leib und Seele. Das ist je greulich gestraft. Darum wohl gesaget ist zum Hebräern [Cap. 10, 31.]: „Schrecklich ist es, in die Hände des lebendigen GOttes fallen.“ Es ist noch gnädig und sanft, wenn GOtt durch Menschen straft; greulich aber ist es, wenn er selbst an Leib und Seele straft.

\begin{quote}
B. 16. Also ging Kain von dem Angesichte des HErrn, und wohnete im Lande Nod, jenseit Eden, gegen dem Morgen.
\end{quote}

45. Wie ist Kain von GOttes Angesichte gegangen, so ihm doch niemand entlaufen kann? Antwort: Es ist so viel gesagt: Er ist geflohen vom Vater ins Elend, und in ein fremdes Land gekommen, da niemand gewohnt hat. Das heißt, von GOttes Angesicht geflohen. Denn wo Christen sind, die GOttes Wort haben und predigen, da ist GOtt gegenwärtig, da sieht er hin, wie die Schrift vielmal sagt. So hat er da müssen fliehen von Adam und der Versammlung der Gläubigen, so viel ihrer auch waren; so ist er auch ohne Liebe und Freundschaft ge-

flohen, und gekommen in ein Land Nod. Da ist der Zweifel, ob das Wort Nod heiße flüchtig, oder ob es ein Name des Landes sei, dahin er geflohen ist. Ich halte es aber für des Landes Namen, das jenseit Eden gelegen sei. Denn wir haben oben [Cap. 2, § 19 ff.] gehört, daß das Paradies gepflanzet sei in Eden, daß es dahin will bringen, daß das Paradies an einem sonderlichen Orte sei gewesen. Nun, was finster ist, das bleibe finster; wir wollen auch nicht viel davon machen.

46. So haben wir nun die zwei Exempel von Kain und Abel vom Anfange der Welt und noch immer kräftig bis ans Ende, was Christen und Unchristen seien, wie GOtt sich derer annimmt und lohnt, die an ihn glauben, und wie greulich er die Ungläubigen straft. Folgt nun weiter:

\begin{quote}
B. 17. 18. Und Kain beschlief sein Weib, die ward schwanger, und gebar den Hanoch. Und er bauete eine Stadt, die nannte er nach seines Sohnes Namen, Hanoch. Hanoch aber zeugete Irad. Irad zeugete Mahujael, Mahujael zeugete Methusael, Methusael zeugete Lamech.
\end{quote}

47. Da beschreibt Moises Kains Geschlecht bis ins siebente Glied; sonderlich sagt er, daß Kain eine Stadt gebauet habe. Diese Stadt hat er gebauet als ein Bürger auf Erden. Denn, wenn es also stünde, daß wir alle Christen wären, dürfte man nicht des weltlichen Schwerters und Schutzes. Die bei Adam geblieben sind, haben keine Stadt gebauet, noch sich gedacht zu schützen und wehren. Dieser aber hat freilich darum gebauet, daß er etwas sicher möchte sein, weil er sich fürchtete und zagte. Da ist abermal angezeigt, was die Welt thut; die müssen am ersten beschrieben werden, und hoch empor kommen. Denn was GOttes Volk ist, kommt langsam hervor, ist geringe; das andere bricht hervor, und steigt bald hoch, bauet, verwahret, und mehret sich zeitlich.

\begin{quote}
B. 19—22. Lamech aber nahm zwei Weiber, eine hieß Ada, die andere Zilla. Und Ada gebar Jabal, von dem kamen, die in Hütten wohnten und Vieh hatten. Und sein Bruder hieß Jubal; von dem kamen, die mit Harfen und Pfeifen umgingen. Die Zilla aber gebar auch, nämlich, den Thubalkain, der war ein Polirer in allen Meisterstücken Erz und Eisens. Und die Schwester des Thubalkain hieß Naema.
\end{quote}
\quellseite{72}

48. Hier sind unsere Lehrer aber[mal] klug gewesen, daß sie sagen, Lamech sei der erste Ehebrecher gewesen (wie sie achten die digamos, das ist, die zwei Weiber haben). Ich halte nicht, daß dies die Meinung sei. Denn der Text sagt schlecht, er habe zwei Eheweiber gehabt. Ob er auch der erste gewesen sei, weiß ich nicht. Aber damit ist nicht geschlossen, daß er unrecht thue, zwei Weiber haben; denn man dasselbe hernach von vielen, auch heiligen Leuten liest. Es ist aber nur so viel angezeigt, wie das Weltvolk heraus bricht, vor die Kinder des Lichtes [Luc. 16, 8.], wollen sich bald ausbreiten und groß machen; dazu sind sie immerdar geschickter denn Gottes Kinder.

49. Also sagt er, wie Jabal, des Lamech Sohn, ein Mann war, gerüstet und auf die Nahrung geschickt, ist damit umgegangen, daß er reich würde, denkt auf Vorrath, daß er beiseite lege und genug habe von Vieh und allerlei Habe. Adam und Eva haben sich damit nicht bekümmert. Aber der andere Sohn, Thubal-Kain, ward ein Meister, der mit Erz und Eisenwerk umging; hat es von erstem aus der Erde gegraben, ist ein Kriegsmann worden, der erste, der die Leute mit Schwert und Eisen gedachte zu zwingen und unter sich [zu] bringen; denkt auf Leute und Land, ihm zu unterwerfen, wie jener auf die Nahrung. Der dritte Sohn, von der andern Frau, ist mit Tanzen, Springen und Hofiren umgegangen, hat sein Datum auf gute Tage, Lust und Freude gesetzt. Das ist der Welt Lauf. So ist es bei einander: nach großer Nahrung und Reichthum trachten, Land und Leute unter sich bringen, Gewalt und Ehre haben, und in Lust und Wohltagen leben.

50. Auch sagt man von der Schwester Naema, wie sie erfunden habe Seidenstücken, Nähen und Wirken, daß da beides, von Mann und Weib, geschickte Leute sind worden. Ist alles geschrieben, ehe Gottes Kinder angehen, zu einem Spiegel, was die Welt sei; wie sie da angefangen hat, thut sie auch noch immerdar; Gott leuget nicht, so lässet sie von Art nicht. So ist nun beschrieben, wie von Kain das weltliche Regiment kommen ist, und seine Kinder geschickte und kunstreiche Leute worden sind, und sich gesetzt haben auf der Welt, sich verpaart mit Städten und allen Dingen. Da ist schon kein Geist mehr gewesen, wiewohl äußerlich ein großer, herrlicher Schein ist eines ehrbaren Lebens, als wir aus dem folgenden Texte sehen werden. Denn so schreibt er von Lamech, als sie nun im Regimente gesessen sind:

\begin{quote}
V. 23. 24. Und Lamech sprach zu seinen Weibern Ada und Zilla: Ihr Weiber Lamech, höret meine Stimme, lasset meine Rede zu euren Ohren kommen. Ich habe einen Mann erschlagen mir zur Wunde, und einen Jüngling mir zur Beule. Kain soll siebenmal gerochen werden, aber Lamech siebenundsiebenzigmal.
\end{quote}

51. Es sind jetzt auf diesen Tag wenig Leute, die so fromm sind als dieser Lamech, sonderlich unter Fürsten und großen Herren. Denn du siehst, daß er sich selbst verdammt und [das] Urtheil fällt, daß der Todtschlag nicht recht sei; wiewohl er vom weltlichen Geschlecht ist, darinne kein Geist war. Dabei zu merken ist, daß die vom Kain ein fein äußerlich Wesen und Regiment gehalten haben, viel besser, denn man jetzt immer hält, auf daß bestehe, wie oben [§ 50] gesagt: Außen fromm, innen Schalk; in seinen Werken, aber ohne Glauben gewandelt.

52. Wie nun dies Geschicht zugegangen sei, gibt der Text nicht; aber so viel gibt er, daß der Lamech den Kain erschlagen habe, denn er es selbst mit seinen Worten bekennt. Nun haben wir oben [§§ 43. 44] gehört, wie Kain gesagt hat, daß, wer ihn todtschlüge, sollte siebenfältig gerochen werden. Das Wort ist unter ihnen gegangen, und haben es für ein gemein Gesetz gehalten. Denn sie haben auch Gottes Wort gehabt; aber nicht das rechte, heilsame Evangelium, sondern nur das Gesetz. Aber Adam und Eva hatten das Evangelium. Wie denn Gott immerdar die Welt regiert, daß er viel Leute unter dem Gesetze hält, wenig aber unter dem Evangelio. Also ist das Wort, zu Kain gesagt, auch Gottes Wort, aber nicht ein gnädiges Wort; keine Verheißung, sondern allein ein Gebot gestellt, was man nicht thun sollte. Nach dem haben sie sich gerichtet, dasselbe gefaßt und getrieben, welches jedermann wohl wußte in dem Geschlechte.

53. Weil wir nun nicht können aus dem Texte haben, wie es gekommen und zugegangen ist, daß Kain von Lamech erschlagen ist, müssen wir es lassen im Zweifel bleiben. Die jüdischen Schreiber sagen, es sei also zugegangen, daß Lamech sei einsmals mit einem Knaben, den er bei sich gehabt, auf die Jagd gezogen, und Lamech
\quellseite{73}
sei blind gewesen; daß er nirgend hin geschossen habe, denn wo er hin gewiesen ward. Nun sei Kain im Walde gekrochen und habe da gebüßt, und als er \unsicher{taub} war, soll er also von Lamech, wie ihm sein Knabe weisete, für ein Wild erschossen sein; als er aber des inne sei worden, habe er im Zorn den Jungen auch erschlagen. Es lautet aber so hin; ich will es nicht für Wahrheit sagen, denn die Juden pflegen auch gute, starke Lügen zu thun.

54. Es sei ihm nun, wie ihm wolle, so ist der Text klar, daß Lamech den Kain erschlagen habe, und den Mord begangen; dem mögen wir trauen. [Es] liegt nicht viel daran, ob man nicht wisse, wie sich’s begeben habe. Aber da liegt die Macht an, warum es geschrieben ist, nämlich, daß Gott mit diesem Exempel hat wollen anzeigen, wie fromme Leute wir sind, wenn wir durchs Gesetz regiert werden, wenn das Evangelium nicht dazu schlägt. Denn niemand ist je vor Gott durch das Gesetz fromm worden, schwerlich auch vor den Leuten. Vor Gott ist damit nicht allein der Bosheit gewehrt, sondern wird nur stärker und größer. Also thut die Natur, wenn man ihr allein die Gebote vorhält, und nichts mehr, daß man sehe, wie es vom Anfang also gegangen ist, daß niemand mit Geboten geholfen ist. Denn Lamech hat je Gottes Gebot gehabt, wie gesagt, daß er Kain kein Leid thun sollte, wie auch zuvor Kain hatte, seinem Bruder nicht Böses zu thun;\footnote{Wittenberger und Jenaer: „gethan“ statt: zu thun.} noch hat es keinen geholfen. Denn daß die Juden sagen, Lamech habe es ohngefähr gethan, ist nicht wohl zu glauben, sintemal hernach im Gesetze Mosis solcher Todtschlag, der unwissend geschieht, gar sanft gestraft wird, daß solchen etliche Städte ausgesondert werden, dahin zu fliehen und sicher zu sein [4 Mos. 35, 11. 15. 5 Mos. 19, 5.]

55. Nun siehe, was er thut. So übel lohnt er dem alten Manne, seinem Ahnherrn, der nun das sechste Glied erreicht hatte, und seit bei fünf oder sechshundert Jahren alt war, und hoch gehalten gewesen ist von jedermann auf dieser Seite, als der Oberste und der Klügste. Darum muß Lamech freilich eine große Bosheit an ihm begangen haben, weil es die Schrift so aufnimmt vor andern. Denn er nicht meldet, daß mehr Wort hier geschehen sei; und doch, weil sein Sohn, Thubalkain, Erz und Eisen, das ist, Harnisch und Gewehre erstlich erfunden und gemacht hat, muß viel Leid und Haß unter ihnen gewesen sein, so, daß es Schwerter zu machen Ursache hat gegeben, daß einer den andern vom Leben zu bringen gedacht hat; wie der andere Bruder sich auf Geiz gegeben hat, und Naema sich auf den Schmuck, davon Hoffahrt, auch Neid und Zorn erwachsen, und ein ganz weltlich Wesen, wie man es noch sieht.

56. Als nun Lamech die Uebelthat begangen hat, da kommt auch das böse Gewissen und heißt ihn, daß er sich noch viel mehr fürchtet, zappelt und zagt denn Kain. Darum fährt er zu, und sagt es seinen Weibern; zeigt fein an, wie er verzagt gewesen ist; hat es sonst niemand wollen sagen, daß er auch Sorge hatte, es würde ihn jedermann erschlagen. Denn es war ein Großes, daß er einen solchen Mann, seinen alten Herrn und Vater, wider Gottes Gebot erwürgte, der das Regiment unter ihnen hatte, und von jedermann geehrt ward; daß ihn auch das natürliche Gesetz verdammt; darum sagt er: „Kain soll siebenmal gerochen werden, aber Lamech sieben und siebenzigmal.“

57. Daß dies Gott geredet habe, sagt der Text nicht; ich halte auch nicht dafür, daß es wahr sei. Darum ist es vergebens, daß wir uns damit bekümmern, wie er sieben und siebenzigmal gerochen sei. Welches etliche ziehen bis in das siebenzigste Glied, bis zur Sündfluth, in welcher nicht allein die von Kain kommen, sondern auch von Adam und Seth, umkommen sind, daß nur acht Seelen behalten sind, als die Kinder Gottes auch gar Fleisch und Blut worden waren. Das ist aber das Allersicherste, weil es die Schrift nicht meldet, daß wir es lassen Gott befohlen sein; er hätte es wohl können offenbaren, wenn etwas daran gelegen wäre. Das ist aber gewißlich geschehen, daß Lamech siebenfältig gestraft ist, weil es Gottes Wort war.

58. So weit haben wir nun gehört von Kains Gliedern. Da schreibt er nicht mehr von, und zeucht nun hinfort an das andere Geschlecht, und theilt sie beide also, wie wir hören werden im sechsten Capitel, daß er sie ein Theil heißt „Menschen Kinder“, eines Theils „Gottes Kinder“. So haben wir gesehen, wie sich die Weltkinder gestellt haben, wie sie noch immerdar thun, daß sie alle Bosheit treiben, und doch Gottes Gesetz vorgeben. Nun sind Gottes Kinder auch nicht
\quellseite{74}
gewesen ohne Mann und Weib, denn es da nicht galt, daß man ohne Ehe wäre. Also haben sie sich beide ausgebreitet; aber Gottes Kinder hatten nicht das Gesetz allein, sondern auch das Evangelium, davon sie auch Gottes Kinder worden sind und heißen. Davon hernach weiter.

V. 25. 26. Adam beschlief abermal sein Weib, und sie gebar einen Sohn, den hieß sie Seth; denn Gott hat mir, sprach sie, einen andern Samen gesetzt für Abel, den Kain erwürget hat. Und Seth zeugete auch einen Sohn, und hieß ihn Enos. Zu derselben Zeit fing man an\footnote{Randglosse: Fing man an. Nicht, daß zuvor nicht auch Gottes Name wäre angerufen, sondern nachdem durch Kains Bosheit der Gottesdienst gefallen war, ward er damals wieder aufgerichtet, und irgend ein Altärlein gebaut, dahin sie sich versammelten, das Gottes Wort zu hören und zu beten.} des Herrn Namen anzurufen.

59. Da beschreibt Moses, was sich geflissen haben die Gottes Kinder auf der andern Seite, gar Widersinns. Denn die Menschenkinder bauen Städte, richten sich auf die Nahrung, auf den Streit, denken, wie sie sich hier setzen, wollen hier guten Muth haben und wohlleben, geben doch daneben vor, sie seien fromm.

60. Aber hier zeugt Adam einen andern Sohn, den heißt er Seth, der soll ein ander Wesen führen, und gibt Ursach des Namens, daß Eva gesagt hat: „Gott hat mir einen andern Samen gesetzt für den Abel.“ Hier sagt St. Augustinus (wer ihn betrogen hat), das Wort soll resurrectio, das ist, Auferstehung heißen. Es heißt aber eigentlich „gesetzt“, als man einen Grund legt oder setzt, so da stehen soll.

61. Dieser zeugt auch einen Sohn, und nennt ihn Enos. Was das heiße, wollen wir hören; erstlich das sehen, daß der Text sagt: „Zu der Zeit fing man an Gottes Namen anzurufen“; das ist, daß damals der Gottesdienst ist angegangen. Das verstehe nicht also, daß er zuvor auch nicht gewesen sei; denn das muß man bleiben lassen, daß sint der Zeit, da das Wort Gottes ist kommen zum Weibe, der Glaube nie untergegangen, denn das Wort ist nicht aufgehoben noch gefallen. Es muß allezeit die christliche Gemeinde bleiben, wenn gleich nur zwei Menschen auf Erden wären; denn um deren willen, so da glauben, läßt er es allein ausgehen. Darum ist kein Zweifel, daß es noch nie ist untergegangen. Es leidet aber zuweilen einen Abbruch; dennoch bleibt es stehen.

62. Was will denn Moses mit diesem Worte? Antwort: Es ist so viel gesagt, daß das Evangelium, so Adam gesagt ward, wieder angegangen ist und öffentlich gepredigt worden. Denn die Menschenkinder, Kain mit seinem Gesinde, als sie Tyrannen wurden und die Leute unterdrückten, machten, daß das Evangelium keinen Raum hatte, noch bekennen werden mochte, bis Gott wieder Glück gab mit dem Enos, daß er anfing, ein Häuflein Gläubige zusammen zu bringen an Einen Ort, da man Gottes Wort predigte, und vielleicht auch einen Altar aufrichteten, da sie Gott lobten, dankten, und baten. Welches denn heißt, Gottes Namen anrufen, wie Paulus sagt [Röm. 10, 13.]: „Wer da Gottes Namen anruft, der wird selig werden.“ Item, A. 14.: „Wie werden sie anrufen ohne Prediger?“ Daß also das Anrufen muß folgen nach dem Glauben, doch damit auch andere dazu [zu] bringen.

63. Sollte es nun dahin kommen, wie hier der Text klingt, daß man Gottes Namen öffentlich anriefe, mußte\footnote{Jenaer und Erlanger: müßte.} es so zugehen, daß Adam und seine Nachkommen haben müssen an etlichen Orten predigen. Darum muß es verstanden werden vom äußerlichen Gottesdienst; nicht wie wir jetzt Gaukelwerk treiben, sondern wie David im 102. Psalm, V. 22. 23., sagt: „Auf daß man zu Zion predige seinen Namen, und sein Lob zu Jerusalem; wenn die Völker zusammen kommen, und die Königreiche, dem HErrn zu dienen“; das ist, daß sie bei einander predigen, taufen, Sacrament handeln, und daneben beten. So haben wir gehört, daß die Väter nicht gehabt haben Taufe noch Sacrament, sondern ein ander Zeichen, nämlich das Opfer, da das Feuer vom Himmel genommen, und dasselbe angezündet hat. Solcher Gottesdienst ist da wieder aufgerichtet, nicht an Einem, sondern an mehr Orten, da sie gewesen sind, daß wieder der Gottesdienst zugenommen hat, und gebessert ward.

64. Daher halte ich nun, daß Enos den Namen muß haben, denn Enosch heißt hebräisch „ein Mensch“. Dieselbige Sprache hat wohl drei Worte und mehr, daß „Mensch“ bedeutet. Isch heißt ein Mann, aber oft in der gemein hin ein jeglicher Mensch. Das andere ist Adam,
\quellseite{75}
hat den Namen von der Erde, als eine röthliche Erde. Also heißt Enoch auch einen Menschen, aber sonderlich von dem Jammer und Herzeleid, das der Mensch leidet, daß man es deutlich möchte: ein geplagter, elender, betrübter Mann, daß es bedeute sein Elend und armes Wesen. So hat er müssen heißen, daß er ein betrübter, armseliger Mensch ist gewesen. Und alle, die Gottes Namen anrufen, die müssen sein nach dem äußerlichen Wesen im heiligen Kreuze, das er hat müssen tragen und erstlich den Grund dazu legen.

\section{Das fünfte Capitel}

Dies ist das Buch von des Menschen Geburt, zur Zeit, da Gott den Menschen schuf, und machte ihn nach dem Gleichniß Gottes. Er schuf sie ein Männlein und Fräulein, und segnete sie, und hieß ihren Namen Mensch, zur Zeit, da sie geschaffen wurden. Und Adam war hundert und dreißig Jahr alt, und zeugte einen Sohn, der seinem Bilde ähnlich war, und hieß ihn Seth. Und lebte darnach achthundert Jahre, und zeugte Söhne und Töchter; daß sein ganzes Alter ward neunhundert und dreißig Jahre, und starb. Seth war hundert und fünf Jahre alt, und zeugte Enos. Und lebte darnach achthundert und sieben Jahre, und zeugte Söhne und Töchter, daß sein ganzes Alter ward neunhundert und zwölf Jahre, und starb. Enos war neunzig Jahre alt, und zeugte Kenan. Und lebte darnach achthundert und fünfzehn Jahre, und zeugte Söhne und Töchter, daß sein ganzes Alter ward neunhundert und fünf Jahre, und starb. Kenan war siebenzig Jahre alt, und zeugte Mahalaleel. Und lebte darnach achthundert und vierzig Jahre, und zeugte Söhne und Töchter, daß sein ganzes Alter ward neunhundert und zehn Jahre, und starb. Mahalaleel war fünf und sechzig Jahre alt, und zeugte Jared. Und lebte darnach achthundert und dreißig Jahre, und zeugte Söhne und Töchter, daß sein ganzes Alter ward achthundert fünf und neunzig Jahre, und starb. Jared war hundert zwei und sechzig Jahre alt, und zeugte Henoch, und lebte darnach achthundert Jahre, und zeugte Söhne und Töchter, daß sein ganzes Alter ward neunhundert zwei und sechzig Jahre, und starb. Henoch war fünf und sechzig Jahre alt, und zeugte Methusalah. Und führte einen göttlichen Wandel dreihundert Jahre darnach, und zeugte Söhne und Töchter, daß sein ganzes Alter ward dreihundert fünf und sechzig Jahre. Und dieweil er einen göttlichen Wandel führete, nahm ihn Gott hinweg, und ward nicht mehr gesehen. Methusalah war hundert sieben und achtzig Jahre alt, und zeugte Lamech, und lebte darnach siebenhundert zwei und achtzig Jahre, und zeugte Söhne und Töchter, daß sein ganzes Alter ward neunhundert neun und sechzig Jahr, und starb. Lamech war hundert zwei und achtzig Jahre alt, und zeugte einen Sohn, und hieß ihn Noah, und sprach: Der wird uns trösten in unsern Werken, und in der Kümmerung unserer Hände auf Erden, die der Herr verflucht hat. Darnach lebte er fünfhundert fünf und neunzig Jahre, und zeugte Söhne und Töchter, daß sein ganzes Alter ward siebenhundert sieben und siebenzig Jahre, und starb. Noah war fünfhundert Jahre alt, und zeugte Sem, Ham und Japhet.

1. Das ist ein schlecht Capitel anzusehen, weil darinne nicht mehr erzählt wird, denn der Erzväter Geburt, und wie Gottes Gebot ist von statten gegangen, als er die Menschen gesegnet und geschaffen hatte, daß sie sollten Früchte zeugen. Und zählt zehn Glied nach einander, das ist, zehn Patriarchen, nämlich: Adam, Seth, Enos, Kenan, Mahalaleel, Jared, Henoch, Methusalah, Lamech und Noah. Das er aber vorne sagt: „Das ist das Buch“, das ist, die Beschreibung, „von des Menschen Geburt“ etc., ist alles zuvor geredet, er holt es aber noch einmal wieder, auf daß er etwas hinzu setze, und weiter beschreibe, wie es forthin ist gegangen. Darum ist es nicht noth, daß man hier lange stehe.

2. Was da nach Gottes Bilde geschaffen heiße, haben wir im ersten Capitel [§ 56] gesagt, nämlich, daß da Gott gleichförmig sei, und was man Gottes Natur Gutes zuschreibt,
\quellseite{76}
daß der Mensch auch darnach, und ihm gleich gemacht sei. Aber dasselbige Bild ist nun untergegangen und verderbt, und an dessen Statt des Teufels Bild aufgerichtet. Aber durch Christum ist es wieder gebracht und erneuert; durch welches Blut wir errettet sind von Sünde, Tod und Teufel, und durch den Heiligen Geist, durch ihn erworben, gerecht, wahrhaftig, grundgut im Herzen, und ins ewige Leben gesetzt werden.

3. Item ist auch genug gesagt von dem, daß er hier wieder anzeigt, daß Gott den Menschen habe geschaffen, ein Männlein und Fräulein, und gesegnet, daß die Gewalt nicht bei uns stehe, zu machen oder zu wehren, daß ein Mann kein Mann, oder Weib kein Weib sei, und wiederum; also auch zu hindern, daß ein Mensch nicht fruchtbar sei. Aber diese Worte werden nicht vergeblich angezogen, sind bald gesagt, aber eitel Worte des Glaubens, und eines hohen Verstandes; verstehet auch niemand, denn der im Glauben ist, was Mann und Weib sei.

4. Das heiße ich aber im Glauben Mann oder Weib erkennen, wenn ein Mensch gewiß und ungezweifelt dafür hält, daß beide, Mann und Weib, Gottes Werk seien. Aber man findet ihrer wenig, die also zusammen kommen, oder bei einander sind. Fleisch und Blut und böse Lust fühlt man wohl, dazu auch Unglück im ehelichen Stande; aber daß es Gottes Werk sei, sieht man nicht. Derhalben geht es also, wie man täglich sieht, daß man sein bald müde wird und darüber schreiet. Das macht nichts denn die Blindheit, daß sie allein darum hinan gehen, daß sie ihren Vormiß und Lust büßen; Gott aber hat es darum eingesetzt, daß die Creaturen gemehrt würden. Aber davon genug. Nun wollen wir die Patriarchen überlaufen.

5. Adam, schreibt Moses, ist alt worden hundert und dreißig Jahre, und zeugete einen Sohn, der seinem Bilde ähnlich war, und hieß ihn Seth etc. Möchte jemand fragen: was er die Zeit gethan habe? und läßt sich das Capitel sonderlich ansehen, als wären die Väter nichts denn Fleisch und Blut gewesen, daß Moses nichts von ihnen schreiben kann, denn wie sie Hochzeit gehabt und Kinder gezeuget haben. Aber es ist freilich kein Buch in der Schrift, das die Keuschheit so hoch lobt, als dieses. Ist es nicht hoch genug gepreiset, daß Adam hundert und dreißig Jahr gegangen ist in Zucht und keuschem Wesen? Dazu auch die andern Patriarchen alle, die so gegangen sind bis in sechzig, siebzig, hundert Jahre? Wo findet man jetzt solche Leute, die es thun könnten?

6. Die Natur ist hernach immer mehr und mehr verderbt, daß jetzt niemand so lange kann keusch bleiben noch leben. Es haben wohl viel Leute geträumt und gesagt, Adam habe die hundert und dreißig Jahre gebüßt in einem Loche in dem Lande, und dergleichen mehr; sind aber eitel Fabeln, stracks wider den Glauben. Denn er hat Vergebung seiner Sünde, sammt Eva, durch die Verheißung von dem Samen bekommen. Der Glaube hat die Sünde ausgelöscht; denn Gott läßt die Sünde nicht mit Werken büßen.

7. Nun ist weiter zu denken, daß Kain auf jener Seite auch Kinder gezeuget hat. Adam aber auf dieser Seite regiert seine Kinder christlich, ist ihr Vater, Bischof, Pabst und Prediger gewesen, hat aber wider das geistliche Recht gehandelt, darinne der Pabst schreibt und setzt, man könnte nicht Gott dienen, und ehelich sein. Dazu hat er nicht allein gepredigt seinen Söhnen, sondern fortan getrieben, und ist lange Zeit der oberste Vater gewesen, und wird gar eine feine, schöne Welt gewesen sein. Denn wo man die zehn Glieder zählt, so finden wir tausend sechshundert und sechs und fünfzig Jahre bis auf die Sündfluth. Nun hat Adam gelebet neunhundert und dreißig Jahr; so ist nach seinem Tode die Sündfluth nicht viel über siebenhundert Jahr gewesen: so nahe ist er hinan gekommen, daß er gelebet hat mit seinen Söhnen bis an Lamech, Noahs Vater. So haben die lieben Väter unter einander gelebt, alte, betagte Leute, in großem Geiste, Weisheit und Verstand die Welt regiert. Da war es recht im Lenz, daß es daher grünt; nahm aber bald hernach ab.

8. So haben wir hier die erste Tugend und edle Frucht des Glaubens, die da heißt die Keuschheit. Die Väter haben sich nicht mehr Fleisches und Blutes gebraucht, denn Kinder zu zeugen; haben Gott gedient mit dem Vortheil, daß sie nicht aus Büchern gelehrt haben. Denn es ist ein groß Zeichen der Gebrechlichkeit, und böses Verstandes und Gedächtnisses der Natur, daß wir müssen Bücher haben. Da ist man zu den Vätern gegangen, hat gefragt nach allen Sachen, und sie [haben] geantwortet aus dem Geiste. Ist fein regiert durch lebendige Stimme
\quellseite{77}
ohne alle Schrift. Jetzt müssen wir uns mit Büchern behelfen. Auch haben sie, weil sie so lange gelebet haben, gar viel Dinge können lernen und erfahren vom Gestirne am Himmel, von Kräutern und Thieren, und allerlei auf Erden, daß sie hochverständig sind worden in allen Creaturen, davon wir gar wenig erlangen mögen. Es ist uns zu bald entfallen; ehe wir recht zu Kräften kommen, müssen wir sterben.

9. Item, dagegen werden sie auch unzählig viel böser Stücke gesehen haben von den Weltleuten, sonderlich auf jener Seite; weil wir jetzt, wenn wir dreißig Jahr erleben, so manche böse That sehen, daß Sünde und Schande ist. Wiewohl die Welt dazumal nicht so böse sein konnte, als sie jetzt ist, und sie den Vortheil hatten, daß unter so viel Patriarchen Gottes Wort im Schwange ging.

10. Summa: Es liegt nicht an Büchern noch Vernunft; es liegt daran, daß Gott Leute auf Erden schickt. So sehen wir in allen Historien und der ganzen Schrift: wenn Gott einem Volk hat wollen helfen, hat er es nicht mit Büchern gethan, sondern nicht anders, denn daß er einen Mann oder zwei hat aufgeworfen; der regiert besser, denn alle Schrift und Gesetze. So muß nun das ein Ausbund eines guten Regiments gewesen sein, da nicht Ein Mann, sondern so viel seiner Leute und alter Väter, dazu voll Geistes gewesen sind und regiert haben. Es ist wahrlich ein seltsamer, theurer Schatz. Man findet noch verständige und gelehrte Leute; aber die auch Gottes Geist haben, sind wenig. Es hat aber müssen sein, daß der Glaube erhalten würde, und die Leute nicht von Gottes Worte fielen. Das sei nun in der gemeine von den Patriarchen gesagt.

11. Aber unter andern müssen wir sonderlich vor uns nehmen den im siebenten Gliede, nämlich Henoch. Von dem schreibt er etwas Sonderliches, daß er zeitlicher, denn die andern, gefreiet habe, im fünf und sechzigsten Jahre, und einen Sohn gezeugt, so er doch der Allerheiligste gewesen ist, daß der Text sonderlich sagt: Er führte einen göttlichen Wandel, und lebte darinne hernach dreihundert Jahr. „Und dieweil er einen göttlichen Wandel führte, nahm ihn Gott hinweg, und ward nicht mehr gesehen.“ Von dem allein wird geschrieben, daß er nicht gestorben sei, sondern hinweg genommen, daß niemand weiß, wo er geblieben sei. Und ist ein großer, hoher Preis und wunderlicher Text, daß er dem Henoch die Ehre gibt, daß er Gott sonderlich gefallen hat. Darum muß er ein hoher Mann im christlichen Glauben gewesen sein, und hohe Predigt geübt [haben]. Sonst, wo er nicht sonderlich hoch wäre gewesen, hätte ihn Gott nicht so eben abmalen lassen, und solches Wunder mit ihm gethan.

12. Zum ersten habe ich gesagt, daß alle heiligen Väter eben denselben Glauben und denselben Christum gehabt haben, den wir jetzt haben. Daß nun dieser hinweg genommen ist, ist ohne Zweifel zu Troste geschehen denen, so an Christum glaubten. Denn Gott ist so fleißig bei seinem Worte, daß er, die es haben, täglich stärkt, beide mit Worten und Werken. Also hat er hier das gethan, der Väter Glauben zu stärken, daß ein ander Leben nach diesem Leben wäre. Als wollte er dadurch anzeigen und sagen: Sehet, damit will ich sehen lassen, daß ich etwas Anderes habe, da ich die Leute behalte, daß sie nicht sterben, wenn ich sie gleich von der Welt nehme. Darum ist die Summa, daß er mit dem Gesicht hat angezeigt die Auferstehung von [den] Todten, und geschieden dies sterbliche Leben und das unsterbliche, und uns einen Trost zu geben, daß, die wir von Fleisch und Blut sind, können dahin kommen, daß wir ewig leben.

13. Da erhebt sich nun die Frage: Wo dieser Henoch hingekommen? oder, wo er jetzt sei? Wir lesen nur von zweien in der Schrift, die Gott also hinweg genommen habe, diesen Henoch und Eliam [2 Kön. 2, 11.]; auch sagt man es von St. Johanne dem Evangelisten. Das lasse ich in seinen Würden, denn die Schrift davon nichts sagt. Nun muß man ungezwiefelt dafür halten, daß Henoch nicht heimlich hinweg sei gestohlen, sonst wäre es ihnen kein Trost gewesen, sondern daß er sichtiglich und scheinbarlich vor ihren Augen hinweg genommen sei, und daß Gott solches zuvor mit Zeichen beweiset habe.

14. Wo ist denn nun die Stätte, da er geblieben ist? Man spricht, er sei im Paradies; und ist auch zu glauben. Was aber das Paradies sei, weiß ich nicht; [es] ist genug, daß man glaube, daß Gott einen Raum habe, da er noch vielleicht auch Engel behalte. Und steht darauf, daß Henoch und Elias noch selbst nicht wissen, wo sie sind. Denn, wie ich oft gesagt habe, wir müssen die Augen gar weit aufthun, wenn wir in jenes Leben wollen sehen. Es geht nicht
\quellseite{78}
also zu, wie hier. Er ist ein solcher GOtt, daß er außer der Welt jemand behalten kann, und in der Luft, oder wo er will, daß er es selbst nicht sehe, geschweige denn ein anderer. So ist genug zu wissen, daß ihn GOtt weggenommen hat, und bei GOtt bleibt; wo er aber ist, wollen und sollen wir nicht wissen.

15. Ob er auch wieder kommen werde, wie man bisher gesagt hat, laß ich fahren; ich glaube es nicht, und habe die Ursache: Wenn GOtt ein solch scheinbarlich Werck thun wollte, daß dieser Patriarch kommen sollte und den Endechrist umpredigen, würde er es ohne Zweifel zuvor an einem Orte in der Schrift gesagt haben. Denn solche Dinge thut er nimmermehr ohne Wort und Zusagung, wie im Amos [Cap. 3, 7.] geschrieben ist: GOtt wird kein Ding thun, er habe es denn zuvor offenbar durch seine Knechte, die Propheten. Dieweil nun das nicht ist, so halten wir es für eine Fabel, wie andere viel mehr bisher erdichtet.

16. Siehe, die lieben Väter alle werden am jüngsten Tage mit uns hertreten, da werden wir sie sehen bei uns stehen, und wir bei ihnen. Da werden wir wissen, wo sie gewesen sind. Jetzt ist es eine lange Zeit, da er hinweg genommen ist; aber am jüngsten Tage wird er eben da stehen, als wäre er dieselbige Stunde hinweg gefahren. Es ist vor GOtt alles vor Augen in einem Hui; da ist keine Zeit, wie bei uns, da eines nach dem andern geht. Das sei genug von diesem Patriarchen.

17. Der neunte Altvater „Lamech“, spricht Mosis, war hundert und zwei und achtzig Jahr alt, und zeugete einen Sohn, und hieß ihn Noah, und sprach: Der wird uns trösten in unsern Werken, und in der Bekümmerung unserer Hände auf Erden, die der HErr verflucht hat“. Was dieser Lamech mit diesen Worten im Sinne habe gehabt, ob ihm vielleicht etwas verkündigt ist oder offenbart, kann ich nicht gewißlich sagen. Mich dünkt, es sei die Meinung, daß er’s dafür gehalten habe, dies sollte der Same sein des Weibes, der der Welt helfen sollte; wie oben [Cap. 4, 1.] Eva von Kain gesagt hat. Darum sagt er eigentlich, wie er wollte errettet werden von der Plage und dem Fluche, [der] über die Erde gegangen.

Es ist auch wohl etwas daran gewesen, aber nicht alles. Denn der Mann Noah hernach der Welt groß geholfen hat.

18. Nun siehe aber darauf, Noah lebte fünfhundert Jahr, und hernach noch einhundert Jahr, bis auf die Sündfluth. Was muß er Arbeit und Mühe gehabt haben auf Erden! Hier auf dieser Seite waren GOttes Kinder, und der rechte Same: noch ist es so böse worden, daß auch der GOttes Kinder viel abgefallen sind. Auf dieser Seite ist auch \unsicher{eingewissen}, daß unter ihnen böse Leute worden sind, und so viel, daß schrecklich ist zu hören, daß von dieser Seite nicht mehr denn acht Menschen erhalten sind von der Sündfluth.

19. Es ist ein großes, greuliches Abnehmen gewesen, von so viel heiligen Vätern, die ohne Zweifel wohl gelebt und regiert haben; was soll denn jetzt in der Höllenlumpen werden? Methusalem hat gelebt bis an\footnote{Erlanger: in.} dasselbige Jahr, da die Sündfluth kam, und bald davor gestorben. Hätte er gelebt, wäre er auch in [den] Kasten kommen. Der hat müssen sehen, daß die Welt so böse worden war, und sich niemand kehrte an sein alt, gut Leben. Also schreiet dieser Lamech auch darüber; als wollte er sagen: Ist es doch jetzt nun alles böse und verflucht.

20. Das wird nun so zugegangen sein, wie mir im nächsten Capitel hören werden, daß sich GOttes und der Welt Kinder zusammen geworfen haben, daraus Tyrannen im Lande geworden, gewaltige und berühmte Leute. Das sahen die Väter wohl. Wenn nun ein Land so böse ist, und die Büberei überhand nimmt, ist [es] zumal ein gewiß Zeichen, daß es müsse brechen und anders werden. Wie es jetzt auch ist, daß jedermann sagt, es stehe alles am allerhöchsten; darum muß entweder der jüngste Tag vorhanden sein, oder ja sonst eine große Plage und Veränderung. Denn die Bosheit ist zu groß, kann nicht höher kommen. Also meint hier Lamech auch, der Sohn, Noah, muß es ändern und besser machen, oder [es] muß gar untergehen. Was nun Noah gethan hat, wollen wir sehen im folgenden Capitel.
\quellseite{79}

\section{Das sechste Capitel}

V. 1—4. Da sich aber die Menschen begannen zu mehren auf Erden, und zeugeten ihnen Töchter, da sahen die Kinder Gottes nach den Töchtern der Menschen, wie sie schön waren, und nahmen zu Weibern, welche sie nur wollten. Da sprach der Herr: Mein Geist\footnote{Handglosse: Mein Geist. Das ist, es ist umsonst, daß ich durch meinen Geist ihnen predigen, sagen und strafen lasse, sie sind zu gar fleischlich worden, verachten und lästern meines Geistes Wort; darum soll er aufhören, und ich will sie lassen fahren, und nicht mehr mich mit ihnen zanken und strafen.} wird nicht immerdar Richter sein unter den Menschen, dieweil sie Fleisch sind. Ich will ihnen noch Frist geben hundert und zwanzig Jahr. Es waren auch zu den Zeiten Tyrannen auf Erden. Denn da die Kinder Gottes\footnote{Handglosse: Kinder Gottes. Das waren der heiligen Väter Kinder, die in Gottesfurcht auferzogen, darnach ärger denn die andern worden, unter dem Namen Gottes; wie allezeit die Geistlichen die ärgsten Tyrannen und Verkehrtesten zuletzt worden sind.} die Töchter der Menschen beschliefen, und ihnen Kinder zeugeten, wurden daraus Gewaltige in der Welt und berühmte Leute.

1. Da fähet Moses an zu beschreiben die große, mächtige Veränderung, zu der Zeit geschehen, als die Welt nun gestanden hatte tausend sechshundert und sechs und fünfzig Jahr, die mit solchem greulichen und schrecklichen Ernst über sie kommt, daß, wer es bedenkt, davor erschrecken muß, daß alle Menschen ersoffen sind, ohne allein Noah selbst achte. Nun haben wir gesagt, daß Moses ihm vorgenommen hat zu schreiben, wo die Menschen herkommen, und derhalben angezogen zwei Schnüre und Linien, der Frommen und Bösen, Geistlichen und Weltlichen. Ist aber oben [Cap. 4, § 47 ff.] genug angezeigt, wie sich Kains Geschlecht bald auf weltliche Dinge begeben hat mit allerlei Hantierung; was zur Nahrung, Ehre, Herrschaft, Schmuck, guten Tagen und Lust dient, ist eines bald nach dem andern eingerissen. Wie denn die Natur thut, wenn man ihr eine Handbreit läßt, daß sie eine ganze Elle nimmt. Hat sich allmählig\footnote{In den alten Ausgaben: meilich.} angefangen, aber so überhand genommen, daß es zu viel und groß worden ist.

2. Als nun Moses spricht, daß die Sünde daher eingerissen ist, daß die Gottes Kinder nach der Menschen Töchtern gesehen haben 2c., zeigt er genugsam an, daß die Welt eitel Fleisch und Blut geworden ist, wie er auch hernach sagt, und bereits dazu gekommen war, daß man solches in eine Gewohnheit gebracht, und für eine Landsitte hielt. Da ist keine Zucht geblieben, bis so lange nicht mehr zu rathen und helfen gewesen ist, daß die Welt mußte verderben; wie auch jetzt vor Augen ist.

3. Solches sagt er nicht von jenen, so von Kain gekommen sind, sondern die von den heiligen Vätern geboren waren. Denn, daß er nennt Gottes Kinder, ist eigentlich zu deuten auf die Menschen, auf den Stamm und Linie von Seth, Adams Sohn. Wiewohl man aber sonst wenig Schrift findet im Alten Testament, daß die Menschen Gottes Kinder heißen, doch ist es damals schon herüben\footnote{berufen = bekannt.} gewesen, daß daher zu nehmen ist, daß es keine hochverständige Leute gewesen sind, und das Evangelium oder die Verheißung, Adam gegeben, wohl verstanden haben, daß, wer da glaube an den versprochenen Samen, alles hätte, was Gott selbst hat, und Gottes Kind würde.

4. Das rede ich darum, daß unsere Schreiber hierüber gerne gewiesen und mancherlei Dinge erdacht haben, wer die Gottes Kinder gewesen wären, weil sie nicht gewohnt waren, daß man Menschen auf Erden sollte Gottes Kinder und heilig heißen, so doch das Vater-Unser, so wir täglich beten, uns das in [den] Mund gibt, daß wir Gottes Kinder sind. Aus solchem Unverstand träumen etliche, daß die Engel dadurch gemeint werden, welche rechte Gottes Kinder sind, als seien sie zu Menschen Töchtern gegangen, und sie beschlafen, daraus denn große Riesen oder Giganten sollen geboren sein. Es ist aber Narrentheidung. Möglich ist es wohl, wie man sagt, daß der böse Geist sich zu den Zauberinnen thun kann, und sie auch schwängern und alles Unflat anrichten. Daher haben sie gedacht, es wäre hier mit den Engeln auch so zugegangen.
\quellseite{80}

5. Darum müssen wir der Schrift gewohnen, daß wir einen rechten Verstand fassen, daß die Gottes Kinder heißen, so da glauben, das ist, die Frommen und Rechtfertigen. Also ist es zu der Zeit fast unterschieden gewesen, Gottes- und Menschen-Kinder, wie wir unterschiedlich genannt haben die Geistlichen und Weltlichen. Wie es aber nun in Mißbrauch gekommen ist mit dem geistlichen Stande, also ist es da auch gegangen. Die Schrift bittet an vielen Orten, daß Gott wolle einsetzen, die da predigen und sein Wort führen, darum, daß, wenn er es nicht thut, wohl ungepredigt bleibt. Aber so ist es allezeit gewesen: wenn erstlich fromme, heilige Leute das geistliche Regiment, zu predigen, rechtschaffen geführt haben, daß nach ihnen fleischliche Leute aufgestanden sind, und desselben Amtes mißbraucht, doch das Regiment, Statt und Namen, so jene gehabt haben, behalten.

6. Also ist im jüdischen Volke erstlich gewesen Moses, Josua und etliche mehr, welche, da sie lebten und regierten nach und durch Gottes Wort, ging es recht zu. Aber nach ihnen kamen andere, die von jenen nichts wußten, hießen dennoch auch Priester und Leviten, erbten und behielten allein den Namen. Desgleichen auch, als die Apostel und ihre Jünger hinweg waren, folgte bald das Volk hernach, die den Namen behielten; sind aber so weit davon als Himmel und Erde. Also ist es auch hier geschehen. Seth, Enos, und ihre Kinder bis in das zehnte Glied haben es fein gehalten; aber als sie nach einander hinweg kamen, hat es ausgebrochen und aufgehört. Da behielten sie wohl den Namen, daß sie auch Gottes Kinder hießen, waren aber gar fleischlich und weltlich.

7. Und sonderlich zeucht er an das, als sie nach der Menschen Töchtern gesehen haben, und zu Weibern genommen, welche sie wollten. Mit diesen Worten gibt er vor, daß kein Gesetz, Zucht noch Ehre gewesen sei, sondern daß sie nach allem Muthwillen gelebt haben, haben sich unter die Weltleute gemengt, Freundschaft bei ihnen gesucht, sind zugefallen und [haben] genommen, welche sie gelüstet, und nirgend nach gefragt. Summa, sie sind gar sicher und ohne Gottesfurcht gewesen.

8. Das ist der erste Fehl, daß sie sind fleischlich worden. Das andere, so aus jenem folgt, ist: wenn die Leute gute Tage haben und wohl leben, kostet es freilich Gewalt, daß sie andere Leute schinden und schaben. Darum sagt er: „Es waren auch zu den Zeiten Tyrannen auf Erden, die waren gewaltige und berühmte Leute“, und zeigt auch an, woher sie kämen, nämlich, wie gesagt, daß die Kinder Gottes der Menschen Kinder beschliefen und Kinder zeugten; davon sind sie gekommen. Das ist es, daß der Name Gottes muß der Schanddeckel sein, darunter alles Unglück angerichtet wird; wie man auch sagt: In Gottes Namen hebt sich alles Unglück an. Was der größte Greuel ist, muß allezeit Gottes Name beschönen und auf sich nehmen, als bis auf diesen Tag geschieht mit unsern Stiften und Messen, Pfafferei und Möncherei, und allem,\footnote{In den alten Ausgaben: alles.} so man Gottesdienst heißt. Wenn man nur den Namen daran hängt, so fällt jedermann darauf, und hält es für köstlich. Die sich nur des Namens dürfen rühmen, daß sie sagen, sie seien Gottes Kinder, da fürchtet sich und erschrickt jedermann vor; da reißt es denn ein, und werden gewaltige Herren unter dem Schein und Namen Gottes.

9. Das will Moses damit, daß er sagt: „Es waren Tyrannen zu den Zeiten“, die nennt er auf hebräisch Nephilim, das haben die Lehrer lateinisch genannt gigantes. Und schreiben, daß sie darum also hießen, daß sie groß gewesen sind, daß die andern gegen ihnen gestanden sind, als wären sie gefallen. Es sind aber auch Fabeln, denn sie sind da alle gleich groß gewesen; wiewohl Moses hernach auch von Giganten oder Riesen wird sagen. Aber darum heißen sie Nephilim, daß sie die Leute überfallen und gezwungen haben. Denn Naphtal heißt hebräisch fallen, davon dies Wort kommt; wie im 10. Psalm, V. 10., steht: „Er fällt herein mit seiner Macht über die Armen“, oder fället den armen Haufen mit seiner Gewalt. Item, im 20. Psalm, V. 9.: „Sie sind getrümmet und gefallen.“ Man hat von ihnen gesagt, wie die Heiden von ihrem Hercules, als von großen Leuten und Helden, die viel Schlachten und Wunder gethan, jedermann haben gezwungen, und niemand gehorcht haben noch unterthan gewesen sind. Daß also die Welt voll Frevels ist worden, wie der Text sagt.

10. Nun siehe, was Gott dazu sagt, da solches begann überhand zu nehmen. „Mein Geist (spricht er) soll nicht immerdar Richter sein unter
\quellseite{81}
den Menschen, weil sie Fleisch sind.“ Die Worte aber und Predigt hat er ohne Zweifel durch Noah gesagt, wie zuvor durch Adam. Da ist nun das Urtheil aufs kürzeste gefällt und gesprochen, aber fast greulich und erschrecklich.

11. Man muß aber die Worte heraus ziehen in das Amt, das der Heilige Geist treibt durch das Wort, das da gepredigt wird; und ist so ein klarer und gewaltiger Spruch, daß genug wäre, wenn er im Neuen Testament stände von dem Heiligen Geiste, wie sein Amt sei, die Welt zu strafen; wie Christus im Johanne Cap. 16, 8. sagt: „Wenn der Heilige Geist kommt, wird er die Welt strafen um die Sünde, Gerechtigkeit und Gericht.“ Und wird hier fein angezeigt, was die Welt sei, wenn der Heilige Geist nicht Richter ist und straft. Könnte sie etwas Gutes thun, und durch ihr Wesen und Werke GOtt gefallen, so dürfte der Heilige Geist nicht strafen und richten. Soll es nun recht stehen und göttlich Wesen sein, so muß der Heilige Geist da sein, sein Amt vollführen, und allein alles thun durch das Wort; daß [es] eine große Gnade ist, wo er so straft und mit den Leuten zankt, und wiederum, die höchste Ungnade und Strafe, wo er sich entzieht und die Welt nicht straft.

12. Darum will er hier so viel sagen: Es ist vergebens, was ich bisher und noch durch meinen Geist in den Patriarchen, Adam, Seth, Henoch, predige, strafe, und mich bemühe, der Leute Bosheit zu steuern und mehren. Sie sind zu gar im fleischlichen Wesen ersoffen, wollen nicht hören noch glauben, sondern verachten und lästern das Wort, das ich durch meine Kinder, die meinen Geist haben, predigen lasse; darum will ich auch ablassen und aufhören zu strafen, und sie hin lassen gehen. Dies ist die greuliche Plage, die GOtt über die Welt gehen läßt, wenn er sie vertilgen will, daß er zum ersten seine Hand abzeucht und sein Wort aufhebt, daß ein jeglicher nach seinem Wahn und Dünkel hingeht, tappt wie ein Blinder, und nicht weiß, was er thut; da folgt denn alle Sünde und Schande auf einem Haufen, so lange bis alles muß untergehen und in Grund verderben.

13. Das ist es, daß Moses im fünften Buch [Cap. 28, 28.] heißt Wahnsinnigkeit, Blindheit und Rasen des Herzens, da er die Flüche erzählt, so über die, die GOttes Wort verachten und nicht hören wollen, gehen sollen. Diese Plage dräuet auch der Apostel Paulus 2 Thess. 2, 10—12. demselbigen Haufen: „Dafür, daß sie die Liebe der Wahrheit nicht haben aufgenommen, daß sie selig würden, wird ihnen GOtt kräftige Irrthümer senden, daß sie glauben den Lügen, auf daß gerichtet werden alle, die der Wahrheit nicht glauben, sondern haben Lust an der Ungerechtigkeit.“ Wie es denn auch über die Welt gegangen ist, wie wir jetzt vor Augen sehen, daß\footnote{Jenes: da.} wir meinten, wir hätten den rechten Gottesdienst und wäre alles recht, ist aber der Teufel und Tod gewesen. So wird es zuletzt wiederum gehen: weil man das Evangelium nicht hören noch annehmen will, sondern auf das allerärgste handelt und lästert, und niemand darnach thun will, wird es GOtt wiederum hinweg nehmen, und die Welt wieder mit größerer Blindheit strafen, denn zuvor je, so lange bis er endlich alles vertilgt. Das ist der Verstand und Meinung dieses Spruches.

14. Weiter spricht GOtt: „Ich will ihnen noch Frist geben hundert und zwanzig Jahr.“ Das redet er auf die Zeit, die er der Welt noch geben wollte, bis auf die Sündfluth, daß sich die Leute indeß bekehren und bessern sollten. Nun war Noah daselbst, wie der Text [Cap. 5, 32.] sagt, fünfhundert Jahre alt, und wird hernach angezeigt, daß nur hundert Jahre auf die Sündfluth waren, als er den Befehl kriegte, die Arche zu bauen, daß es eben zusammen sechshundert Jahr werden, als die Sündfluth kommen ist. Ist nun die Frage: wo denn die zwanzig Jahr bleiben, die GOtt in diesen Worten hinzu setzt? Ich weiß nichts darauf zu antworten, noch aufzulösen, ohne, daß es wohl sein mag, daß die Bosheit so trefflich überhand genommen habe, daß GOtt geeilt habe mit der Sündfluth, und die zwanzig Jahre abgebrochen; oder daß es per anticipationem gesagt sei, also, daß diese Worte, „zwanzig Jahre“, zuvor geredet sind, ehe Noah die drei Söhne gezeugt hat, oder je ehe er fünfhundert Jahre völlig alt worden ist.

15. Es geht aber, wie gesagt ist, eben auf die Erfristung der Jahre, die GOtt der Welt zur Besserung gegeben hat bis zur Sündfluth: nicht, wie etliche sagen, daß er damit habe bestimmt, wie lange eines Menschen Alter und Leben sich strecken solle. Denn wo der Verstand recht wäre, würde daraus folgen, daß jeglicher hundert und
\quellseite{82}
zwanzig Jahr lebete, so doch der meiste Haufe noch jung dahin stirbt, und für ein Wunder gehalten wird, wenn man einen hört, der hundert und zwanzig Jahr lebe; dazu man in der Schrift nicht findet, daß Gott vorstelle, wie lange einer leben solle. Darum ist das Ziel der ganzen Welt vorgestellt, durch Gottes Langmüthigkeit, ob sie sich bekehren wollten; [es] ist aber nicht besser, sondern von Tag zu Tag ärger worden.

16. So ist uns nun diese Historie wohl zu merken; denn es ist eine große, schreckliche Veränderung, darinne man sieht, wie Gott auf Erden zu regieren hat. Er macht es heimlich, sieht durch die Finger so lange, daß die Leute so blind und thöricht werden, daß sie sagen, es sei kein Gott, und sicher verachten alles, was Gott predigen läßt. Also haben wir ein Exempel und Warnung, daß wir gewißlich glauben, daß Gott das Unrecht und Böse nicht ungestraft noch ungerochen läßt, ob er schon eine Zeitlang verzeucht. Das ist ein Stück dieses Textes. Folgt weiter:

\begin{quote}
V. 5—10. Da aber der Herr sahe, daß der Menschen Bosheit groß war auf Erden, und alles Dichten und Trachten ihres Herzens nur böse war immerdar, da reuete es ihn, daß er die Menschen gemacht hatte auf Erden, und es bekümmerte ihn in seinem Herzen, und sprach: Ich will die Menschen, die ich geschaffen habe, vertilgen von der Erde, von den Menschen an bis auf das Vieh, und bis auf das Gewürme, und bis auf die Vögel unter dem Himmel; denn es reuet mich, daß ich sie gemacht habe. Aber Noah fand Gnade vor dem Herrn. Dies ist die Geburt Noah. Noah war ein frommer Mann und ohne Wandel, und führte ein göttlich Leben zu seinen Zeiten, und zeugete drei Söhne, Sem, Ham, Japhet; aber die Erde war verderbt etc.
\end{quote}

17. Hier fragt man und bekümmert sich fast, wie es zugehe, daß der Text sagt, es habe Gott gereuet, daß er den Menschen gemacht hatte, weil man spricht, es könnte ihn nichts reuen, was er thut, sintemal er die höchste Weisheit ist? Antwort: Ich halte, daß dieses Stück geschrieben sei, anzuzeigen der lieben Väter Klage und Geschrei, die da gelebt haben, daß sie solches Dräuen und schreckliches Urtheil Gottes, der es alles wandeln wollte auf Erden, gefühlt haben. Desgleichen auch hernach [1 Mos. 18, 23—32.] angezeigt wird, da Gott die Städte Sodoma und Gomorra wollte umkehren, wie sich Abraham ihrer annahm, sich ängstete, und Gott so hart anlag mit Flehen, daß er sie gern errettet hätte, wie hernach folgen wird. Also ist auch hier wohl zu denken, wie des frommen Vaters Noah Herz gestanden ist, daß er alles, was kommen würde, gefühlt, und sich darum bekümmert und geängstiget habe; daß man es nicht darf Gotte zueignen, als [ob es] ihn reuete und leid wäre, daß er es so ändert und wandelt; denn es bei Gott schon alles ist beschlossen gewesen. Aendert er es aber, so ist es nicht anders, denn daß er es will umkehren, wie er von Ewigkeit versehen hat. Aber die Aenderung thut er, daß es die frommen Leute zuvor fühlen; die haben so gedacht: Ach! Gott hat es nun im Sinn, alles zu ändern und umzustoßen.

18. Dieser Regel muß man oft in der Schrift wahrnehmen, daß von Gott geredet wird, wie wir es fühlen. Denn wie wir ihn fühlen, so ist er uns. Denkst du, er sei zornig und ungnädig, so ist er ungnädig. Also, wenn die Schrift sagt, Gott sei zornig, ist es nicht anders, denn daß er so gefühlt wird. Also steht oft in Psalmen: Wache auf, Herr, warum schläfest du? Wie verbirgst du dich? [Ps. 44, 24. 25.] und dergleichen, so doch keine Natur und Wille seine Veränderung hat, ohne daß er sich so und so stellt, und uns sehen und fühlen läßt. Darum wie es hier Noah gefühlt und sein Herz getroffen hat, so ist es gewesen. Also wollte ich die Frage auflösen; denn es ist das Sicherste, daß man erniedert\footnote{Erlanger: ernieder.} bleibe, und ist dennoch Gottes, nicht Menschen Gedanke. Denn er solche Betrübniß und Bekümmerniß in das Herz gibt, daß es nicht von Menschen, sondern aus dem Heiligen Geist kommt.

\begin{quote}
V. 13—16. Da sprach Gott zu Noah: Alles Fleischliches Ende ist vor mich kommen; denn die Erde ist voll Frevels von ihnen. Und siehe da, ich will sie verderben mit der Erde. Mache dir einen Kasten von Tannenholz, und mache Kammern drinnen, und verpiche sie mit Pech, inwendig und auswendig, und mache sie [die Arche] also: Dreihundert Ellen sei die Länge, fünfzig Ellen die Weite, und dreißig Ellen die Höhe. Ein Fenster sollst du daran machen, oben an, einer
\end{quote}
\quellseite{83}
Elle groß. Die Thür sollst du mitten in ihre\footnote{Jenner: der.} Seite setzen. Das untere Theil sollst du zweischichtig und dreischichtig machen.

19. Daß wir den Text verstehen, wollen wir erstlich einfältig dahin sehen, wie die Arche gebaut ist; darnach, was Noah für einen Glauben gehabt habe; zum dritten, was damit bedeutet ist. Zum ersten, daß man die Archa nennt, ist ein lateinisch Wort, in die deutsche Sprache gerathen; als wir sonst viel mehr lateinische und welsche Worte haben. Wir heißen es auf unsere Sprache eigentlich einen Kasten, oder eine lange Lade. Wenn man es nun abmißt nach dem, wie es hier beschlossen ist, und faßt es in Haufen: es ist sechsmal länger denn breit, und zehnmal länger denn hoch gewesen. Das heißen die Gelehrten, proportionem geometricam et arithmeticam, wenn man sagt, ein Ding ist zweimal so lang, als es breit ist; da man nicht nennt und ausdrückt, wie viel Ellen oder Spannen es in sich hat.

20. Also sollte der Kasten länger denn breit sein, daß, wenn er kleiner gewesen wäre, dagegen zu rechnen gewesen ist wie ein Sarg. Denn wenn man einen Mannesleib vor sich nimmt, und recht messen will, ist er zehnmal länger denn hoch, und sechsmal länger denn breit. Also haben auch die Lehrer die Arche abgemessen, und also gedeutet, daß sie daraus machen den Herrn Christum. Das lassen wir aber gehen; ist genug, daß man so viel wisse der Historie nach, als ich gesagt habe. Nun soll er eine Thür in den Kasten machen, mitten in die Seite, also, daß zehn Ellen drüber und zehn Ellen drunter blieben, daß sie fast vierzehn Ellen oben, und vierzehn unten gestanden ist. Denn der Kasten wird fast nahe gegangen sein zehn Ellen im Wasser, so blieben zwanzig Ellen drüber in der Höhe. Ueber die Thür aber, die mitten in der Seite war, sollte das Fenster stehen, und nicht größer sein, denn einer Elle in die Länge und Breite.

21. Dabei ist abzunehmen, was für Licht in der Arche gewesen sei, und wohl zu spüren, daß die Schrift hier anzeigen will, daß es nicht darum gemacht ist, daß sie Tag davon hätten, noch sich nicht nach der Sonne möchten richten. Denn es oben auf allen Seiten ist zu gewesen, daß der Tag nur eine Elle weit hinein gegangen ist über der Thür; so will daraus folgen, daß sie in dem Kasten immer Licht gebrannt haben, sintemal sie sich nicht konnten von der Sonne Licht behelfen, und das so lange getrieben, so lange sie darinne sind gewesen, daß es eigentlich zugegangen ist wie im Tabernakel Mosis, da kein Tag hinein kommen konnte, sondern ohne Unterlaß sieben Lichter brennen mußten [2 Mos. 25, 37. 27, 20. 21.], wie wir bei Nacht thun. Also ist es nur darum gebauet, daß sie möchten heraus sehen, ob es Tag oder Nacht wäre.

22. Weiter sagt er: daß untere Theil soll er zweischichtig und dreischichtig machen. Denke also, daß der Kasten drei Boden habe gehabt, unten zehn Ellen, mitten auch zehn, da die Thür gestanden, und oben auch so viel, da das Fenster ist gewesen. Im untersten hat er Kämmerlein gemacht, hier zwei, da drei, und so fort mancherlei Wohnungen, darinne dieses und jenes Viehe und Thiere stehen sollten; für dieselben hat er einen solchen Bau gemacht, oder vielleicht für sich auch; also, daß der Kasten eine Behausung wäre allerlei Thiere, auf daß von allen Samen erhalten würde. Das ist das Gebäu der Arche, für das erste.

23. Zum andern ist in dieser Historie angezeigt der Glaube Noah, welches das beste Stück hierinnen ist. Da lässet uns sehen, wie ein christlicher Glaube solle geschickt sein, und wäre wohl recht, daß man solche Historien ausstriche, wie sie wohl werth sind; es fehlt aber daran, daß man es nicht mit Worten erlangen kann, und daß wir nicht so gänzlich nachdenken können, wie ein großer, mächtiger Glaube dies gewesen sein muß, daß es überaus ist; daß auch, wenn man es recht ansieht, nicht Wunder ist, daß nicht mehr denn Noah, selb achte, in der Sündfluth erhalten\footnote{Erlanger: behalten.} ist. Denn siehe darauf: Noah hat nun drei Söhne, die haben drei Weiber; das sind zusammen, mit seinem Weibe, acht Personen, die sind ohne Knecht, Gesinde, ohne Freundschaft nicht gewesen, wie es noch in der Welt ist, und zu derselben Zeit größer, denn jetzt war. Von denen allzumal, wie große Freundschaft er hatte, ist ihm und seinen Söhnen niemand angehangen, sondern schlecht alle von ihm gewichen; daß es je ein Jammer gewesen ist, daß sein Wort und alle Predigt, so er gethan hat, nichts geholfen, und niemand geglaubt noch angenommen hat.
\quellseite{84}

24. Dagegen ist es je ein Großes, daß die acht Menschen allein so fest stehen sollen, und gewiß dafür halten, sie haben allein recht, und so viel Leute in der ganzen Welt allesamt unrecht wären. Wenn es sich heute so begäbe, daß ein Mensch käme, der sagen und schließen dürfte, er hätte allein recht, und aller Anderer Thun und Meinung falsch und verloren wäre; nicht angesehen, wie weise, heilige, mächtige Leute, Pabst, Bischöfe, Fürsten etc. sie sein möchten: wer könnte es doch glauben? Sehen wir doch, wie sie jetzt wider das Evangelium schreien und toben, daß ihnen niemand sagen noch wehren kann. Darum ist kein Zweifel, es sei dazumal auch so gegangen.

25. So rechne es nun eben nach: Noah war nun fünfhundert Jahre, als er den Kasten anhebt zu bauen, und bauet darüber ganzer hundert Jahre. Da hat er täglich gepredigt, getrieben und gesagt: Bessert euch, GOtt ist zornig, wird die Welt mit der Sündfluth verderben, darum baue ich den Kasten, wie mir GOtt hat befohlen. Das ist freilich weit erschollen in die ganze Welt, darum ihm auch GOtt so lange Zeit und Weile dazu gegeben hat. Was ist aber geschehen? Er hat immer fort und fort gepredigt und nicht abgelassen. Daher ihn Petrus [2. Ep. 2, 5.] nennt praeconem justitiae, einen Ausrufer und Ausschreier oder Prediger der Gerechtigkeit, welchen Titel die Schrift seinem Patriarchen so hoch gibt.

26. Aber die Welt that, wie sie noch immerdar thut, verhöhnte und verachtete GOttes Wort und seine Prediger, daß je mehr er davon predigte, je mehr er hören mußte, daß sie sagten: Du bist ein großer, alter Narr, hast Sorge, der Himmel falle auf dich, predigst und schreiest nun zehn, zwanzig, dreißig, vierzig, fünfzig Jahre etc. von der Sündfluth; wann kommt sie denn? Meinst du, daß du allein der Frommste bist, daß GOtt die ganze Welt verderben sollte, und deiner\footnote{Erlanger: dir.} allein schonen? oder, daß du der Klügste seiest, daß er es niemand denn dir sagen sollte?

27. Solches mußte er täglich hören, und noch viel unzählig größeren Hohn und Schmach; läßt sich aber nicht erschrecken, predigt immer anhin, bis ins sechzigste, siebenzigste, achtzigste und neunzigste Jahr. Da wurden sie erst sicher und gewiß. Ja, ja, sollte GOtt so lange harren, wenn er die Welt vertilgen wollte? wenn er so zornig wäre, würde er nicht zehn Jahre geharrt haben. So ließ GOtt der Welt so lange predigen, und sein Wort für Thorheit und Lügen halten, und auf das allerschändlichste und schmählichste handeln. Da siehe, was die Welt mit ihrer Klugheit und Vernunft ist, wie fein sie sich zu GOttes Wort stellen kann. Nun ist es noch Wunder, daß sein Weib, seine Söhne mit ihren Weibern geglaubt haben, und an ihm allein gehangen, dennoch so viel mit ihm blieben sind,\footnote{Der Sinn ist: Es ist ein Wunder, daß dennoch alle, die bei ihm geblieben sind, an ihm allein gehangen haben.} sonderlich weil es sich so lange verzog, als er immer sagte, es würde kommen. Wie viel sind auch zuvor Leute da gewesen, die Geld genommen haben, und den Kasten gemacht, aber nichts geachtet; noch ist er darauf geblieben, und hat geglaubt, es würde so gehen, GOtt würde alles ersäufen.

28. Darum lerne nun, was des Glaubens rechte Art sei, daß er keinem Ding folgt, ohne dem bloßen Wort. Ist es doch alles widersinnig, treibet wider alle Vernunft der Menschen, ja, wider seine eigenen Sinne, und alles, was er nur sieht, fühlt und hört, daß er wohl denken möchte: Wie, sollte denn GOtt so wunderlich fahren, daß er mich allein erwählt, und alle Menschen auf Erden verdammt? wie auch solches Noah ohne Zweifel gefühlt hat und große Anfechtung erlitten.

29. Denn das ist auch die große Anfechtung in der Welt, ausgenommen die hohen geistlichen, daß man denkt: Siehe, die ganze Welt hält es anders, solltest du denn allein recht haben und der Klügste sein? Lieber, wie wenn du irrst? meinst du nicht, daß GOtt barmherzig sei, er werde viel Leute erhalten, und nicht alles so gar vertilgen? Wenn er dahin hätte gesehen, und nicht die Augen zugethan, würde er nicht sein stehen geblieben. Doch hat er solcher Stöße und Püffe oft viel müssen leiden.

30. Er sahe, daß GOttes Wort da stund, von acht Personen allein zu erhalten. Das war so viel gesagt, als spräche er: Nur todt, todt, alles, was da lebt. Da steht der Glaube und spricht: GOtt kann und will nicht lügen. Darum hat er die Augen zugethan und alle Vernunft gedämpft, und allein an dem Wort gehalten. Also muß der Glaube alle Vernunft, Sinne und Verstand mit Füßen treten, und alles, was er sieht, aus den Augen setzen, und nichts wollen wissen
\quellseite{85}
denn Gottes Wort. So lasset uns nun die Vernunft und unser natürlich Licht preisen und rühmen; was ist es anders, denn Blindheit und Finsterniß? Was kann es Besseres, denn wider den Glauben und Gottes Wort fechten? Was gilt nun lange Gemohrheit, und daß so viel Leute in der Welt anders halten und glauben?

31. Solches alles hat er nun hundert Jahr getrieben durch solchen Glauben, den niemand genug erlangen noch ermessen kann. Hilf Gott! könnten wir doch kaum, wenn es aufs Treffen geht, eine Stunde stehen, und er ist da so lange Zeit ohne Unterlaß im höchsten Kampfe gestanden und nicht gewankt; und nicht allein das, sondern hat dazu täglich gepredigt, und Hohn und Spott zu Lohn gekriegt, daß sein Herz recht wohl gekreuzigt ist worden.

32. Also müssen wir auch thun. Wer selig will werden, muß ihm nicht vornehmen, daß er wolle glauben, was Papst, Bischöfe und die ganze Welt glaubt, oder die Concilia beschließen. Wer so glaubt, ist schon verloren. Denn so lange der Haufe steht, steht er auch; steht nicht auf Gottes Wort, sondern auf Menschen Wahn; denkt, sollten jene irren und fehlen, so viele und große Leute? und sieht nicht, daß Gottes Wort stracks dawider sagt: Viele und große Leute sollen irren [Ps. 62, 10.]. Zu Noah Zeiten sind ja unzählig viel gelehrter, kluger, frommer, ehrbarer Leute gewesen, mehr denn jetzt, weil die Welt noch jung war; dennoch mußte er seinen Sinn brechen, und allein halten, was Gott sagt.

33. Solches wird sich fein finden, wenn wir sterben sollen. Da muß der Glaube bloß stehen, und kämpfen wider den Teufel und alle Sinne und Vernunft. Da wird er dir vorhalten: Das sagt der Papst und so glaubt die ganze Welt; wie dünkt dich, wenn es fehlte und irrte? Da muß das Herz, das sich darauf verlassen hat, anheben zu zappeln, steht so bald im Zweifel; wenn es denn zweifelt, so fährt es dahin und muß verloren werden. Darum müssen wir dahin kommen, daß der Glaube allein auf Gottes Wort stehe, daß wir können sagen: Gott gebe, die Welt glaube sonst oder so, da habe ich Gottes Wort, wie Noah hatte; ist jemand in der Welt, der es anders hält, den lasse ich fahren, Gottes Wort muß recht haben, das trügt nicht [Ps. 24, 13.].

34. So sehen wir hier den großen, mächtigen Ernst und Zorn Gottes, daß er nur die acht Personen erhält; dazu das große Werk, daß ein Menschenherz das sollte glauben. Es ist beides, Gnade und Zorn, unbegreiflich und unglaublich aller Vernunft und Menschenkräften. Darum sagt der Text [1 Mos. 7, 1.] wohl: Noah war rechtfertig und ohne Wandel zu seinen Zeiten. Item [1 Mos. 6, 8.]: „Noah fand Gnade bei dem Herrn.“ Die Worte sind so gesetzt, daß sie etwas sollen gelten, nämlich so große Dinge, wie wir gehört haben. Aus dem folgt auch, daß die acht Personen sind heilig gewesen. Acht lebendige Heiligen, und blieben doch nur sechs; denn Ham fiel nach der Sündfluth, und [ich] halte, auch sein Weib. Es ist schrecklich, daß einer in solchem Glauben steht, und dennoch dahin fallen soll. Davon wir hernach hören werden.

35. Zum dritten wollen wir auch die heimliche Deutung dieser Historie sehen. St. Petrus [1. Ep. 3, 21.] hat sie so ausgelegt, daß die Sündfluth bedeutet habe die Taufe, denn dieselbige eine geistliche Sündfluth ist, wie jene äußerlich war. Und wie Noah im Kasten selb achte behalten ist, so werden wir selig in der Taufe. Nun werden durch die Taufe viel mehr Leute ersäuft, denn durch jene Sündfluth, denn sie durch die ganze Welt ausgebreitet ist, und noch immerdar geht. Wie nun dort alle Menschen ersäuft sind, ohne die achte, die im Wasser erhalten sind; so wird auch in der Taufe alles ersäuft, was fleischlich ist. Denn wir werden also getauft, daß wir sterben der Welt, dem Fleisch und Blut; wie Paulus [Röm. 6, 3. 4.] sagt: „Alle, die in Christum getauft sind, die sind in seinen Tod getauft; so sind wir je mit ihm gestorben durch die Taufe in den Tod“ etc. Daß gleich so viel ist, wenn man ein Kind in die Taufe steckt, als wenn ich spräche: Ich ersäufe und würge an dir alles, was Fleisch und Blut ist; daß, wenn man es heraus zeucht, ist es ohne alle Sünde und Unglück, also, daß in der Taufe untergehen muß alles, was weltlich und fleischlich ist, auf daß allein der Geist lebe. Also hält gegen einander die zwei: dort ersaufen die Menschen, die Fleisch und Blut sind; hier ersäuft geistlich alle Sünde und alles Unglück, was der Mensch mit sich bringt. So reimt es sich auf die Taufe.\footnote{Wittenberger: „Tag“, was in der ersten Ausgabe steht.} Darauf geht der Spruch im Psalm: Dominus diluvium inhabitare facit,
\quellseite{86}
der Herr sitzet, eine Sündfluth anzurichten [Ps. 29, 10.], nämlich eine bleibende Sündfluth, die nicht vergehe bis an der Welt Ende.

36. Was bedeutet aber nun die Arche? Sie ist also und darum gemacht, daß Noah darein thun sollte allerlei Thiere, beiderlei Geschlechtes, reine und unreine, auf daß Same auf Erden bliebe, und sich die Welt mehren könnte. Das ist, die christliche Kirche, die Versammlung aller deren, die getauft werden, berufen zum christlichen Glauben, so das Evangelium hören, sind alle in dem Kasten, sind aber noch wenig gegen dem andern großen Haufen. Denn ich sage nur von denen, so das Evangelium hören und rechtschaffen haben, nicht die allein den Namen haben, daß sie Christen heißen, sondern da keine falsche Predigt ist. In derselben Gemeine sind allerlei Thiere, Sie und He, darnach Noah selb achte; das ist, in der Predigt von der Taufe im Evangelio sind verfaßt mancherlei Menschen, nicht allein so da glauben, sondern auch die nicht glauben; haben dennoch alle das Evangelium, das sie hören und wissen. Denn die andern, die es nicht hören, sind nicht im Kasten, sondern ersaufen im Wasser.

37. So sind nun unter diesen etliche rein, und etliche unrein. Was das ist, werden wir hernach hören, jetzt ist es genug so viel davon, daß die christliche Gemeinde der Art ist, daß sie nicht eitel Heilige hat; davon ich oft gesagt habe. Und wäre gut, daß man ein Sprichwort daraus machte, denen zu antworten, so sich an unsrer Schwachheit ärgern, und sagte: Hat doch der Kasten Noah beide, unreine und reine Thiere, haben müssen; daß so viel gesagt sei: Wollen wir Christen sein, müssen wir uns darein begeben, daß wir viel gebrechliche Brüder und Schwestern unter uns haben müssen, auf daß wir uns gegen einander gleichen, und einander tragen, und sprechen: Lieber, wir werden nicht alle rein sein, es müssen doch im Kasten unter einander sein Raben und Tauben, Wolf und Schaf. Daher auch die Sophisten gesagt haben: Omnis anima in arca Noae, allerlei Seelen sind in dem Kasten Noah. Wäre wohl gesagt, wenn sie es verstanden hätten. Solches siehst du auch im Evangelio, daß Christus mit Zöllnern und Sündern umgeht [Matth. 9, 10. 11, 19. Luc. 5, 29.], läßt die andern fahren, die nur reine Thiere wollten haben, und waren es doch selbst nicht. So muß es in der Christenheit stehen.

[Es] steht auch darinne die größte Kunst, höchste Weisheit und Tugend, so die Christen haben. Wer die nicht kann, der halte sich nicht für einen Christen.

38. Weiter wird die Arche gemacht von glattem und gehobeltem Tannenholze. Damit aber bezeichnet ist der Christenheit höchste Tugend, die sie hat, daß sie sanftmüthig, gelinde ist, wie ein fein schlecht, glatt Holz, das ohne Aeste ist; ist keine Ungeduld, Zorn, Unwille, Neid darunter, nicht rauh, knorrig, noch ungebrochen; kürzlich, eitel Liebe, daß eines dem andern allerlei gut halten kann.

39. Ueber das soll sie auswendig und inwendig verpicht sein, oder verleimt. Das ist beides Liebe und Geduld. Wenn man ein Schiff nicht verpicht, kann es bald spalten und zerreißen. Nun ist Pech ein unflätig, schwarzes Ding, das nicht schmückt, aber wohl bewahrt. So ist auch das Leiden nicht lieblich anzusehen, es hält aber und bewahrt wohl. Summa, es ist Geduld und das heilige Kreuz, davon Paulus [Röm. 5, 4.] sagt: Patientia probationem operatur. Geduld bringet Erfahrung. Es ist auf allen Seiten Anfechtung; oben hat es hinein geregnet, und das Wasser daran geschlagen; nun ist es auswendig dagegen überall mit Pech verwahrt, daß es keinen Schaden nähme. Das soll sein Demuth und Geduld, daß man leide alles, was uns anficht, inwendig und auswendig.

40. Darum ist es in der Summa alles dahin gerichtet, daß die höchste Tugend ist, leiden und tragen alle Gebrechlichkeit unserer Brüder, darnach alle Anfechtung des Teufels und Todes. (Es ist nicht das geringste Leiden an\footnote{Jenner: an den. — Der Sinn ist: Das Leiden, welches uns von gebrechlichen Leuten zufließt, ist nicht ein geringes.} gebrechlichen Leuten; wir wollen immer gerne bei heiligen Leuten sein, die nicht wunderlich sind, sondern sich fein nach uns lenken, daß es immer gehe nach unserm Sinn. Darum sind viel gewesen, die daher und dorthin sind gelaufen; meinten, sie wollten einen Ort finden, da es nicht ungleich zuginge; hat sich ein jeder wollen aus der Schleife ziehen, und immer nach der Schnur haben, daß jedermann nach seinem Willen lebte, und er niemand nachgeben dürfte. Eben als jener Narr, der in der Sonne stand, und sich krümmte, und warf zornig, daß sich der Schatten auch krümmte. So wollen wir andere
\quellseite{87}
ich unrecht haben, und richten uns selbst nicht, machen uns damit unser Leben selbst sauer.

41. Weiter ist hier auch zu sehen, was die Thüre in der Seite, das Fenster über der Thüre, und die zweischichtigen und dreischichtigen Gemächer unten am Boden des Kastens bedeuten. Nun habe ich oft gesagt, man solle vor allen Dingen bei der Schrift bleiben im einfältigen Verstande, so die Buchstaben geben, lauter und rein; wenn das geschehen, mag man darnach mit Figuren und Deutung spielen; wie wir zum ersten nach den einfältigen Worten hin diese Historie gehandelt. Auch soll man solche Deutung nicht führen als gewiß und sicher, es sei denn, daß man aus der Schrift beweisen könnte, daß sie solches soll deuten. Wo das nicht ist, magst du wohl hin deuten; aber es soll sich niemand darauf verlassen. Zu mehren ist es nicht, daß einer nach seinem Geiste spiele; aber was den Glauben lehren soll, muß so gegründet und gewiß sein, daß man auch das Leben darf darüber lassen.

42. So haben wir oben [§ 35] aus St. Petri Spruche [1. Ep. 3, 21.] gegründet, daß die Sündfluth die Taufe bedeute; aber davon haben wir keinen Spruch, daß die Thüre auf der Seite deute die Wunde in der linken Seite Christi (wie man es bisher gedeutet hat), oder auch, daß die Arche den Leib Christi deute; darum darauf nicht zu bauen ist. Es ist wohl an ihm selbst wahr, daß dadurch die Deutung angezeigt wird; ob aber auch die Deutung recht sei und sich dazu reime, kann man nicht für einen Artikel des Glaubens haben. Wenn du es aber so deutest, daß der Kasten die christliche Kirche sei, die Thür aber bedeute das Wort, dadurch man in die christliche Gemeinde kommt, wie auch St. Paulus pflegt die Predigt zu heißen, als 1 Cor. 16, 9.: „Die Thür ist mir aufgethan, und sind viel Widersacher“; desgleichen auch Christus im Johanne [Cap. 10, 9.] sagt, er sei die Thür, dadurch man in den Schafstall aus- und eingeht; das nehme ich an, als der Schrift nicht ungemäß. Wer es aber nicht annehmen will, lassen wir seines Sinnes walten; denn wir müssen bei dem Hauptverstande bleiben, so der Buchstabe gibt; das andere mögen wir schmücken und ausstreichen, wie wir können.

43. Item, daß das Fenster oben angemacht ist, ohne Zweifel, wie [§ 21] gesagt, nicht dazu, daß sie Licht davon hätten; denn wie sollte ein groß Gebäu ein Licht empfahen von solchem kleinen Loch? Darum Noah in der Arche ohne Licht hat müssen bleiben, und ihm Licht machen mit Lampen, wie wir in der Nacht. Man wollte denn sagen, daß GOtt dem Licht so viel Glanz gegeben habe zu dem Kasten, daß es sie [die Arche] gar erleuchtet habe; ist aber nicht wohl zu glauben. Auch, hätte er gewollt, daß der Kasten vom Fenster erleuchtet würde, hätte er ihrer wohl mehr machen lassen. Was nun das deute, gebe ich auch einem jeglichen heim. Achte, was du willst, wir wollen es nach dem Buchstaben bleiben lassen, daß er es darum lasse machen, daß man könnte heraus sehen, wie es am Himmel stände, wann es Tag oder Nacht wäre, die Wolken aufhöreten, klar oder finster würde. Und hat ihm eben ein Maß bestimmt, daß es nicht größer denn einer Elle sein sollte.

44. Das wollte ich nun deuten auf dies, so Paulus scientiam, das ist, Erkenntniß heißt, daß die Christenheit, die inwendig lebt vom Licht der Gnaden, nicht vom Licht der Welt, sieht und richtet doch auch auswendig, was gut oder böse ist, wie es zugeht nach dem äußerlichen Wesen; was böse ist, daß man sichs nicht annehme; was recht und gut ist, daß man sich darnach richte, und GOtt danke. Das sei meine Deutung; hat jemand ein anderes, [der] mag ihm folgen; denn ich habe deß nicht gewissen Grund aus der Schrift.

45. Daß er nun weiter sagt von dem Untertheil, zweischichtig und dreischichtig zu machen, dabei ist zu sehen, wie der Kasten inwendig gemacht ist zu unten, innen ist er im Wasser gestanden, da läßt er ein Gebäu für allerlei Thiere; [es] müssen viel Gemächer mit vielen Thüren gewesen sein, denn die Thiere hat er alle müssen weiden, und jegliches sonderlich bei Paaren. Das will ich dahin deuten, daß in der Christenheit mancherlei Gaben sind. GOtt hat jeglichem seine Gabe geschenkt, damit er der Christenheit diene; davon Paulus schreibt Röm. 12, 6. 1 Cor. 12, 28. 29. Eph. 4, 11.: „Er hat etliche zu Aposteln gesetzt, etliche zu Propheten, etliche zu Evangelisten, etliche zu Hirten und Lehrern“ etc. Diese Deutung habe ich auch nicht aus der Schrift, doch reimt sie sich wohl mit der Schrift. Das sei genug von den Deutungen. Spricht nun Moses weiter:

\begin{quote}
V. 17—22. Denn siehe, ich will eine Sündfluth mit Wasser kommen lassen auf Erden, zu verder-
\end{quote}
\quellseite{88}
ben alles Fleisch, darinnen ein lebendiger Odem ist unter dem Himmel; alles, was auf Erden ist, soll untergehen. Aber mit dir will ich einen Bund aufrichten. Und du sollst in den Kasten gehen mit deinen Söhnen, mit deinem Weibe, und mit deiner Söhne Weibern. Und du sollst in den Kasten thun allerlei Thiere von allem Fleisch, je ein Paar, Männlein und Fräulein, daß sie lebendig bleiben bei dir. Von den Vögeln nach ihrer Art, von dem Vieh nach seiner Art, und von allerlei Gewürm auf Erden nach seiner Art; von den allen soll je ein Paar zu dir hinein gehen, daß sie lebendig bleiben. Und du sollst allerlei Speise zu dir nehmen, die man isset, und sollst sie bei dir sammeln, daß sie dir und ihnen zur Nahrung da seien. Und Noah that alles, was ihm Gott gebot.

46. Also geht Gottes Werk immerdar, daß es närricht ist anzusehen, auf daß er je die Welt zu Schanden mache. Konnte er nicht so viel verschaffen, da er im Sinn hatte, Noah zu erhalten in solcher greulichen Sündfluth, daß er nicht dürfte so viel Speise eintragen? Mußte er erst solchen Vorrath sammeln von zeitlichem Gut, als könnte sie Gott nicht lange genug speisen und nähren? [Er] hebt an und läßt Noah sorgen, daß er Essen und Trinken habe. Warum thut er nun das? Erstlich ist keine andere Ursach, denn sein Wille. Aber doch ist es auch geschrieben, uns zu warnen vor Vermessenheit, die da heißt Gott versuchen. Wahr ist es, daß Gott will, daß alle unsere Werke im Glauben gehen sollen, daß man nichts thue ohne Glauben; aber dennoch will er auch nicht, daß man das lasse anstehen, das vorhanden und zuvor von ihm gegeben ist.

47. Als, daß wir es in andern Exempeln sehen: Er hat lassen Brücken über die Wasser bauen, darüber ein jeglicher sicher gehen könne. Wenn du nun auf den Glauben wollest\footnote{In den Ausgaben: wollest.} so trotzen, und nicht über die Brücke, sondern auf dem Wasser gehen, würdest du billig ersaufen und zum Teufel fahren. Denn das wäre gethan, das er nicht geboten noch befohlen hat. Es ist wohl beides sein; aber das hat er gegeben, jenes hat er nicht gegeben. Darum, weil du nicht im Geist fühlst, daß du Ursach hast, warum du müssest auf dem Wasser gehen, sollst du es nicht thun.

48. Also hat auch Christus gethan, wie im Matthäo [Cap. 4, 7.] steht. Als ihn der Teufel auf des Tempels Spitze führte, daß er sich herab sollte lassen, sprach er: „Es steht geschrieben, du sollst Gott nicht versuchen“ [5 Mos. 6, 16.]. Dergleichen wäre, wenn du wolltest müßig gehen und nicht arbeiten, und sagen: Ich glaube, Gott wird mich ernähren. Nicht also. Es ist wahr, daß er alles gibt, alle Dinge ernährt und erhält; aber daß du nicht wolltest brauchen, das du wohl brauchen kannst, wäre Gott versucht. Denn er will, daß du brauchest, was du vor dir hast liegen, das dir schon gegeben und vorhanden ist; nicht, daß du das Maul gen Himmel aufsperrest, und die Creatur fahren lässest, die er dir gegeben hat. Er wird um deinetwillen ohne Noth kein Wunder thun.

49. Wo sich’s aber also begibt, daß es aus deiner Gewalt kommt, daß du nicht arbeiten könntest, kein Korn, Geld noch Barschaft hättest; da mußt du dennoch glauben, Gott werde dich ernähren und speisen. Also, wenn du gerne über die Brücke gehen wolltest, aber nicht könntest, und doch gebrungen würdest, über das Wasser zu gehen, so gehe in Gottes Namen in starkem Vertrauen.

50. Item, kurzum ist es um den Glauben also gethan, daß er beide im Vorrath und Mangel auf rechter Mittelstraße fährt, also, daß er des Vorraths braucht mit Dankbarkeit, wo es Gott gibt; auch wiederum, wo Mangel ist und doch fest trauet, durch solche Zuversicht im Mangel keine Noth leidet.

51. Item, desgleichen ist auch von der Keuschheit zu reden, ja, von allen Sachen. Es ist ja eine große Gabe die Jungfrauschaft, aber die vom Himmel kommt; niemand ists angeboren. Denn Gott hat es so gemacht, daß dies ein Männlein, dies ein Mägdlein sei. Wenn du nun das willst fahren lassen, daß du im Vorrath hast und wohl thun kannst, und in Himmel gaffen, nach der hohen, seltenen Gabe warten, wirst du fehlen; wie denn jetzt alle Mönche und Pfaffen Gott versuchen, ja, schänden mit ihrer gelobten Keuschheit. Wenn es des gelten sollte, möchtest du auch sagen: Was darf man Gottes Wort in der Gemeinde mündlich predigen, oder sonst hören und lesen; ich will lieber warten, bis Gott selbst mit mir redet vom Himmel; er wird wohl ohne Wort mein Herz erleuchten und anzünden, und mir den Geist geben etc.; welches er doch verboten hat.
\quellseite{89}

52. Warum ist das Gottes Wille, daß wir deß, das er geschaffen oder geordnet, brauchen, und nicht in Wind schlagen; du wirst dennoch zu glauben genug haben. Das will dies Stück lehren, und ist fast noth, daß mans wohl wahrnehme. Da stellt sich Gott also, daß ihn die Weisen wohl möchten für einen Narren halten, daß er sich solches Dinges annimmt; er thut es aber eben darum, daß er dieselbigen zu Narren mache in ihrer Klugheit, wie Paulus sagt [1 Cor. 3, 19. 20.]. Folgt nun:

\section{Das siebente Capitel}

\begin{quote}
B. 1. Und der Herr sprach zu Noah: Gehe in den Kasten, du und dein ganzes Haus.
\end{quote}

1. Da siehst du das rechte Urtheil des allmächtigen Gottes, das schrecklich zu hören ist, aber viel schrecklicher gemeinen anzusehen. St. Petrus zeucht es auch an, und beweget es hoch: „Gott hat der vorigen Welt (spricht er [2. Ep. 2, 5.]) nicht verschonet, sondern bewahrete Noah, den Prediger der Gerechtigkeit, selb achte, und führte die Sündfluth über die Welt der Gottlosen.“ Alles sollte er sagen: Dieselbige Welt war dieselbige Zeit viel besser denn jetzt, noch hat er sie vertilgt; wie schrecklich wird er denn jetzt, am Ende der Welt, strafen?

2. So siehe nun den Noah an, wie er da mitten in der Hölle steht. Es ist eine schlechte Historie anzusehen, aber so gewaltig, daß [es] niemand genug erlangen kann. Wenn es\footnote{In den Ausgaben: er.} heutiges Tages so käme, daß wir also führen, wie er in dem Kasten, wer würde sein, der nicht verzweifelte? (Es ist der Vernunft und Sinnen zu greulich anzusehen, daß sie so müssen schweben, haben keinen Aufenthalt, [es] soll oben und unten eitel Wasser herplatzen und schlagen.

3. Zum andern sehen sie dazu nichts vor Augen, wo sie sich hinwenden, denn eitel Sterben und Tod; da sollte erst einem das Herz zappeln. Denn in solchem greulichen Zorn, da Gott so streng fährt, daß er um der Menschen willen auch seines Thieres verschont, möchten sie denken, Gott würde auch unter sie schlagen, sonderlich, weil sie das Wasser rings herum sahen schlagen.

4. Darum siehe, was das für ein Glaube gewesen ist, der in solchem Ansehen des grausamen Zorns hat können stehen. Da ist ein rechter Kampf mit dem Glauben und Unglauben gewesen, und muß ihr Herz manchen starken Puff erlitten haben. Da sehen sie noch wissen keinen Rath, daß sie heraus kommen und sicher seien, weder über sich noch um sich, müssen aushalten vierzig Tage lang. Wer ihm nun da hätte selbst wollen helfen, hätte müssen verzweifeln; denn da war keine Hülfe zu sehen, mußten allein am bloßen Worte hangen, und durch den Glauben wider alle Sinne und Vernunft fechten. So siehst du, welche große, mächtige Gewalt des Wortes ist, wenn es mit dem Glauben gefaßt wird, daß sie fünf Monden lang den Tod ohne Unterlaß vor Augen hatten, und doch verachten konnten.

5. Darum ist es zu Preis des Glaubens geschrieben, wie denn die ganze heilige Schrift dazu gegeben ist, den Glauben zu stärken. Das ist der Fische einer, damit Christus das Volk in der Wüste speisete, und der Ochs oder Mastvieh, zu der Hochzeit geschlachtet, wie im Matthäo [Cap. 22, 4.] steht. Solches soll man predigen und aufwerfen; da kann kein Werk noch einiges Vornehmen helfen, da muß man nackend allein am bloßen Worte hangen, so Gott gesagt hat.

\begin{quote}
Denn dich habe ich rechtfertig ersehen vor mir zu dieser Zeit.
\end{quote}

6. Das war so stark, daß es sie alle mit der Arche im Tode erhielt, die sonst nicht Einen Tag hätten mögen bleiben. Der Glaube mußte lebend bleiben, wenn auch eine ewige Sündfluth gewesen wäre. So haben sie sich frei dahin müssen mitten in Tod geben, daß sie verloren wären mit Leib und Seele, wo das Wort nicht wäre gewesen.

7. So lasset uns nun auch lernen, wenn die Stunde kommt, daß man sterben soll, und der Tod vor Augen ist, und schreckt uns mit seinem Blicke Teufels Hinterlist und Gottes Zorn, daß du nicht denkest, du müßest untergehen, nicht um dich tappest und sehen wollest, wo du bleiben und hintreten sollst (wie die thun, so da sagen: Ich sterbe und weiß nicht wann; ich fahre, und weiß
\quellseite{90}
nicht wohin). Da mußt du Augen und alle Sinne zuthun, nichts wissen noch hören wollen, denn was Gottes Wort sagt, nicht achten, was du fühlest, oder je dasselbe überwinden. Das Wort ergreife und lasse dir es nicht nehmen, daß du sagest: Hier bin ich in Todesnoth und Angst, ich weiß aber, daß ich getauft bin, und mir Gott das und das verheißen hat; rufe ihm also sein Wort auf, so sehr und hart der Tod herein bricht.

8. Da wird nun nicht gelten, daß man sage: Ach! wer nun Gott gedient und viel gute Werke gethan hätte. Es sind eitel tolle, unsinnige Worte, damit der Teufel die armen Leute nur verzagt macht. So sollte man aber sagen: Ich habe gesündigt und viel Böses gethan, das ist mir leid; aber du bist ein solcher Gott, der nicht ansieht, wie fromm oder wie böse man ist, wenn man nur auf deine Güte sieht und trauet. Also kann man im Tode genesen und bleiben, wie die Acht mitten in der Sündfluth geblieben sind. Das ist es, das man hieraus soll lernen. Das ist nun die Zusagung, Noah gegeben: „Gehe in den Kasten, du und dein ganzes Haus“. Und sind eitel lebendige Worte, darinne sie viel tiefer erhalten sind, denn in dem Kasten. Denn hätte er das nicht gehabt, so hätte er müssen sinken, tiefer denn unter die Erde, wo es sein könnte. Uns aber sind sie zum Exempel geschrieben, den Glauben zu stärken und [zu] trösten.

9. Und hier muß man gewohnen, daß die Schrift heißt ein „Haus“, da wir Weib und Kind“ sagen. Er schließt aber aus das Gesinde, begreift allein, was Eines Fleisches und Blutes und der Sippschaft ist. Also liest man der Weise nach im andern Buch Mosis [Cap. 1, 17. 21.] von den Wehmüttern: „Weil sie Gott fürchteten, und der Juden Kinder leben ließen, machte er ihnen Häuser“, das ist, er gab ihnen Männer und Kinder; nicht, daß er ihnen von Holz und Steinen Häuser gebauet hätte. Also wirst du finden hernach in diesem Buche, Cap. 30, 3., wie Rahel, als sie nicht Kinder hatte, ihrem Manne Jakob ihre Magd gab, und sprach: „Lege dich zu ihr, auf daß ich doch durch sie erbauet werde“, das ist, daß mir uns züchtigen und mehr werden im Hause. So sagt er hier: „Du und dein ganzes Haus“, das ist, dein Weib, drei Söhne und ihre Weiber.

\begin{quote}
V. 2. 3. Aus allerlei reinem Vieh nimm zu dir je sieben und sieben, das Männlein und sein Fräulein; von dem unreinen Vieh aber je Ein Paar, das Männlein und sein Fräulein. Desselben gleichen von den Vögeln unter dem Himmel, je sieben und sieben, das Männlein und sein Fräulein, auf daß Saame lebendig bleibe auf dem ganzen Erdboden.
\end{quote}

10. Das ist ein wunderlicher Text, daß ich nicht wüßte, warum es gesagt wäre, wenn nicht mysteria, das ist, geistliche Deutungen, darinne wären. Oben [Cap. 6, 20.] hat er gesagt von Paaren und Paaren, wie auch hernach in diesem Capitel, und spricht doch hier mitten im Texte, er solle einen Unterschied machen unter den reinen und unreinen, und sonderlich drückt er die Vögel aus, daß er je sieben und sieben darein nehmen solle.

11. Zum ersten sieht man hier klärlich, daß des Gesetzes Mosis viel genommen ist aus den alten Historien der Patriarchen. Denn so liest man im dritten Buche Mosis [Cap. 11, 3. 9. 10.], daß Gott einen Unterschied macht unter reinen und unreinen, die man essen oder nicht essen soll, und sagt, was das Zeichen sei der reinen Thiere, nämlich, was die Klauen mitten entzwei spaltet, und wiederkäuet; item, von den Fischen, welche Schuppen und Floßfedern haben etc. Dies hat schon hier Noah gehalten, und die Thiere unterschieden. Damit Moses anzeigt, wie das Gesetz bereits in vielen Stücken ganghaftig ist gewesen bei den ältesten Vätern, das er hernach zusammen getragen, gemehrt und gebessert hat durch Gottes Befehl.

12. Zum andern, was nun die Ursache sei, warum er hier je sieben und sieben heißt nehmen von den reinen Thieren und Vögeln, und oben am sechsten Capitel, V. 20., nur ein Paar von allerlei Vieh und Vögeln, lassen wir fahren bis in die mysteria. Denn nach dem schriftlichen Sinne ist genug die Meinung, daß der unreinen Thiere nicht zu viel in dem Kasten wären. Darnach, als Noah hinein kam, sollte er vielleicht einmal oder zwei opfern; wie es denn schon ging, wie hernach bei Mose im Gesetze. Sollte er nun nur Ein Paar reiner Thiere mitgenommen haben, wäre nichts Uebriges geblieben, damit ein lebendiger Saame wäre erhalten.

\begin{quote}
V. 4. Denn noch über sieben Tage will ich regnen lassen auf Erden vierzig Tage und vierzig Nächte, und vertilgen alles, das das Wesen hat, von dem Erdboden, das ich gemacht habe.
\end{quote}
\quellseite{91}

13. Das ist die letzte Predigt, die Noah auf Erden gethan hat vor der Sündfluth, als er nun alle Dinge hatte eingetragen und zugerichtet. Aber die Leute waren es nun gewohnt, daß er sich so lange mit ihnen bläuete,\footnote{sich bläuen = sich abmühen, plagen.} daß es nicht halb umb sie nur verdroß worden. Darum beschleußt Moses:

B. 5—10. Und Noah that alles, was ihm der HErr gebot. Er war aber sechshundert Jahr alt, da das Wasser der Sündfluth auf Erden kam, und er ging in den Kasten mit seinen Söhnen, Weibe, und seiner Söhne Weibern, vor dem Gewässer der Sündfluth. Von dem reinen Vieh, und von dem unreinen, von den Vögeln, und von allem Gewürm auf Erden, gingen zu ihm in den Kasten bei Paaren, je Männlein und Fräulein, wie ihm der HErr geboten hatte. Und da die sieben Tage vergangen waren, kam das Gewässer der Sündfluth auf Erden.

14. Droben [Cap. 5, 32.] sagt er, er sei fünfhundert Jahr alt gewesen, jetzt sagt er sechshundert, daß er eben hundert Jahre gepredigt hat, und den Kasten zugerichtet, daß sie Zeit genug hatten, sich zu bessern. Sie haben zu derselben Zeit lange gelebt, darum hat ihnen Gott auch so lange Frist gegeben. Jetzt gibt er nicht so lange Frist, weil wir so lange nicht leben. Das hat sie auch ärger gemacht, und verhärtete sie auch, daß sie gedachten: Gott gibt uns so lange Frist, läßt es uns wohl gehen, das ein Zeichen ist, daß er uns gnädig ist, und noch nicht Noth hat; wenn er so zornig wäre und so greulich plagen wollte, ließe er uns so lange nicht leben; wie denn alle Ungläubigen gedenken. Aber Gott läßt sie nur dadurch verblendet und verstockt werden, daß sie es ja wohl verdienen.

15. Da es ihnen aber auf den Hals kam, da flohen sie, und konnten nirgend aus. Da wollte Gott nicht mehr hören bitten, denn sie seiner Predigt gehorcht hatten; sie hatten ihn lassen umsonst rufen, so läßt er sie auch wieder umsonst rufen [Micha 3, 4.]. Also sollen wir hierin lernen, beide die Langmüthigkeit und den Zorn und Strafe Gottes.

B. 11. In dem sechshundertsten Jahre des Alters Noah, am siebenzehnten Tage des andern Monden.

16. Deß muß man gewohnen, wer in Mose liest, daß er feindlich\footnote{feindlich = gar sehr.} viel Worte führt, damit er immer Ein Ding treibt, und zuweilen über die Maße reich ist, und daher schwemmt mit Worten; wiederum auch zuweilen so krank, daß er kaum tröpfelt, und viel Dinges mit Einem Wort ausrichtet, daß wohl möchte verdrießlich sein einem, der der Sprache Art nicht weiß. Aber Gott läßt es eben darum so geschrieben werden, daß er die Vernunft zum Narren mache; denn wo er am närrischten redet, da ist er am klügsten; darum ist es nicht zu verachten. Er thut es aber hier in dem Capitel sonderlich darum, daß Noah Macht daran gelegen war, da er in solcher Angst stund. Wenn man in Nöthen liegt, ist es nicht zu viel, daß man ein Wort drei- oder viermal hört, da kann man nicht genug predigen und trösten; wenn man aber gute Ruhe hat, wird man bald verdrossen, als hör einer deuten möchte: Du Narr, was machst du so viel Worte? darum, daß es scheint, wie er möchtig sei; macht freilich nichts Anderes, denn daß wir nicht in dem Sinne stecken, den er hier gehabt hat.

B. 11—24. Das ist der Tag, da aufbrachen alle Brunnen der großen Tiefe, und thaten sich auf die Fenster des Himmels, und kam ein Regen auf Erden, vierzig Tage und vierzig Nächte. Eben am selben Tage ging Noah in den Kasten mit Sem, Ham und Japhet, seinen Söhnen, und mit seinem Weibe, und seiner Söhne dreien Weibern. Dazu allerlei Thier nach seiner Art, allerlei Vieh nach seiner Art, allerlei Gewürm, das auf Erden kreucht, nach seiner Art, und allerlei Vögel nach ihrer Art; alles, was fliegen konnte, und alles, was Fittig hatte, das ging alles zu Noah in den Kasten bei Paaren, von allem Fleische, da ein lebendiger Geist inne war; und das waren Männlein und Fräulein von allerlei Fleische, und gingen hinein, wie ihm denn Gott geboten hatte; und der HErr schloß hinter ihm zu. Da kam die Sündfluth vierzig Tage auf Erden, und die Wasser wuchsen, und huben den Kasten auf, und trugen ihn empor über der Erde. Also nahm das Gewässer überhand, und ward sehr auf Erden, daß der Kasten auf dem Gewässer

Lutheri Werke. Bd. III.
\quellseite{92}
fuhr. Und das Gewässer nahm überhand, und wuchs so sehr auf Erden, daß alle hohen Berge unter dem ganzen Himmel bedeckt wurden; fünfzehn Ellen hoch nahm das Gewässer überhand über die Berge, die bedeckt wurden. Da ging alles Fleisch unter, das auf Erden kreucht, an Vögeln, an Vieh, an Thieren, und an allem, das sich regte auf Erden, und an allen Menschen; alles, was einen lebendigen Odem hatte im Trocknen, das starb. Also ward vertilgt alles, was auf dem Erdboden war, vom Menschen an, bis auf das Vieh, und auf das Gewürm, und auf die Vögel unter dem Himmel, das ward alles von der Erde vertilget. Allein Noah blieb über, und was mit ihm in dem Kasten war. Und das Gewässer stand auf Erden hundert und fünfzig Tage.

17. Da sagt er, wo die Wasser herkommen, nämlich von unten und oben. Was sind denn die Brunnen der Tiefe, und Fenster des Himmels? Es ist mit verblümten Worten geredet, und so viel gesagt: Wo Tiefe war, als große Teiche und Seen, da ging es aus, quellte und ergoß sich, daß es daher schwemmete; also, daß die Brunnen sind, daher es quellte, und unten heraus ist gebrochen; oben aber ist es hereingebrochen mit eitel Wolkenbrust,\footnote{Wolkenbrust = Wolkenbruch.} und mit Gewalt greulich geschwemmt, hat plötzlich viel Wasser gegeben; denn da gehört viel Wasser zu, daß es fünfzehn Ellen über alle Berge gegangen ist, wie er hier schreibt.

18. Von der heimlichen Deutung dieses Capitels will ich nicht viel sagen; denn was die Sündfluth deutet, ist oben [Cap. 6, § 35] genugsam gehandelt; aber von den reinen und unreinen Thieren gehört in ein ander Buch Mosis, da geboten wird, von keinem unreinen zu essen, noch zu opfern. Denn dieselbige Deutung ist zu weit, gehört zum Evangelio und Gesetze. Wir haben es auch noch in vielen Stücken, daß wir etliche Thiere nicht essen, welcher Unreinigkeit nicht von Natur herkommt, sondern aus Gottes Gesetze. Also sind von Anfang der Welt immerdar etliche Satzungen gewesen. Denn es fast unmöglich ist, daß ein Regiment ohne mancherlei äußerliche Satzung bestehe. So haben schon die Väter Unterschied unter den Thieren gehabt, aber sind doch ohne Zweifel damit frei umgegangen, nach der Freiheit des Gewissens, wie auch viel Propheten hernach unter dem Gesetze gethan haben.

\section{Das achte Capitel}

V. 1—12. Da gedachte Gott an Noah, und an alle Thiere, und an alles Vieh, das mit ihm in dem Kasten war, und ließ Wind auf Erden kommen, und die Wasser fielen, und die Brunnen der Tiefe wurden verstopfet, sammt den Fenstern des Himmels, und dem Regen vom Himmel ward gewehret, und das Gewässer verlief sich von der Erde immer hin, und nahm ab nach hundert und fünfzig Tagen. Am siebenzehnten Tage des siebenten Mondes ließ sich der Kasten nieder auf das Gebirge Ararat. Es verlief aber das Gewässer fortan, und nahm ab, bis auf den zehnten Monden. Am ersten Tage des zehnten Monden sahen der Berge Spitzen hervor. Nach vierzig Tagen that Noah das Fenster auf an dem Kasten, das er gemacht hatte, und ließ einen Raben ausfliegen, der flog aus und kam wieder, so lange als das Gewässer vertrocknete auf Erden. Darnach ließ er eine Taube von sich ausfliegen, auf daß er erführe, ob das Gewässer gefallen wäre auf Erden. Da aber die Taube nicht fand, da ihr Fuß ruhen konnte, kam sie wieder zu ihm in den Kasten; denn das Gewässer war noch auf dem ganzen Erdboden; da that er die Hand heraus, und nahm sie zu sich in den Kasten. Da harrte er noch andere sieben Tage, und ließ abermal eine Taube fliegen aus dem Kasten, die kam zu ihm um Vespertzeit, und siehe, ein Oelblatt hatte sie abgebrochen, und trug es in ihrem Munde. Da vernahm Noah, daß das Gewässer gefallen wäre auf Erden. Aber er harrte noch andere sieben Tage, und ließ eine Taube ausfliegen, die kam nicht wieder zu ihm.
\quellseite{93}

1. Das Capitel hält inne, wie die Sündfluth hat abgenommen und hinweg gefahren ist, nachdem es eine greuliche lange Zeit gestanden und gewachsen hatte, und auch etliche Tage gestocket hat, bis es vergangen ist; hundert und fünfzig Tage hat es gestanden, sagt der Text, das ist beinahe ein halb Jahr. Da kannst du denken: sollte Noah in der Zeit mit den Seinen nicht gejagt haben und gezappelt, daß es nicht wollte aufhören? So greulich und schrecklich geht Gottes Wort herein; noch ist der Glaube viel stärker, der es kann aushalten, und des Endes erwarten. Das anzuzeigen, beschreibt er es hier so fleißig mit so viel Worten, daß er auch die Tage so eben zählt. Erstlich sagt er, es habe angefangen abzunehmen nach hundert und fünfzig Tagen, daß der Kasten im siebenten Monden auf dem höchsten Gebirge Ararat gestanden ist, so das Wasser zuvor fünfzehn Ellen hoch über alle Berge gegangen war. Danach hat es gefallen vierzig Tage, bis auf den zehnten Monden. Da harrte er noch bis in den andern Monden des folgenden Jahres. Also ist er eben ein Jahr und bei zweien Monden in dem Kasten geblieben. Das sei darum geredet, daß man sehe die Gewalt des Glaubens und Gottes Wunderwerk.

2. Danach setzt er zwei Stücke, von dem Raben und von der Taube. Und ist wohl wunderlich anzusehen, gerade als hätte sich nichts in dem Kasten begeben denn das. Von ersten ließ er einen Raben ausfliegen, spricht Moses, der flog aus und kam wieder, so lange das Gewässer vertrocknete. Das ist so viel geredet: er machte es so lange mit dem Wiederkommen, bis es gar trocken ward, daß er noch wieder kommen soll. Wie wir deutsch möchten sagen: Er kommt und kommt, er harrt und harrt; das ist, es wird so lange, bis er gar außen bleibt.

3. Was gelüstet aber den Heiligen Geist, solch Ding zu reden? Es ist eine lose Historie, menschlich zu rechnen. Aber Gott richtet nicht, wie Menschen; er ist der Meister, der weiß, was er soll schreiben lassen. Summa Summarum: Die Historie hat das in sich, daß Noah natürlich braucht des Raben; denn weil er der Art ist, wo er auf ein faul Aas kommt, sitzt er darauf und frißt; [so] dachte Noah, er würde sich voll fressen und wieder kommen. Er wollte aber nicht, und blieb außen.

4. Danach schickte er die Taube aus, die sollte wieder kommen und sagen. Das ist ein rein Thier, ißt kein Aas noch Unreines; daher ein Sprichwort ist: Es ist so rein, als hätten es die Tauben erlesen. Die kommt nun wieder, und bringt nichts. Die andere bringt ein grün Oelblatt; die dritte bleibt außen.

5. Wenn man dies will deuten, kann man es nicht besser thun, denn auf die Prediger, denen da befohlen ist, Gottes Wort zu predigen. Die christliche Kirche, durch den Kasten Noah bedeutet, hat zweierlei Prediger: zum ersten die Bauchdiener und Freßlinge, die sind durch den Raben bedeutet, die durch das Predigen nichts suchen, denn daß sie genug haben, und ihren Bauch füllen; wo das Wort nicht will dienen zu ihrem Bauche, da drehen sie ihm eine Nase, daß es sich dazu lenken muß. Dazu, wenn es zum Treffen geht, daß sie es sollen bekennen, so leugnen sie; das sind unreine Schelme, die Aas fressen; die müssen wir auch haben. Denn so ist es gegangen zu allen Propheten- und Apostel-Zeiten bis hieher, daß neben der rechten Lehre ausliegen falsche Lehrer, dazu soll man diesen Raben hoch heben, daß er predige das rechte Gottes Gesetz, wiewohl er’s nicht recht auslegt.

6. Unsere Lumpenprediger, die allein Menschentand predigen, sind noch nicht so gut. Nun, diese kommen auch aus dem Kasten, haben den Namen, daß sie in die Christenheit gehören, sind Prediger und Bischöfe. Was thun sie aber? Fressen eitel unreines, todtes Aas; das ist, ihre Lehre ist ganz von Fleisch und Blut und äußerlichen Werken; da ziehen sie die Schrift hin, da bleiben sie inne, kommen nicht wieder, wie Paulus [1 Tim. 5, 12. 6, 4. 5. Tit. 3, 11.] sagt: „Sie haben ihr Urtheil dahin.“ Der Rabe bringt kein gut Zeugniß, noch nichts Tröstliches; solche Lehrer predigen, da niemand seine Freude von hat, machen eitel Angst des Gewissens.

7. Die aber durch die Taube bedeutet werden, sind, die einfältig daher gehen in Gottes Wort; wie St. Paulus [2 Cor. 1, 12.] seine Predigt rühmt und preist, und spricht: „Wir haben gewandelt und gehandelt in rechter Einfältigkeit und Lauterkeit, nicht mit fleischlicher Weisheit“; das ist, wir Prediger haben so einen einfältigen Sinn, wie die Tauben, suchen nichts Anderes, denn die reine Lehre des Evangelii, begehren weder Ehre noch Gunst, wagen alles daran, daß nur darum da zu thun ist, daß die Lehre lauter bleibe, es gehe ihnen, wie Gott
\quellseite{94}
will. Die andern aber sind nicht einfältig, sondern stecken voll eigenes Muthes; das ist ein schalkhaftig Auge, wie Christus [Matth. 6, 23.] sagt. Das ist nun die liebe Taube, die fleugt erstlich aus, findet nichts, [es] ist noch nicht Zeit, daß man das Evangelium predige, man muß Gott lassen auszürnen, bis daß der Rabe sich voll Aas fresse; darum kommt sie wieder in Kasten.

8. Zum andern schickt er aber[mal] eine andere Taube aus, die kommt wieder und bringt ein Oelblatt im Munde zur Vesperzeit. Es hat die ganze Schrift allenthalben angezeigt, daß das Evangelium die letzte Predigt sein solle; darum man in der Schrift hin und wieder findet, daß der Abend immer bedeutet und anzeigt die letzte Zeit und Predigt, wie auch im Evangelio, da es Christus [Luc. 14, 16.] dem Abendmahl vergleichet. Wiewohl sich aber der Abend lange verzeucht, so ist es dennoch die letzte Predigt, darnach keine andere kommen soll. Darum kommt die Taube eben um Vesperzeit, als der Tag schier will ein Ende haben.

9. Was ist aber das Oelblatt in ihrem Munde? Nichts Anderes, denn das lautere Evangelium in dem Munde der reinen Prediger. Oel bedeutet durch die ganze Schrift Gnade und Barmherzigkeit. Denn es hat die Art, daß, wie der Wein den Leib innwendig geschicht und fröhlich macht, also macht Oel den Leib auswendig hübsch, gelenk, glatt, und geschicht. Weil es nun so fein gelinde macht, heilt und sanftet mit seinem lieblichen, feinen Safte, so wird dadurch bedeutet göttliche Gnade und Barmherzigkeit, gepredigt durch das Evangelium. Darum führt die Taube das Oelblatt im Munde; denn das Blatt bedeutet die Predigt, wie die Frucht die Werke. So siehe, was das Evangelium ist: nichts Anderes, denn wie das Oelblatt, das die Taube in den Kasten bringt.

10. Aus solchem Exempel und Figur sollten wir Christen billig ein täglich Sprüchwort machen, wider alle Menschenlehre und Gesetze, daß wir in unsern Reden ganghaftig sein ließen: Die Taube bringt nichts im Munde, denn das Oelblatt, das ist, man soll in der Christenheit nichts Anderes predigen, denn das lautere Evangelium. Sonst, wo etwas Anderes zu predigen wäre, würde sie irgend ein Blatt von einem Holderbusche,\footnote{Holder = Holunder.} oder ein Stück von einem faulen Aas bringen.

11. Als nun das geschehen ist, hat Noah gesehen, wie Moses im Anfange dieses Capitels sagt, daß Gott an ihn gedachte; da ist nun der Jammer aus, und der Zorn hinweg. Das ist nun die letzte Predigt, wie die Taube der letzte Bote ist, darnach keine mehr kommen soll. Wenn die ausgeht, haben wir ein gewisses Zeichen der Gnade, daß Gott den Zorn abwendet; da kommt der Trost, und hört die Anfechtung auf, und kommt das Herz hervor, vom Tode in das Leben. Bisher ist der fromme Noah immer in großem Kampfe gestanden, daß die Taube mit dem Oelblatt freilich ein tröstlicher Bote gewesen ist.

\begin{quote}
B. 13. 14. Im sechshundertsten und einem Jahr des Alters Noah, am ersten Tage des ersten Monden, vertrocknete das Gewässer auf Erden. Da that Noah das Dach von dem Kasten, und sahe, daß der Erdboden trocken war. Also ward die Erde ganz trocken, am sieben und zwanzigsten Tage des andern Monden.
\end{quote}

12. Da siehst du, daß sie ein ganzes Jahr im Kasten gewesen sind. Droben [Cap. 7, 11.] sagt er, wie er im sechshundertsten Jahr, am siebenzehnten Tage des andern Monden ist hineingegangen. Hier sagt er, daß die Erde im sechshundertsten und einem Jahre, am sieben und zwanzigsten Tage des andern Monden ganz trocken ist worden, daß Noah heraus gegangen ist. Am ersten Tage des ersten Monden, spricht er, vertrocknete die Erde; darnach, am siebzehnten und zwanzigsten Tage des andern Monden, ward sie gar trocken, das ist, er hat noch drei Monden geharrt nach der Taube Ausfliegen, und noch länger.

\begin{quote}
B. 15—19. Da redete Gott mit Noah, und sprach: Gehe aus dem Kasten, du und dein Weib, deine Söhne und deiner Söhne Weiber mit dir, allerlei Thier, das bei dir ist, von allerlei Fleische, an Vögeln, an Vieh und an allerlei Gewürm,
\end{quote}
\quellseite{95}
das auf Erden kreucht, das gehe heraus mit dir. Und reget euch auf Erden, und seid fruchtbar und mehret euch auf Erden. Also ging Noah heraus mit seinen Söhnen und mit seinem Weibe und seiner Söhne Weibern. Dazu allerlei Thiere, allerlei Gewürm, allerlei Vögel, und alles, was auf Erden kreucht, das ging aus dem Kasten, ein jegliches zu seines gleichen.

13. Das ist auch ein rechter Text, der zum Glauben dient, denn da siehst du, daß er nicht aus dem Kasten gegangen ist, noch gegangen wäre, ehe ihm GOtt solches befahl. Durch GOttes Gebot war er hinein gegangen, durch GOttes Gebot geht er wieder heraus, und nicht allein er, sondern alles, was da lebte auf Erden. Und siehe, wie es nach einander geht. Zum ersten gebeut GOtt ihm, seinem Weibe, Söhnen und ihren Weibern, und Thieren, daß sie heraus gehen. Darnach sagt er: „Reget euch, seid fruchtbar und mehret euch.“ Denn es war anzusehen, als sollte es gar aus sein, daß kein Leben mehr auf Erden würde bleiben; darum heißt er sie, sich wieder züchtigen. Und das ist nicht allein ein Gebot, sondern strenger Ernst, wie oben gehört ist.

14. Das ist aber hierinne angezeigt, welches durch den ganzen Mosen immerdar angezeigt wird, mehr denn in der ganzen Schrift, daß vor GOtt kein Leben noch Werk gilt, das er nicht befiehlt, entweder öffentlich oder heimlich durch Menschen. Darum läßt er Noah nichts thun, er habe es ihm denn befohlen; wiewohl es hier also gethan war, daß solches Werk geschehen mußte, daß er mit allem, was bei ihm war, aus dem Kasten ginge, weil er sah, daß GOtt die ganze Erde hatte lassen trocken werden; dennoch durfte er es nicht eher thun, ehe er es ihm hieß, auf daß er ja gewiß wäre, daß sein Werk ginge in GOttes Gefallen.

15. Demnach soll ein jeglicher seinen Stand und alle Werke also stellen, daß er gewiß sei, daß sie GOtt gefallen, und also leben, daß wir immer bereit seien zum Tode. Welches niemand thun kann, denn der deß sicher ist, daß sein Leben, und was er thut, GOtt gefalle. Denn wer das nicht sagen kann, ist in einem verdammlichen Stande. Daß wir aber deß gewiß werden, wird man nicht mit Werken ausrichten, sondern der Glaube, der macht den Menschen angenehm, und versichert das Herz, daß alles GOtt gefalle; und ob es gleich etwas thut, das ihm nicht gefiele, als oft geschieht, daß er ihm dennoch solches zugute hält. Denn des Menschen Leben steht also, daß es gar viel thut, das es nicht gerne wollte gethan haben; wie Paulus Röm. 7, 19. sagt. Es wird aber nichts Anderes daraus,\footnote{In den Ausgaben: nicht anders aus.} kann auch nicht anders werden, es müssen immer Sünden und grobe Gebrechen mitlaufen; darum wird der Mensch dem Leben feind, und begehrt immer des Todes. Also muß ein Christenmensch immer mit den Werken umgehen, die der Gnade und Schonens dürfen; er wollte gerne mit guten Werken umgehen, so kann er nicht.

16. Darum haben wir gelehrt, daß das ganze Leben im Glauben und GOttes Worte hergehen muß, soll es GOtt gefallen. Denn der Glaube ohne Wort ist nichts. Denn was soll ich glauben, wenn mir niemand zusagt? Soll mir etwas zugesagt werden, so muß das Wort zu mir kommen. Darum schließt Moses, daß vor allen Dingen, die wir thun, müsse zuvor GOttes Wort sein, und wir darnach müssen daran hangen durch den Glauben. So kann denn der Mensch fröhlich sprechen: Das hat GOtt gesagt, deß bin ich gewiß und kann mir nicht trügen, so weiß ich, daß ihm mein Thun gefällt; und wenn ich gleich sündige, so verschont er, und sieht durch die Finger, so lange bis ich todt bin.

17. Dies Stück muß ihn finden sonderlich in Mosis Büchern, daß mich dünkt, er habe es darum wollen thun, daß er uns lehrete, kein Werk zu thun, wir haben denn Zeugniß von GOtt, daß es ihm wohlgefalle. Damit liegt schon danieder alles Leben und Wesen, außer GOttes Gebot. So sage mir nun, wo wollen die bleiben, die eigen Wesen erdacht und erwählt haben, davon GOtt nicht geboten hat? Wie kann jemand sagen, daß des sogenannten geistlichen Standes Werke GOtt gefallen, weil ihrer keiner sagen kann, daß er es befohlen habe? Wenn du denn nicht sicher bist, daß er es befohlen hat, stehst du so bald in Gefahr, daß es ihm nicht gefalle; oder, wenn es ihm gleich gefiele, so weißt du es nicht. Derhalben bauen sie allzumal auf Sand und ungewiß. Ja, sagen sie, die christliche Kirche und heiligen Väter haben es so geordnet, darum sind wir gehorsam; so wird GOtt wiederum sagen: Ja, wo ist mein Wort? wo habe ich es befohlen, also zu ordnen?
\quellseite{96}
Solches wird dir der Teufel vorgeben, wenn du sterben sollst; so bist du denn unsicher, ob es recht geboten sei, weil sie selbst nicht wissen, ob es Gott gefällt. Aber „ungewiß“ kann den Fuß nicht erhalten, noch dem Teufel entlaufen.

18. Derhalben habe ich so oft gepredigt, daß niemand nichts ansage, das vor Gott gelten solle, er sei denn zuvor gewiß, ob es ihm gefalle. Wirst du etwas aus eigenem Vornehmen thun, daß dich es also gut dünkt, so wisse nur gewißlich, daß [es] des Teufels sei.

19. Darum merkte diesen Text eben, wie Noah, der heilige Mann voll Geistes, thut nicht einen Tritt aus der Arche ohne Gottes Wort; ja, daß kein Thier soll Junge hecken, ehe es Gott befohlen hat. Darum kommt er zuvor, ehe jemand darum redet, und spricht: „Reget euch, seid fruchtbar und mehret euch.“ Kurzum, er will keine Creatur lassen ein Werk thun, er habe es denn geheißen. So haben wir droben im ersten Capitel immerdar gehört: „Und Gott sprach“ etc., daß kein Ding also könnte gehen, wie es geht, er müßte es zuvor heißen; noch wollen wir klüger und mehr denn Gott sein, daß wir alles thun nach unserm Dünkel, und denken nicht einst auf Gottes Wort.

\begin{quote}
V. 20—22. Noah aber bauete dem HErrn einen Altar, und nahm von allerlei reinem Vieh, und von allerlei reinem Gevögel, und opferte Brandopfer auf dem Altar. Und der HErr roch den lieblichen Geruch, und sprach in seinem Herzen:\footnote{Die Worte: „in seinem Herzen“ fehlen hier in den Ausgaben; doch nachher bei der Auslegung (§ 22) sind sie gesetzt.} Ich will hinfort nicht mehr die Erde verfluchen um der Menschen willen, denn das Dichten des menschlichen Herzens ist böse von der Jugend auf. Und ich will hinfort nicht mehr schlagen alles, was da lebet, wie ich gethan habe. So lange die Erde stehet, soll nicht aufhören Same und Ernte, Frost und Hitze, Sommer und Winter, Tag und Nacht.
\end{quote}

20. Das Opfer, so geschah im alten Testament, war gemeiniglich von Tauben oder Turteltauben, und von Vögeln. Vom Vieh aber Ziegen, Schafe und Farren oder Ochsen. Die wollen wir sparen in das vierte Buch Mosis. Und unter andern Opfern ist vornehmlich das Brandopfer, davon hier gesagt wird. Das war ein solches Opfer, davon die Priester nichts kriegten, sondern mußte gar verbrannt werden, mit Mist und Fell. Welches alles auf den Glauben deutet; davon anderswo zu sagen ist.

21. Was ist aber das, daß die Schrift redet, als habe Gott eine Nase, und rieche von Schafen und Rindern, und spricht, es sei ihm ein süßer Geruch, wie man hernach oft in Mose liest? Es ist ja eine kleine Lust, zu riechen, daß man einen Ochsen oder Schaf zu Pulver brennt. Darum ist es ein geistlicher Geruch; und daß er vor Gott kommt, ist, daß es Gott erfährt, oder ihm kund wird, und nimmt sich des Opfers an, gleich als man vom Ruche, den man von einem Dinge hat, das man sonst nicht sieht, prüfen kann, was es sei. So will er sagen, daß ihm etwas Liebliches und Angenehmes geschehe, das er gerne hat und ihm wohlgefällt, davon ein gut Geschrei ausgeht und vor ihn kommt. Und ist sonderlich so geredet auf unsern Glauben. Wenn ich vor Gott etwas thue, das da angenehm ist, und mein Herz fühlt es, wie es denn gewiß muß sein, so ist mir's noch verborgen; darum, wie es vor mir im Herzen ist, so ist es vor Gott. Das Herz muß es merken, daß es Gott gefalle, angenehm und ein lieblicher Ruch sei, aber weil ich es noch nicht sehe, ist es noch zugedeckt im Glauben. Darum redet die Schrift davon, wie es zwischen Gott und uns steht, daß es also sei, wie wir es halten, wie mein Herze wandelt im Glauben. Also schreibt die Schrift davon, daß er jetzt rieche, jetzt nicht, darum, daß ich es fühle, jetzt also, jetzt anders.

22. Da kommen nun wieder die tröstlichen Zusagungen Gottes, da der Text also sagt: „Gott sprach in seinem Herzen, ich will nicht mehr die Erde verfluchen um der Menschen willen“ etc. Wer hat das Noah oder Mosi gesagt? Ist so viel: er hat es gefühlet im Grunde seines Herzens, daß Gott nicht mehr wollte die Erde verfluchen um der Menschen willen. Warum das? „Denn das Dichten des menschlichen Herzens ist böse von der Jugend auf.“ Das ist wahrlich wohl gelobt, und uns eine große Ehre. Wer sich für einen Menschen hält, mag den Reim wohl lassen stehen. Er ist uns allen in [den] Schild gemalt, daß wir wohl möchten das Angesicht niederschlagen, und nimmer fröhlich aufrichten. Und will so sagen: Was hülfe es, wenn ich sie ohne Unterlaß schlüge; ich sehe doch wohl, daß es gar verderbt ist; es ist nichts
\quellseite{97}
Gutes da von Jugend auf, daß auch kein junges Kind rein und unschuldig ist; denn er redet von denen, die nicht getauft sind. Ist nun alles Vornehmen und Dichten des Menschen, was aus dem freien Willen kommt, gar böse, was sollten denn die Werke sein? Denn das muß auch die Welt sagen: Was aus einem argen, schalkhaftigen Herzen kommt, das ist nicht gut, wenn es schon auf das allerreinste scheint, und heißt ein verrätherisch böses Stück. Da sind ja alle unsere Geistlichen, das ist, Weltheiligen, angegriffen, daß Gott dürre heraus sagt: es sind Rüben in der Haut, wo sie am besten sind; wie sie es vornehmen und dichten, ist es alles im Grunde böse.

23. Darum spricht er: Ich will nicht mehr die Welt also schlagen, sammt allem, das da lebt. Denn nach dem Mal schlägt er wohl zuweilen eine Stadt, Land, und einen Haufen, aber solches thut er nicht mehr; deß sind wir sicher bis auf den jüngsten Tag. Bleiben soll Tag und Nacht, Sommer und Winterzeit, daß nicht alles zumal verderbt werde, so lange die Erde steht. Und wird damit angezeigt, daß sie nicht ewig stehen soll, daß einmal der jüngste Tag werde kommen. So lange will er es lassen bleiben, bis er es denn alles umkehre. Indes muß immer Land bleiben, das Samen zu sich nehme, und wachse, und ernten lasse. Das soll auf Einen Tag alles zergehen, aber nicht mehr durch Wasser, sondern durch Feuer verzehrt, daß nicht mehr Tag noch Nacht, Winter noch Sommer, Same noch Ernte wird sein, sondern neue\footnote{Erlanger: neus.} Himmel und Erde, und ein ewiger Tag.

\section{Das neunte Capitel}

\begin{quote}
B. 1—3. Und Gott segnete Noah und seine Söhne, und sprach: Seid fruchtbar und mehret euch, und erfüllet die Erde. Eure Furcht und Schrecken sei über alle Thiere auf Erden, über alle Vögel unter dem Himmel, und über alles, was auf dem Erdboden kreucht; und alle Fische im Meer seien in eure Hände gegeben. Alles, was sich regt und lebt, das sei eure Speise; wie das grüne Kraut habe ich es euch alles gegeben.
\end{quote}

1. Dieweil nun Gott, wie wir gehört haben, durch seinen Zorn die ganze Welt vertilgt hatte, daß sich's ansehen ließ, als wollte er keine Gnade mehr erzeigen, darum wiederholt er hier alles, was er zuvor den Menschen gegeben hat. Das erste, daß sie sich züchtigen und mehren sollten. Das andere, daß sie Herren sollten sein über alle Thiere, Vögel und Fische. Und thut hier dazu einen Zusatz. Denn vor der Sündfluth ließ man nicht, daß er ihnen Fleisch habe zu essen gegeben, daß es hier etwas besser wird denn vor.

2. Ist aber alles darum geredet und geschrieben, daß wir sehen, wie genug gesagt ist, wie Gott nichts will gethan haben auf Erden, das nicht zuvor von ihm befohlen sei, daß ja kein Werk gehe, außer und ohne sein Wort oder Befehl. Welches auch alle Creaturen halten, ohne der Mensch und der Teufel. Die andern können nicht anders gehen, denn wie es Gott geordnet hat, ohne daß wir ärger sind und haben übertreten; doch muß es dennoch in dem nach Gottes Ordnung gehen, daß sich die Menschen züchtigen und mehren.

3. Daß er nun sagt: „Eure Furcht und Schrecken sei über alle Thiere und Vögel, und Fische im Meere“, ist nicht zu deuten, daß wir alle Fische im großen Meere sollten haben, und alle Thiere auf Erden. Ich einige Person kann nicht alle Fische und Thiere in Wäldern unter Hand haben. Es ist zu der ganzen Gemeine gesagt, daß alles, was lebt, müßte den Menschen unterliegen, wenn sie zusammen hätten. So ist die Gewalt nicht jeglichem insonderheit, sondern dem Haufen gegeben, daß er alles zu und unter sich bringen kann. „Die Fische im Meere.“ Meer\footnote{„Meer“ fehlt in der Wittenberger und in der Erlanger.} heißt er alle große Haufen Wasser, die wir Teiche, kleine und große Seen, item Sümpfe nennen, darinne Fische gehen und leben, wie der Schrift Sprachweise ist. Daß er ihnen aber Macht gibt Fleisch zu essen, ist, wie gesagt, eine neue Gabe. Aber da sind ausgezogen die unreinen Thiere, wie noch vor Augen, daß man
\quellseite{98}
dieselbigen nicht isset. Doch nimmt er nur einen gemeinen Auszug von allerlei Thieren, und spricht:

\begin{quote}
V. 4. 5. Allein, esset das Fleisch mit dem Blut nicht, darum die Seele ist. Denn ich will euer Blut, darum eure Seele ist, von aller Thiere Händen fordern, und will die Seele des Menschen fordern von des Menschen Händen, eines jeglichen Seele von des andern Händen.
\end{quote}

4. Dies Gebot wird fast getrieben im Alten Testament durch das Buch Levitici. [3 Mos. 17, 12.] Die Juden halten es auch noch strenge, daß sie das Blut alles vom Fleische ausziehen und ausdrücken. Die Ursache aber, warum es verboten ist zu essen, weiß ich nicht, sonderlich nach der Historie, denn daß er uns also hat entwöhnen wollen, daß wir nicht blutsüchtig würden, und auch ein Grauen\footnote{In den Ausgaben: „einen grauen“; doch „Grauen“ ist bei Luther immer sächlichen Geschlechts.} gewinnen, Menschenblut zu vergießen. Denn der Todtschlag war schon angegangen, und in die Welt gekommen von Kain, daß Gott ein sonderliches Verbot darauf schlug, wie er hier auch thut. So wollte er durch dies Gebot die Leute davon reißen, daß sie scheu würden. Wir aber halten solches jetzt nicht mehr, trösten uns, daß wir im neuen Testamente frei sind von solchen äußerlichen Satzungen, wie Paulus in seinen Episteln oft anzeigt, daß alle Aufsätze nicht mehr gelten in der Christenheit.

5. Also setzt er hier die Ursache selbst dazu, und spricht: „Darum die Seele ist.“ „Denn ich will euer Blut, darum die Seele ist“ etc. „Die Seele“, das ist, „das Leben“. Ich wollte gerne, daß wir ein solches Wort in unserer Sprache hätten, damit wir das hebräische Wort Nephesch könnten recht geben. Wenn wir nach unserer Sprache „die Seele“ nennen, so nennen wir den andern Theil des Menschen, der da abscheidet vom Leibe, das heißt die Schrift gemeiniglich „den Geist“. Denn die hebräische Sprache ist viel reicher, denn wir in unserer Sprache. Aber „Seele“ heißen sie das leibliche Leben, das wir führen natürlich, nach den fünf Sinnen. Also wird es gemeiniglich gebraucht in der Schrift, als man oft liest im dritten Buch Mosis [Cap. 5, 2. 7, 21.]: „Wenn sich eine Seele verunreiniget.“ Item [5 Mos. 19, 11. 22, 26.]: „Wenn jemand eine Seele schlägt“, und dergleichen, daß es nicht mehr heißt, denn ein natürlicher Mensch.

6. Also werden wir hören im zwölften Capitel [V. 5.], wie Abraham mit aller Habe, die sie gewonnen hatten, und Seelen, die sie gezeuget hatten, auszog; und Cap. 46, 27., wie Jakob mit siebenzig Seelen in Egypten zog. Wir können es nicht wohl leiden in unserer Sprache, müssen aber der Schrift Weise gewohnen. Also verstehe auch den Spruch im Evangelio [Joh. 12, 25.]: „Wer seine Seele hasset, der bewahret sie ins ewige Leben.“ Item [Joh. 10, 12.]: „Ich lasse meine Seele für meine Schafe.“ Wenn man es recht will deutschen, ist es so viel: Ich lasse mein Leibesleben für meine Schafe. Item: Wer sein Leibesleben hinsetzt und verachtet. Darum denke nicht, wie etliche Heiden gemeint haben, daß der Geist oder Seele im Blut wohne; denn wo das wäre, könnte nimmer die Seele vom Blut abscheiden. „Denn ich will (spricht er) euer Blut von aller Thiere Händen fordern, und will die Seele des Menschen fordern von des Menschen Händen“; das ist so viel gesagt: Wer Blut vergießt, wer es auch sei, davon will ich es fordern, auch wenn es ein Thier ist. Das ist ein seltsamer Spruch, dabei man sieht, daß dies Gebot eigentlich dahin sieht, daß es dem Mord und Todtschlag wehre. Darum folgt:

\begin{quote}
V. 6. 7. Wer Menschen Blut vergießt, des Blut soll auch durch Menschen vergossen werden, denn Gott hat den Menschen zu seinem Bilde gemacht. Seid fruchtbar und mehret euch, und reget euch auf Erden, daß euer viel darauf werden.
\end{quote}

7. Das2) ist das erste Gebot von dem weltlichen Schwert: Wer Menschen Blut vergießt, der soll schuldig sein, daß sein Blut wieder vergossen werde; doch nicht, daß er [Gott] es selbst wolle thun, sondern [es] soll durch Menschen geschehen. In den Worten ist eingesetzt die weltliche Obrigkeit, und das Recht von Gott, das Schwert, ihr in die Hand gegeben. Wie auch droben [1 Mos. 4, 14.] von Kain, da er sprach: „So wird mir es gehen, daß mich todtschlägt, wer mich findet“, welches er nicht gesagt hätte, wenn Adam nicht hätte das Schwert geführt. Darum wird Noah das auch gepredigt und geführt haben. Das ist eins, das hier gegründet ist.\footnote{Erlanger: Da.}
\quellseite{99}

8. Aber den Spruch verstehe also, daß wahr ist, daß aller Todtschläger Blut durch Menschen Hände vergossen wird; es sei denn Sache, daß Gott zuvor kommt. Denn er setzt hier nur Recht und Macht, ob es wohl oft durch die, so das Recht haben, wird nachgelassen. Damit wird aber nicht das Recht aufgehoben, sondern geht gleichwohl im Schwange; denn er redet nicht de facto, was geschieht, sondern de jure, was geschehen soll, was man thun müsse; aber darum geschieht es nicht so bald allezeit.

\begin{quote}
V. 8—17.
\end{quote}

\begin{quote}
Weiter sagte Gott zu Noah und seinen Söhnen mit ihm: Siehe, ich richte mit euch einen Bund auf, und mit eurem Samen nach euch, und allem lebendigen Thier bei euch, an Vögeln, an Vieh und an allen Thieren auf Erden bei euch, von allem, das aus dem Kasten gegangen ist, waserlei Thiere es sind auf Erden. Und richte meinen Bund also mit euch auf, daß hinfort nicht mehr alles Fleisch verderbt soll werden mit dem Wasser der Sündfluth, und soll hinfort keine Sündfluth mehr kommen, die die Erde verderbe. Und Gott sprach: Das ist das Zeichen meines Bundes, den ich gemacht habe zwischen mir und allem lebendigen Thier bei euch hinfort ewiglich. Meinen Bogen habe ich gesetzt in die Wolken, der soll das Zeichen sein meines Bundes zwischen mir und der Erde. Und wenn es kommt, daß ich Wolken über die Erde führe, so soll man meinen Bogen sehen in den Wolken. Alsdann will ich gedenken an meinen Bund zwischen mir und euch, und allem lebendigen Thier in allerlei Fleisch, daß nicht mehr hinfort eine Sündfluth komme, die alles Fleisch verderbe. Darum soll mein Bogen in den Wolken sein, daß ich ihn ansehe, und gedenke an den ewigen Bund zwischen Gott und allem lebendigen Thier, in allem Fleisch, das auf Erden ist. Dasselbe sagte Gott auch zu Noah: Dies sei das Zeichen des Bundes, den ich aufgerichtet habe zwischen mir und allem Fleisch auf Erden.
\end{quote}

9. Siehe die Ordnung an. Zum ersten hat er es wieder nach dem äußerlichen Wesen angesetzt, mehr Geschenk gethan, Freiheit Fleisch zu essen, und neue Gesetze gegeben. Denn er sieht wohl, daß nichts Gutes an uns ist; darum, kann er nicht inwendig, so will er doch damit äußerlich mehren mit dem Schwert den Bösen; denn die Frommen dürfen’s nicht. Nun macht er auch den Glauben für die Seinen, gibt erstlich das Wort, darnach das Zeichen. Denn er hat von Anfang so gehandelt: wenn er das Evangelium gibt, läßt er es nicht bei dem Wort bleiben, sondern thut ein Zeichen dazu. Als wir im Neuen Testament haben das Wort [Marc. 16, 16.]: „Wer da glaubet“ etc., dazu auch das Zeichen: „Wer getauft wird.“ Desgleichen haben wir Christi Leib und Blut im Brod und Wein neben dem Wort [Matth. 26, 26. 27.]. Also thut er wie ein frommer, treuer Mann, der einen Brief schreibt, und sein Siegel darauf drückt.

10. Aber diese Worte sind noch nicht das rechte Evangelium, denn es ist eine Zusagung, die betrifft nicht allein sie, die Menschen, sondern alle Thiere auf Erden; ist ein leibliches Gut und Wohltat, die er verheißt, ehe Noah darum bat oder denkt. Darum hält er diese Worte, wir glauben oder nicht. Noah hatte den Glauben, und war ein frommer Mensch. Die Thiere konnten nicht glauben, dennoch hat er es denen auch zugesagt. Doch hilft es Noah, daß dadurch sein Glaube gestärkt und getröstet wird. Er redet aber mit Unterschied, nicht also, daß sein Fleisch sollte umkommen, sondern daß [es] mit Gewässer und Sündfluth nicht mehr solle verderbt werden. Als wollte er sagen: Ich will hernach ein ander Verderben anrichten, nicht mit Wasser, sondern mit Feuer am jüngsten Tage [2 Petr. 3, 7.]. Item, daß er spricht: „alles Fleisch“, zeigt an, daß wohl geschehen mag, daß eine Stadt oder Land ersäufet werde; item, daß die Thiere im Wasser ersaufen; aber ganz insgemein soll es nicht untergehen. Das sind die Worte; nun vom Zeichen.

11. Hier haben der Heiden Meister, die Philosophi, klug wollen sein, und sich fast darüber bekümmert und gebrochen, was der Regenbogen sei, und feindlich getrachtet,\footnote{Das heißt, gar sehr darnach getrachtet. Jenaer und Wittenberger: „fleißig betrachtet“.} wie es zugehe, daß er allemal gegen der Sonne, und gemeiniglich gegen dem Morgen [stehe], wenn die Sonne gegen dem Abend steht. Item, daß er nur halb rund sei, haben etliche also gesagt, es\footnote{Erlanger: er.} komme davon, daß die hohlen Wolken den Strahl von der Sonne zurück stoßen gegen die Sonne, so breche sich derselbe, und mache einen solchen Bogen von mancherlei Farben. Etliche bringen auch andere Ursachen, meinen, wir sollen es
\quellseite{100}
glauben. Aber wenn man es eben ansieht, sind es ja lose Fabeln, als hätten es eitel alte Weiberln gesagt. So laß du solche Träume fahren und folge der Schrift, die sagt, es komme davon, daß hier Gott spricht: „Wenn ich den Himmel mit Wolken überziehe, so soll man meinen Bogen sehen in den Wolken“; nicht daß man es allemal sehe, so oft Wolken am Himmel sind, denn das sagt er nicht; sondern wenn es ihn lüstet, so macht er, daß wir ihn sehen.

12. So oft wir ihn nun sehen (sagt er), so will er denken an den Bund, den er gemacht hat mit uns. So sollen wir wissen und verstehen, wie er so viel bedeutet, daß Gott seiner Zusagung gedenkt, die Welt nicht mehr durch Sündfluth zu ersäufen, und uns erinnern des großen Zorns und Urtheils, das zuvor über die ganze Welt gegangen ist, daß kein schrecklicheres ist von Anfang bis an das Ende der Welt. Also ermahnt er uns des schrecklichen Zorns, den die Welt verdient hatte, und auch noch wohl verdient, und tröstet uns doch daneben, daß er sie nicht mehr so strafen wolle, weil die Erde steht.

13. Was ist aber das, daß er sagt: „Alsdann will ich gedenken an meinen Bund“; item: „Darum soll mein Bogen in den Wolken sein, daß ich ihn ansehe, und gedenke meines ewigen Bundes“? Wann ist er einmal so schläfrig, daß er es vergesse, so doch alle Dinge immerdar vor seinen Augen sind, daß er es sehen muß, und nicht vergessen kann? Es sind abermal Sprüche, die des Glaubens Art anzeigen. Das heißt die Schrift, Gott gedenke, wenn wir fühlen und inne werden, daß er es gedenke, nämlich, wenn er sich äußerlich so stellt, daß wir dabei merken, wie er daran gedenkt. Denn weil der Bogen nicht immerdar steht, scheint es äußerlich, als gedenke er nicht allezeit daran; wenn man ihn aber sieht, so können wir sagen: Siehe, da gedenkt Gott seines Wortes. Also geht es alles dahin: wie ich mich gegen Gott stelle, so stellt er sich gegen mir; er kann sonst nicht mit uns handeln, denn mit Worten und Zeichen. Derhalben heißt das „gedenken“, wenn es so geht, daß ich es sehe und höre. Deß mußt du gewohnen durch die ganze Schrift. Also haben wir die Sündfluth hinweg, welches alles um unserwillen geschrieben ist, daß wir lernen, was Gottes Werke sind.

14. Ist aber noch der Punkt da, was da bedeutet, daß man kein Blut essen soll; das müssen wir ein wenig handeln, weil es hier erstlich steht, und oft in der Schrift wiederholt wird. Ich habe aber vor gesagt, wo man nicht die Deutung erweisen kann, soll man nicht darauf fest bauen und stehen; also sage ich hier, das, was gedeutet wird, ist recht und wahr, aber daß sich's hieher reime, da will ich nicht auf sterben. Zum ersten: Essen und Trinken heißt in der Schrift glauben, oder, die Lehre mit dem Glauben fahren; wie im Evangelio Johannis das Brod bedeutet die Lehre oder Gesetz, die Fische gute Exempel. Das ist recht und gewiß durch die Schrift durch und durch; denn Christus Joh. 6. 54. klärlich sagt: „Wer mein Fleisch isset, und mein Blut trinket, der hat das ewige Leben“, ist nicht von leiblichem, sondern von geistlichem Essen gesagt, wie er es, V. 35., selbst deutet, das ist: „Wer an mich glaubet.“ Also heißt essen immerdar glauben; aber darnach die Speise ist, darnach ist auch das Essen, wie die Lehre ist, so ist auch der Glaube. Ist die Lehre recht, so wird auch ein rechter Glaube; ist sie falsch und vergiftet, so ist es auch ein falscher, todter Glaube.

15. Dieweil denn auch hier von Essen gesagt wird, muß [es] freilich die Lehre und den Glauben betreffen, also, daß das Blut bedeute die Lehre, die Gott nicht haben will. Welche ist aber die? Ich habe oft gesagt, das christliche Leben bestehe in zweien Stücken. Wenn man predigt vom Glauben, soll man zusehen, daß man auf das lauteste predige, denn er kann nicht leiden, daß man etwas daneben einführe; die Liebe aber kann viel Dinge leiden, wie Paulus [1 Cor. 13, 7.] sagt: „Die Liebe duldet alles“, auch die bößesten Buben, wie Christus uns getragen hat; sondern der Glaube trägt nichts, spricht also: Daß die Leute unvollkommen leben, da mag man Geduld haben, aber mit unrechter Lehre habe ich keine Geduld.

16. Derhalben ist großer Unterschied unter der\footnote{„der“ fehlt in der Erlanger.} Lehre und Leben. Mit der Lehre gilt es nicht Scherzens, die muß rein und recht bleiben, aber mit dem Leben halten wir es nicht so strenge. Wie man auch im Evangelio sieht, daß Christus Geduld hat mit den Jüngern und durch die Finger sieht, wenn sie gleich gröblich straucheln. Doch straft er sie zuweilen auch, und spricht: Das ist nicht recht; das ist die Lehre. Mit den Phariseern aber und Heuchlern
\quellseite{101}
hatte er nimmer seine Geduld; denn es trifft nicht das Leben, sondern die Lehre an. Es liegt die größte Macht an der Lehre; wenn die rein bleibt, so kann man allerlei unvollkommenes Leben und Schwachheit tragen, so ferne, daß man an der Lehre halte, und bekenne, daß das Leben anders sein sollte; wo aber die Lehre verfälscht wird, so ist dem Leben auch nicht mehr zu helfen.

17. Also läßt hier Gott nach, und will, daß man Fleisch esse, verderbt aber das Blut zu essen. Als sollte er sagen: Halte nur den Glauben rein, ob du gleich die Liebe nicht kannst rein und vollkommen halten und haben. Das Blut ist rein im Menschen, denn die Seele ist darinne. Das ist, der Glaube lehrt also: Du bist verdammt von der Scheitel an bis auf den Fuß. Darum taucht und tauft man uns in das Wasser, damit bedeutet wird, daß der Mensch gar ersäufen und untergehen muß in diesem Leben, darum, daß dies leibliche Leben eitel Sünde ist. Wer nun anders predigt, und wirft unser eigen Werk auf, wie die Juden auch thaten, als sie sagten, man müßte das Gesetz Mosis halten, der gibt Blut zu essen, darinne das leibliche Leben ist, das ist, er läßt die Leute stecken in ihrem Wuste und Blute, da wir in geboren sind.

18. Also könnte ich alles wohl tragen, was Pabst, Bischöfe \&c. für Fehler haben außer der Lehre; aber daß sie ihre Lehre wollen gepredigt haben, und uns verdammen, das ist uns nicht zu leiden. Denn sie haben die Welt verschlemmt, haben nichts gepredigt, damit der alte Adam getödtet würde, und wer es hat predigen wollen, [den] haben sie nicht wollen leiden. Darum ist ihre Lehre nichts anders, denn Blut fressen und sauffen. So siehst du, warum Gott so hart verboten hat, Blut zu essen, das ist, nichts zu predigen, das Menschen vermögen, und die Werke stärcke, sondern nur das, daß der Mensch gar vor Gott sterben müsse. Darum, wenn ich so predige: Lieber Mensch, willst du selig werden, so ist nöth, daß du verzweifelst an allem, was Menschen ist, und sterbest dem allen ab, was dir angeboren ist [1 Petr. 2, 24.], den alten Adam gar ausziehest, und in eine andere Haut kriechest [Col. 3, 9. 10.], das ist rein gepredigt, und kein Blut gegessen.

19. Danider ist nun jene Predigt. Denn sie lassen den Menschen noch leben in seinem Willen und Werken, als die vor Gott etwas sollten gelten. Wenn du nun lehrst, den Menschen also tödten [Röm. 8, 13. Col. 3, 5.], so magst du wohl Fleisch essen; aber vor dem Blut mußt du dich hüten, das ist, daß du lehrest, das Fleisch, das noch nicht in die Grube ist, dämpfen. Das Fleisch kann man nicht gar hinweg werfen; denn was nicht vollkommen ist, müssen wir leiden, und die Schwachheit tragen. Das sei davon genug.

\begin{quote}
V. 18—29.
\end{quote}

\begin{quote}
Die Söhne Noah, die aus dem Kasten gingen, sind diese: Sem, Ham, Japhet. Ham aber ist der Vater Canaan. Das sind die drei Söhne Noah, von denen wurden alle Länder besezt. Noah aber fing an und ward ein Ackermann, und pflanzte Weinberge. Und da er des Weins trank, ward er trunken, und lag in der Hütte aufgedeckt. Da nun Ham, Canaans Vater, sahe seines Vaters Scham,\footnote{Randglosse: Vaters Scham. Diese Geschichte deuten viele dahin, man solle der Prälaten Laster nicht strafen, welches doch Christus und alle Apostel thaten. Aber deute du es recht, daß Noah sei Christus und alle Gläubigen. Die Trunkenheit sei die Liebe und Glaube im Heiligen Geist. Die Blöße sei das Kreuz und Leiden vor der Welt. Ham seien die falschen Werkheiligen und Gleißner, die Christum und die Seinen verspotten, und Lust haben in ihren Leiden. Sem und Japhet seien die frommen Christen, die solches Leiden preisen und ehren.} sagte er es seinen beiden Brüdern draußen. Da nahmen Sem und Japhet ihr Kleid, und legten es auf ihre beiden Schultern, und gingen rücklings hinzu, und deckten ihres Vaters Scham zu, und ihr Angesicht war abgewandt, daß sie ihres Vaters Scham nicht sahen. Als nun Noah erwachte von seinem Wein, und erfuhr, was ihm sein kleiner Sohn gethan hatte, sprach er: Verflucht sei Canaan, und sei ein Knecht aller Knechte unter seinen Brüdern. Und sprach weiter: Gelobet sei Gott, der Herr des Sem, und Canaan sei sein Knecht. Gott breite Japhet aus, und lasse ihn wohnen in den Hütten des Sem, und Canaan sei sein Knecht. Noah aber lebete nach der Sündfluth dreihundert und fünfzig Jahre, daß sein ganzes Alter ward neunhundert und fünfzig Jahre, und starb.
\end{quote}

20. Diese Historie ist an ihr selbst leicht, daß Noah, der heilige Patriarch, die Welt auch lehrt zeitliche Nahrung suchen, als, daß er Weinberge pflanzte und baute. Denn das ist nicht zu verwerfen, daß die Schrift redet von den heiligen Vätern, wie sie äußerlich geringe Werke anzusehen
\quellseite{102}
gethan haben, also, daß es nicht übel, sondern wohl gethan ist, daß man arbeite die Erde, und die Nahrung heraus bringe, was Gott geben will.

21. Wiewohl daneben auch angezeigt ist, daß man desjenigen mäßig brauchen soll. Denn Gott überschüttet uns gar reichlich mit seinen Gütern; aber derselbigen braucht niemand recht denn der Geist; wo der nicht ist, wird es mißbraucht, daß man es ruchlos\footnote{In den alten Ausgaben: rauchlos.} hinweg verschlemmt, oder daß der Geiz zu sehr zu sich reißt, und niemand läßt genießen.

22. Auch gibt dieses Capitel eine fast gute Predigt für uns Deutschen. Denn hier haben wir ein Exempel, daß wir uns trunken und voll saufen. Denn der Text sagt klar, wie der heilige Mann sei trunken worden, wie wir auch hernach [Cap. 19, 33. 35.] vom Lot hören werden. Nun können wir dennoch so bald dies Werk nicht tadeln, daß es so gar übel gethan sei; wiewohl die Schrift nicht meldet, ob es wohl oder übel sei gethan. Aber das sagt sie, daß er ein heiliger Mann gewesen sei, gibt auch wohl zu verstehen, daß er es freilich aus keiner bösen Meinung oder Vorsatz gethan habe, daß er es vielleicht versehen hat, oder nicht versehen, sondern vielleicht hat, als er müde gewesen ist, sich erquicken wollen, und so trunken ist worden; also, daß solch einerley Werk wohl von einem heiligen, frommen Manne geschehen möchte, wie sonst von bösen Buben. Wie denn in vielen andern Werken gemeiniglich geschieht, als, zürnen, sauer sehen, schelten ist gemein in der Welt, noch leiden wir, daß es Christus selbst gethan hat.

23. Das rede ich darum, daß man die Heiligen ansehen soll, wenn man ihr Leben und Exempel führt, nicht nach den Werken, sondern nach der Person und ihrem Glauben. Es ist kein fährlicher Wesen, denn wo man fällt auf Werke, so die Heiligen gethan haben, und sieht nicht an den Glauben, in und aus welchem sie die Werke gethan haben. Als, wenn ich auf das Werk fiele, und wollte mich auch voll saufen, und darnach sagen: Hat es doch der heilige Mann Noah gethan, so führe ich zum Teufel, wie er gen Himmel. Desgleichen richte auch von andern Werken der Heiligen. Darum, wer den Heiligen nachfolgen will, der muß den Glauben fassen, darinne sie gegangen sind. Wenn er den hat, so kann er ohne Fahr ihnen nachfolgen.

24. Also lesen wir in dem Buch der Richter [Cap. 15, 8. 11.] von dem heiligen Richter Simson, wie er die Philister schlug, und sein Volk ihn fragte: warum er das gethan hätte? da sprach er: „Wie sie mir gethan haben, so habe ich ihnen wieder gethan.“ Siehe, das lautet je unchristlich, und wider das Evangelium. Darum, wenn du dem Exempel folgen wolltest, und auch so thun, so führest du zum Teufel. Ursache ist, daß du nicht der Mann bist, der Simson war, und er nicht war, der du bist. Er war im Herzen rechtschaffen, darum war das Werk auch unsträflich. Denn er war so gesinnt, seinethalben zu leiden, daß sie ihn getödtet hätten. Weil er aber Richter geordnet war, und sollte die Philister strafen, und sein Volk aus ihren Händen retten, redet er recht; als sollte er sagen: Was sie meinem Volke gethan haben, das haben sie mir gethan.

25. Darum beschließen wir mit der Epistel zum Hebräern [Cap. 11], da er erklärt, wie die heiligen Erzväter alle Werke aus dem Glauben gethan: „Sehet an eure Vorgeher, wie sie geglaubt haben“; als wollte er sagen, nicht, welche Werke sie gethan haben. Denn es nicht möglich ist, daß wir alle Werke thun, wie die Heiligen. Es würde auch nichts denn eitel Secten und Uneinigkeit daraus werden. Mit der Weile müßten alle Schuster Fischer sein, weil die Apostel sind Fischer gewesen [Matth. 4, 18.]. Sollte nicht ein Schuster fromm können sein, und sein Handwerk treiben im Glauben? Müßte aber ihm darum jedermann nachfolgen? Darum laß einen jeglichen seines Werkes warten. Wenn es aber dazu käme, daß Noth wäre, so folgst du auch seinem Werke nach. Sonst bleibe du bei deinem, so bleibst du auf der Bahn. Mancherlei sind die Werke, einerlei ist der Glaube. So magst du dem äußerlichen Wesen folgen, das gut heißt, daß du wohl thust, wie er gethan hat, aber nicht eben das Werk, die Person, Zeit, Ort etc. [be]stimmest.

26. So sage ich nun, daß Noah trunken ist worden, das kann ohne Sünde geschehen sein. Wenn es ein anderer thäte, so wäre es Sünde. Das macht, daß man die Werke und den Glauben scheiden muß. Gott richtet nach der Person; wie dieselbige ist, so ist auch das Werk.

27. Es haben auch viele gepredigt wider die
\quellseite{103}
Trunkenheit, und zwar wir Deutschen haben sonderlich daß\footnote{„daß“ fehlt in der Erlanger.} Geschrei davon in andern Ländern. Es ist auch nicht eine köstliche Tugend. Aber wie schändlich es ist, so ist das auch wahr, wenn ich eines wählen müßte, wollte ich noch lieber dies Laster leiden, denn daß sie treiben und führen, die uns gar stolz darum verachten, daß ich nicht nennen will. Ich will auch nicht dazu reizen, sondern vielmehr davon, und wollte es auch gerne gar mehren und abbringen, wenn ich könnte. Derhalben darf niemand das Exempel für sich nehmen, und seine Untugend schützen.

28. Nun weiter. Siehe, wie ein unbeständiges Ding ist es um einen Menschen, und wie schrecklich Gottes Gericht ist. Da sind nur acht Personen, die in der Arche erhalten waren auf Erden, vier Männer und vier Weiber, und alle heilig. Denn, wäre Ham nicht fromm gewesen, hätte ihn Gott nicht mit den andern im Kasten erhalten, auch hätte er es nicht können tragen. Denn es müßte ein großer, starker Glaube sein, der vor der greulichen Plage bestehen sollte, und im Wasser schweben, und gleichsam mitten im Tode, der ihn stets vor Augen war, leben. Daß wir uns wohl mögen entsetzen, daß ein solcher trefflicher Mann in hohem Glauben so herunter fällt, daß er verflucht wird mit allen seinen Nachkommen, und ihm Gott so feind wird, daß er hernach alle sein Geschlecht vertilgt. So nun der nicht besteht, der solchen Glauben gehabt hat, wie leichtlich mögen wir fallen, die wir noch so wenig Glaubens haben?

29. Aber ich habe gesagt, es ist Gottes Gericht in allen Historien vorgeschrieben von Anfang der Welt, daß wir sehen, was Gottes Werke sind. Das ist aber Gottes Werk, daß niemand so hoch droben ist, daß er nicht herunter fallen könnte, und wiederum, niemand so niedrig ist, der nicht möge hinauf kommen. Darum ist auf beiden Seiten keine Ursache, eben so wenig zu verzweifeln, als sich zu vermessen.

30. So sehen wir in dem Exempel, daß er und sein Weib in Vermessenheit gefallen, sicher worden sind, und haben gepocht auf ihren Glauben, und also in Nachlässigkeit gefallen. Sind so lange herab gesunken, bis sie gar zu Boden gefallen,\footnote{Erlanger: fallen.} auf daß ein jeglicher, der im Glauben steht, sich vorsehe und fürchte, und nicht sicher werde, wie Paulus [1 Cor. 10, 12.] auch\footnote{„auch“ fehlt in der Erlanger.} sagt: „Wer da stehet, der sehe, daß er nicht falle.“ Man kann nicht sagen, daß dieser nicht rechte Gnade oder Glauben gehabt habe; [er] ist viel köstlicher gewesen, denn vielleicht jetzt aller auf Erden, denn man würde kaum einen finden, der ihm jetzt gleich wäre nach Kunst und Verstand. Was thut aber nun Noah? Da er erfährt die Schalkheit, die ihm sein kleinster Sohn gethan hatte, spricht der Text:

\begin{quote}
V. 25. 26. Verflucht sei Canaan, und sei ein Knecht aller Knechte unter seinen Brüdern. Und sprach weiter: Gelobet sei Gott, der Herr des Sem, und Canaan sei sein Knecht etc.
\end{quote}

31. Das ist auch anzusehen als wider das Evangelium, denn es ist je verboten mit ausgedrückten Worten, man soll sich nicht rächen, niemand fluchen, noch Böses wünschen, sondern eitel Gutes [Röm. 12, 19. 14.]. Ja, wenn es ein anderer gethan hätte ohne Geist, so taugte es nicht, weil es aber Noah thut, so ist es wohlgethan, denn der Heilige Geist thut es, der in ihm ist. Wer will nun sagen: Warum thust du das? Es kann niemand einen Christenmenschen urtheilen, und sagen, da thust du nicht recht; wie Paulus [1 Cor. 2, 15.] sagt: „Der geistliche Mensch richtet alles, und er wird von niemand gerichtet“; was er schließet, das ist recht, denn was er urtheilt, das urtheilt Gott; wer aber über ihn urtheilt, der urtheilt über Gott. So geht die heilige Schrift und ihre Exempel in großen Dingen, darein sich die Welt gar nicht zu schicken weiß.

32. Nun, der Fluch ist nichts Anderes, denn äußerlich maledeien, daß Ham soll seinen Brüdern unterworfen sein. Und das ist die Prophezei, die längst hernach gefolgt und erfüllt ist. Denn wir sehen bald hernach, daß der Ham mit seinem Geschlechte gar hoch gekommen ist, und fast den Herrn und besten Strich des Landes eingenommen hat, dazu auf das herrlichste gebauet. Noch ist hier ein klarer Spruch, er soll ein Knecht der andern sein, daß dafür zu halten ist, daß der Fluch erst hernach, da Gott das jüdische Volk in das Land gebracht, ist angegangen.

33. Das geht nun dahin, daß, wenn Gott ein Wort redet, so scheint es, als habe es ein Narr gesagt, und sei gelogen; desgleichen muß
\quellseite{104}
und kein Werk geben, als er hier sagt, Ham soll unten liegen, und Knecht sein, so sieht man, daß er dennoch gar lange Zeit immer oben liegt. Wo bleibt denn indes Gottes Wort? Aber darum heißt es Gottes Werk, daß seine Art ist: wen er will hoch heben, den drückt er nieder; wen er lassen und niederdrücken will, den hebt er zuvor empor.

34. Also spielet er mit allen Königreichen, Herrschaften, Wesen und Hohen noch immerdar. Also trieb er es mit den Königreichen Assyrien, Persien, Griechen, Alexandrien\footnote{Alexandrien wird wohl das Reich Alexander's des Großen bedeuten.} und zuletzt mit Rom. So thut er noch immerdar: Wer am höchsten schwebt, und meint, er könne nicht sinken, den wirft er zu Boden. Das ist sein Werk, da gehört Glaube zu; der Geist muß sein, der da sage: Es wohl anders scheint, denn die Worte lauten, dennoch wird Gott seinem Worte nachkommen und erfüllen, ob man es gleich nicht sollte gewahr werden. Laß nur den Ham stark und gewaltig genug werden: wenn er meint, er sitze am sichersten, so wird er im Augenblicke im Drecke liegen. Das sei die Historie und Exempel. Wir wollen nun die heimliche\footnote{Zenner: geistliche.} Deutung auch rühren.

35. Bisher hat man den guten Patriarchen zum geistlichen Prälaten gemacht, und dies Exempel so gedeutet, man solle die geistlichen Prälaten nicht antasten noch schelten, oder ihre Laster strafen. Wie der böse Ham des Vaters lacht, da er so schamlos liegt, und mit Fingern auf ihn deutet. Das haben sie so in die Welt getrieben, und sich damit geschützet, daß sich jedermann gewehret hat, wider sie das Maul aufzuthun. Nun ist es je dennoch ein Jammer und jämmerlich Ding, wenn man also Sprüche mit falschem Verstande in die Leute treibt, daß sie es sähen; daß [es] hernach mehr kostet, wie ich selbst erfahren habe, daß man es wieder verlerne, denn daß man es vor gelernt hat.

36. Darum muß man erstlich einen rechten Verstand daraus fassen. Wahr ist es, wir sollen so leben mit dem nächsten Christenmenschen: wenn wir eine Schande wissen, die noch heimlich ist, sollen wir sie zudecken; wie ich gerne wollte, daß man mir thäte [Matth. 7, 12.]. Wenn es aber offenbar ist, daß es jedermann weiß, und nicht mehr verdeckt bleiben kann, sollen wir\footnote{Erlanger: wir es.} abereins helfen schützen und beschönen, so viel man kann. Aber daß man es ungestraft lasse und nicht wehre, das gilt in keine Wege. Denn so sagt Paulus zu Timotheo [1. Ep. 5, 20.]: Du sollst öffentlich strafen, daß sich die andern daran stoßen. Und siehe, wie er selbst verfährt, und schilt auch ihr geistliches Wesen, heißt sie Heuchler und verdammte Secten etc. Warum thut er das? Darum, daß ihr Wesen ein öffentlich Aergerniß war, das das Volk dahin führte, daß er sagen sollte, daß ihr Ding recht wäre. Da galt nicht mehr zudecken noch schweigen; denn damit werden die Seelen verderbt.

37. Derhalben, wenn das Laster so einreißt und thut Schaden, so muß man nicht länger schweigen. Weil es aber heimlich ist, daß es niemand Schaden thut, denn ihm allein, so soll ich schweigen und heimlich strafen, heraus helfen oder zuscharren, daß es nicht weiter ausbreche.

38. So haben uns unsere blinden Leiter vorgebläuet, man soll den Papst noch die Geistlichen nicht schelten, wenn sie ihr Amt nicht führen, und viel weltlicher leben denn kein anderer, und das Exempel dahin gerissen, und ein eigen Capitel darauf gemacht in ihrem heiligen Rechte. Da sagen wir Nein zu. Laß geschehen, daß man schweige, was da heimlich ist, aber das jedermann kund ist, daß ihr Wesen und Regiment wider das Evangelium ist, müssen wir öffentlich sagen, daß das Volk nicht in [den] Wahn komme, als wäre es recht.

39. Das ist aber noch die geringste Strafe, wiewohl sie sich allein davor fürchten. Aber das ist viel ärger, wenn man sie strafen muß um ihr heiliges Leben, das jedermann für köstlich hält. Da sollte Stein und Holz schreien und sagen: Ich will schweigen ihres schändlichen Lebens, so sie führen, aber das ist mit ihnen zu leiden, daß sie mit ihrem schönen Wesen und Werken die Leute wollen fromm machen. Da muß ich stehen und predigen: Ihr seid Seelenmörder, und bringt damit die ganze Welt in Jammer und Noth; ihr seid des Teufels am meisten, wo ihr am allerheiligsten seid.

40. Siehe nun an dies Exempel mit seiner rechten Deutung. Gott hat mehr zu schaffen, denn daß er durch den heiligen Mann sollte böse
\quellseite{105}
Übel deuten. Ist er heilig, so wird er auch einen Heiligen deuten. Die heilige Schrift geht nicht mit ihrer Deutung auf äußern und Buben, sondern gemeiniglich auf Christum. Darum, willst du es ohne Fahr deuten, so deute es auf den Christum. Denn, wie Noah den Weinberg pflanzt, und des Weins trinkt, davon trunken wird und entschläft, und bloß in der Hütte liegt, und wird von diesem verlacht, aber von andern zugedeckt: also ist es auch Christo ergangen. Er bauete einen Weinberg [Jes. 5, 2. 7.], das war das jüdische Volk; davon trank er, das ist, er mußte von Juden leiden und am Kreuze hangen; da machte ihn die große Liebe trunken, daß er starb. Und weil er so leidet vor großer Liebe, davon er trunken ist, daß er sein selbst vergißt, daß er uns helfe, so hängt er da mit allen Schanden, wird gehöhnt, geschmäht und geschändet, wie ein unschuldiger Mensch, ja, wie der ärgste Dieb und Mörder [Matth. 27, 39. ff.].

41. Ueber dem Leiden und heiligen Kreuze scheiden sich nun die zwei, Fromme und Böse. Es ärgert den Ham, aber die andern schämen sich sein. Ham spottet sein, und weiset mit Fingern auf ihn. Das sind die Heiligen, die Juden, die nicht durch Leiden und Schande, sondern durch ihre Werke gen Himmel wollen. Darum ist es nicht also, daß Ham solle deuten die, so die Prälaten der römischen Kirche verachten, sondern daß die Prälaten, Bischöfe etc. dieser Ham selbst sind, die nichts wollen leiden, keine Schande tragen, sondern die Christi und der frommen Christen Leiden verachten, spotten, und lästern. Sie sollten es beschönen, ehren und schmücken, so verlachen sie es und verdammen’s.

42. Die andern aber sind die frommen Christen, die das Evangelium wohl verstehen, das Wort des Kreuzes annehmen durch den Glauben, welches da Scham und Schande mit sich bringt, wie Christus selbst hat müssen tragen. Darum verspotten sie es nicht, sondern halten es in allen Ehren, bekennen es vor jedermann, wenden daran alles, was sie haben, nehmen die Kleider auf die Schulter, und decken’s zu, das ist, bekennen mit Worten und Werken, daß dies das rechte Leben sei, daß Christus muß leiden und sterben, und alle, die in Christo sind.

43. Also ist die Blöße Noah nichts Anderes, denn das schändliche Leiden Christi und der Seinen, das da gepredigt und angenommen wird von den Gläubigen. Ham aber sind die großen Heiligen, nicht grobe Sünder, die das Leiden verhöhnen, wie sie bei dem Kreuze sprachen: Ja, ist das Gottes Sohn? „Ist er Christus, so steige er herab und helfe ihm selbst“ [Matth. 27, 40.]. So haben sie auch allen Märtyrern gethan; die um des gekreuzigten Christi Namens willen\footnote{„wollen“ fehlt in der Erlanger.} gestorben sind. So wird man uns noch thun, so wir es würdig werden. Das ist der fromme Sohn, der trösten sollte und helfen tragen, so ist er fröhlich und stolz darüber, lacht des Unglücks. Wie Christus im 35. Psalm B. 13—15. sagt: Da es ihnen übel ging, da trug ich Leide, und zog einen Sack an; aber nun es mir übel geht, da tanzen und springen sie, und freuen sich meines Schadens. Und ist noch das Aller schwerste, daß dies des heiligen Vaters eigener Sohn thut; das Christo auch widerfahren ist von seinem eigenen Volke, wie der Psalm [55, 14. 15.] klagt. Darum verflucht er sie auch, wie Noah den Ham verflucht, an vielen Orten, als, durch den ganzen hundert und neunten Psalm. So hast du die Deutung dieses Exempels, daß es, wie fast alle andern Figuren, deutet das Wesen, so das Evangelium anrichtet und mit sich bringt, wenn es\footnote{„es“ fehlt in der Erlanger.} uns auch so widerfährt, daß wir wissen, daß es so sein muß. Darum hüte dich vor den schändlichen, verkehrten Lehren und Deutungen, die unsere Papisten bisher in das Volk getrieben haben.
\quellseite{106}

\section{Das zehnte Capitel}

B. 1—12. Dies ist die Geburt der Kinder Noah, Sem, Ham, Japhet; und sie zeugeten Kinder nach der Sündfluth. Die Kinder Japhet sind diese: Gomer, Magog, Madai, Javan, Thubal, Mesech und Thiras. Aber die Kinder von Gomer sind diese: Askenas, Riphath und Thogarma. Die Kinder von Javan sind diese: Elisa, Tharsis, Kittim und Dodanim. Von diesen sind ausgebreitet die Inseln der Heiden in ihren Ländern, jegliche nach ihrer Sprache, Geschlechte und Leuten. Die Kinder von Ham sind diese: Chus, Mizraim, Put und Canaan. Aber die Kinder von Chus sind diese: Seba, Hevila, Sabtha, Raema und Sabtecha. Aber die Kinder von Raema sind diese: Scheba und Dedan. Chus aber zeugete den Nimrod, der fing an ein gewaltiger Herr zu sein auf Erden, und war ein gewaltiger Jäger vor dem Herrn. Daher spricht man: das ist ein gewaltiger Jäger vor dem Herrn, wie Nimrod. Und der Anfang seines Reiches war Babel, Erech, Acad und Chalne im Lande Sinear. Von dem Lande ist darnach kommen der Assur, und bauete Ninive und der Stadt Gassen, und Chalah, dazu Resen zwischen Ninive und Chalah, dies ist eine große Stadt.

1. Hier hebt nun Moses an zu beschreiben das Geschlecht der drei Söhne des Patriarchen Noah, wie von ihm alle Menschen nach der Sündfluth auf Erden kommen sind; und geschiehet darum, daß die Historie bleibe in der Ordnung, daß man wisse, woher es komme und wohin es gehe; darum aber allermeist, daß man sehe, wie Gottes Wort allezeit wahrhaftig bleibt. Denn wir haben gehört, wie Gott durch Noah verfluchte Canaan, den Sohn Ham, aber Sem und Japhet segnete. Das erklärt nun dies Capitel weiter; und wiewohl es schlecht und einfältig ist, so ist es doch mächtig noch denen, so in der Schrift studiren. Es hanget viel aus den Propheten darinnen, da diese Namen geschrieben stehen. Darum müssen wir es auch nicht gar in Wind schlagen und hinwerfen.

2. Aber ein wenig davon zu fassen ist für die Laien genug, daß man ein wenig rohe Erkenntniß davon nehme, nicht so scharf ausgekehrt. So fähet er nun an, von dem jüngsten Sohne zu schreiben, was für Leute von ihm gekommen sind, und beschließet, daß von ihm ausgebreitet sind die Inseln der Heiden in ihren Ländern etc. Willst du nun wissen, wo sie hin gekommen sind, so theile die Welt gegen Mitternacht, Mittag und Morgen.

3. Japhet ist mit seinen Geschlechte hier herein gerathen gegen Mitternacht. Hispanien, Frankreich, Welisch- und Teutschland, Polen, Ungarn, Moscoviter und die Türkei, mit dem ganzen Meere und Inseln, und sind hier ein Theil Namen ausgedrückt der Söhne Japhets, welcher Namen auch die Länder behalten haben, als nämlich:

4. Gomer, der erste Sohn, davon drei Söhne gekommen, Askenas, Riphath und Thogarma. Gomer legt St. Hieronymus aus, daß es die Galater sollen sein. Aber es reimt sich nicht wohl. Wir halten, daß es sind die Riphei, die im Lande und Gebirge der Tartaren wohnen. Madai sind die Meder, ein groß Volk und Land, den Juden gegen der Mitternacht gelegen. Javan, sind eigentlich Jonen, das ist, die Griechen, davon Eliza, die Aeoles. Tharsis, die Thraces, Moscoviter, und da des Türken Sitz ist, Constantinopel. Kittim halte ich für das Land Macedonia, da der große König Alexander regiert hat [1 Macc. 1, 1.], wiewohl sie alle schreiben, es sei Welischland. Kurz\unsicher{…}, der Javan ist ein Vater aller Griechen, daher ihre Sprache auch gekommen ist. Askenas hält man für unsern, der Deutschen, Vater. So geht nun die Schrift herein, treibt den Sohn und sein Geschlecht nicht viel; aber im Propheten Ezechiel [Cap. 38, 2.] wird Gog und Magog vornehmlich angezogen und getrieben, so stark, daß ich gerne wollte den Türken dazu ziehen. [Es] ist aber ungewiß und nicht wohl zu errathen.

5. Ham aber, der andere Sohn, hat das beste Land gekriegt gegen Mittag, Syrian, Judäam, Arabiam, Egyptum, Lybian, Mauritaniam etc. bis an der Welt Ende, alles was gegen Mittag liegt; [er] hat, halte ich, weiter um sich gegriffen und regiert denn alle beide, und sich eben gesetzt in die rechte Schmalzgrube. Wir
\quellseite{107}
sind herein in die kalten, unfruchtbaren Länder gekommen. Dazu ist er so weit hinein gerissen, daß er in Chaldäa und Babylonia weit gegen dem Morgen gegriffen hat. Und meinen viele, daß, weil er ein Mann gewesen ist, geschickt und gerüstet in die Welt, er habe die Buchstaben erfunden; item, wie man streiten und schießen sollte. Das ist wahr, daß er der erste gewesen ist, davon die Weltherren gekommen sind.

6. Nun nimmt Moses vor sich seine Kinder und Kindeskinder. Der erste Sohn heißt Chus, der ist der Mohren Vater. Daher das Wort Chus noch einen Mohren heißt. Der hat sich wohl besamet und ausgebreitet, hat fünf Söhne gezeugt, die haben sich gesetzt in das Land reich Arabien. Mohrenland hat an sich das rothe Meer, und stößt an Egypten. Reich Arabien ist \unsicher{das erhinter}, gegen dem Morgen. Es soll auch noch eine Stadt sein im selbigen Lande, mit Namen Saba, nach dem einen Sohne Chus, davon die Sabäi genannt sind. Darnach, die andern Söhne sind alle daselbst hingekommen in die Grenze an Mohrenland, gegen Mittag.

7. Aber sonderlich zeigt der Text hier an, wie derselbe erste Sohn Chus habe auch gezeugt den Nimrod. Dieser hat nicht das Mohrenland eingenommen, wie die andern; und steht darauf, er sei ein Hurenkind gewesen; die geraten gemeiniglich also. Nimrod heißt auf deutsch ein Abtrünniger; als, wenn einer von seinem Herrn, dem er geschworen und gehuldiget hat, abfällt, und den geschwornen Eid nicht hält. Den Namen muß dieser Nimrod führen, daß er abtrünnig ist, und ein Vater aller, die in der Welt mit Gewalt regieren.

8. Siehe nun, wie der Text die Worte setzt. „Der fing an (spricht er), ein gewaltiger Herr zu sein auf Erden, und war ein gewaltiger Jäger“, ja, „vor dem Herrn“, und kann ein Sprichwort aus von ihm, nicht unter den Heiden, sondern unter den Gläubigen (denn Semes Geschlecht war fromm, darum wird es von diesem Nimrod), daß man sagt von einem Tyrannen: Siehe, da haben wir abermal einen Nimrod, oder einen gewaltigen Jäger. Darum heißt das\footnote{Jener: der.}, „Jäger“, der die Leute unterdrückt und zwingt; also, daß Nimrod mit Gewalt ein Herr geworden ist, und Leute unter sich gebracht mit Zwang, wie noch in allen Fürstenthümern geschiehet. Denn es ist kein Fürst, der sich an dem Seinen\footnote{In den alten Ausgaben: an seinem.} genügen lasse. So hat auch dieser von Anfange immer um sich gefressen, und ein groß Regiment unter sich gebracht. Und ist ohne Zweifel vor der Welt ein trefflicher Mann gewesen, der das weltliche Regiment in köstlichen Schwang gebracht und gehalten hat; sonst hätte ihn die Schrift nicht so fleißig angezeigt, und gesagt, daß er ein starker Jäger gewesen sei vor dem Herrn. Wenn wir jetzt in der Welt solche Fürsten sähen regieren, sollten wir meinen, wir hätten ein gülden Jahr.

9. Denn, wie er hernach sagt, er hat eine treffliche, große Hauptstadt aufgerichtet, die man hieß Babylonia; die hat er in vier Stück getheilt, wie man vor Zeiten Tetrapoles hieß, als Hebron und Rom auch war. Die waren so geordnet, daß der erste Theil waren patricii, Hersten und Gewaltigsten; der andere equites, die Reisigen oder der Adel, die bereit mußten sein zum Kriege; darnach plebs, der gemeine Mann; zum vierten rus, das ist, die Acker- und Bauleute. Darum muß es gar ein feines, ehrbares und tapferes Regiment sein gewesen, daß die Leute Lust und Liebe dazu haben gehabt, sonst hätte er nicht solche Städte können bauen, und so viel Volks zu sich bringen. Es sind dazumal nicht Narren gewesen, haben nicht viel Bücher gehabt, dennoch sein Regiment eingesetzt und gehalten, daß jedermann hat müssen loben.

10. Warum heißt ihn denn die Schrift einen starken Jäger? Darum, daß wiewohl er ein köstliches Regiment führte vor der Welt, so ist doch nie seines so köstlich gewesen, es hat Geiz und Ehrgeizigkeit bei sich, will immer weiter um sich fressen. Also hat er alle seine Gedanken und Begierden darauf gehabt, daß er ein großer, mächtiger Herr würde, viel Land und Leute zu regieren und herrschen hätte. Darum ist er ein Exempel aller Weltfürsten.

11. Warum hat nun Gott das Regiment so bestätigt, daß es dennoch aufgekommen und erhalten ist? Es hat Gott noch nie zu keinem gesagt: Nimm das Land oder Königreich ein, ohne allein zu den Juden. Er hat aber eingesetzt der Eltern Gewalt, welche die größte und beste ist über die Kinder und Gesinde, die ist unserm ersten Vater Adam befohlen mit ausgedrückten Worten. Darnach hat er es wieder durch Mosen
\quellseite{108}
[2 Mos. 20, 12.] geboten: Du sollst Vater und Mutter ehren. Darum ist diese\footnote{Wittenberger und Erlanger: die.} Gewalt größer, denn aller Könige oder Kaiser Gewalt, ja, die nächste nach GOtt. Daher man liest im Alten Testamente [2 Mos. 21, 7.], daß ein Vater Macht hatte, sein Kind zu verkaufen.

12. Diese Gewalt ist nun darum gegeben und eingesetzt, daß man die Kinder ziehen soll und Gottes Wort lehren, GOtt erkennen, fürchten und ihm glauben, also, daß ein Vater eigentlich ein Bischof und Pfarrherr seines Hauses sein soll. Denn ihm eben das Amt gebührt über seine Kinder und Gesinde, das einem Bischof gebührt über sein Volk. Dazu hat GOtt den Vätern [5 Mos. 21, 18—21.], über das vornehmste Amt, auch die Gewalt des weltlichen Schwerters gegeben, daß er möchte, wie gesagt, sein Kind verkaufen, ja, auch tödten, wenn er nicht recht gehorchen\footnote{„gehorchen“, welches in den Ausgaben fehlt, haben wir mit Walch nach 5 Mos. 21, 18. eingefügt.} wollte, daß ihm beide, geistliche und leibliche Gewalt, gegeben ist.

13. Was ist es denn, daß GOtt das weltliche Schwert und Gewalt hat geordnet und geboten, daß man ihr unterthan sei? Das ist [es]: Als die Kinder den Eltern nicht wollten gehorsam sein, wie sich’s noch wohl begibt, daß ein Vater ein ungerathenes Kind hat, das er nicht bezwingen kann, wird zu muthwillig und läuft von den Eltern; item, daß einem die Eltern nun gestorben sind, der frei und ruchlos nach seinem Willen leben will und niemand gehorchen, hat GOtt dennoch die Welt nicht wollen so unordig, ohne Zwang und Oberhand bleiben lassen; darum hat er das Schwert lassen aufkommen, daß man die Buben strafte; sonst dürfte man sein nirgend zu. Wo nun jene Gewalt nichts schaffen noch wehren kann, so hat GOtt diese Gewalt geordnet und bestätiget, und hat Gefallen daran. Wenn jene im Schwange bliebe, dürfte man dieser nicht; weil sie aber zu schwach ist, muß man jene haben, auf daß der Henker zwinge, die sich von ihren Eltern nicht ziehen lassen. So ist es ein nöthig Ding auf Erden. Das sind nun die Jäger, die solche Gewalt haben; aber die ihrer allewege übel gebrauchen, die sind des Nimrods Nachfolger. Köstlich ist die Gewalt an ihr selbst; aber nicht die Räuberei, dadurch sie mißbraucht wird.

14. Spricht nun Moses weiter: „Der An-

fang seines Reiches war Babel, Erech, Acad und Chalne im Lande Sinear.“ Dies sind vier Namen, und ist alles Eine Stadt. Der Hauptname aber ist Babel, die hoch berühmt in der Schrift ist, hat lange das Kaiserthum gehabt im Lande Sinear, das ist, Chaldäa. Item: „Davon ist darnach kommen der Assur“, das ist das andere Königreich, das man nennt das assyrische, von dem babylonischen kommen. Wie es immerdar geht, daß GOtt eines niederwirft und stürzet, und ein anderes aufrichtet. Durch dasselbige Reich ist hernach das jüdische Volk und Reich gefangen, davon die Propheten viel sagen. Dieser Assur hat auch gebauet die große Stadt Ninive und der Stadt Gassen, den Markt und Hauptplatz. Davon lies den Propheten Ezechiel, und andere.

\begin{quote}
B. 13. 14. Mizraim zeugete Ludim, Anamim, Leabim, Naphtuhim, Pathrusim, und Casluhim. Von dannen sind kommen die Philistim und Caphtorim.
\end{quote}

15. Dies ist der andere Sohn von Ham, der hat endlich sechse Söhne, darnach noch zwei. Von diesem hat das Egyptenland den Namen, daß es überall Mizraim heißt, und hat sehr viel Volks gehabt. Ist ein Zeichen, daß er das Land African habe eingenommen, ein groß Theil der Welt. Das ist das andere Volk. Nun kommt erst der rechte Strich, da die Macht anliegt, von dem dritten Sohne.

\begin{quote}
B. 15—19. Canaan aber zeugete Zidon, seinen ersten Sohn, und Heth, Jebusi, Emori, Girgosi, Hivi, Arki, Sini, Arvadi, Zemari und Hamathi.
\end{quote}

Daher sind ausgebreitet die Geschlechter der Canaaniter, und ihre Grenzen waren von Zidon an, durch Gerar, bis gen Gaza, bis man kommt gen Sodoma, Gomorra, Adama, Zeboim, und bis gen Lasa.

16. Das sind elf Söhne, die das gelobte Land inne gehabt haben, aber hernach durch die Juden vertrieben und geschlagen sind. Darum sagt er: „Daher sind ausgebreitet die Geschlechter der Canaaniter“; als wollte er sagen: Dies ist alle das Volk, da wir uns mit schlagen.

\begin{quote}
B. 20. Das sind die Kinder Ham, in ihren Geschlechten, Sprachen, Ländern und Leuten.
\end{quote}

17. So haben wir zwei große Theile der Welt, gegen Mitternacht und Mittage. Der
\quellseite{109}
dritte Sohn aber ist weit eher hin gekommen, gegen dem Morgen; da wissen wir wenig von. [Es] hat mich oft bekümmert, wo doch das Volk hinkommen und blieben sei; [ich] denke auch, ob wohl die meisten Apostel daselbst hinkommen seien? Wir haben einen Apostel der Heiden gehabt, auf Japheth's Seite, nämlich Paulum; sonst ist niemand hieherwärts gekommen. Was Petrus gethan hat, weiß ich nicht, daß mich dünkt, daß die Prophezei Noahs von Japheth niemand erfüllt habe denn Paulus. Denn er hat das gebracht, daß wir den Gott Sem erkannt haben. Von dem haben wir gewisse Schrift, daß er hier herein gekommen ist; von den andern wissen wir nichts; [es] ist aber zu denken, daß sie eherhin in Sems Land sind gekommen. Denn Sem ist hochgebenedeiet, und ein mächtig Volk worden, da wir nichts von wissen; denn das wir wissen, ist die Welt kaum die Hälfte.

18. Diese haben den größten Raum der Welt inne. Ob sie Christen sind, ist mir unbewußt. Ich halte es aber frei dafür, daß auch Christen unter ihnen sind, denn ich halte mich der Schrift, die je den Sem segenet, und ihm zusagt den Samen; und Christus ist über die ganze Welt ein Herr. Darum wird er ebensowohl Christen unter ihnen haben, als bei uns, und vielleicht bessere. So erzählt er nun des Sems Geschlecht.

\begin{quote}
V. 21—32. Sem aber, der Bruder Japheth's, des größten, zeugete auch Kinder, der ein Vater ist aller Kinder von Eber. Und dies sind seine Kinder: Elam, Assur, Arphachsad, Lud und Aram. Die Kinder aber von Aram sind diese: Uz, Hul, Gether und Mas. Arphachsad aber zeugete Salah, Salah zeugete Eber. Eber zeugete zwei Söhne, einer hieß Peleg, darum, daß zu seiner Zeit die Welt zertheilet ward, des Bruder hieß Joktan. Und Joktan zeugete Almodad, Saleph, Hazarmaveth, Jarah, Hadoram, Usal, Dikla, Obal, Abimael, Seba, Ophir, Hevila und Jobab. Das sind alle Kinder von Joktan, und ihre Wohnung war von Mesa an, bis man kommt gen Sephara, an den Berg gegen dem Morgen. Das sind die Kinder von Sem, in ihren Geschlechten, Sprachen, Ländern und Leuten. Das sind nun die Nachkommen\footnote{Wittenberger: Geschlechte.} der Kinder Noah in ihren Geschlechtern und Leuten, von denen sind ausgebreitet die Leute auf Erden nach der Sündfluth.
\end{quote}

19. Dieser Patriarch Sem hat lange gelebet, bei fünfhundert Jahre hinan [1 Mos. 11, 11.], daß er den Patriarchen Jakob erlebt und gesehen hat. [Es] sind auch viele, die da meinen, er sei der Melchisedech, davon hernach [1 Mos. 14, 18.] folgt; ich glaube es aber nicht. Summa Summarum: die Schrift spricht hier, er sei ein Vater aller Kinder Eber, welchen er hier sonderlich nennt, darum, daß die Geschichte unter ihnen geschah, die im folgenden Capitel steht; da wurden die Sprachen zertheilt, aber die rechte Gebräuche blieb bei denen, so sich haben zu dem gehalten. Darum heißt er der Kinder Eber Vater.

20. Seiner Söhne heißt einer Elam, davon [die] Elamiter genannt, das sind die Perser. Der andere Assur, das sind die Assyrier. Darnach Aram, davon machen sie Syriam, ist aber noch nicht gewiß. Die andern kennen wir nicht.

21. Von Arphachad ist kommen Eber, der hat zwei Söhne gehabt, spricht Moses, und hieß einen Peleg, das ist, Zertheilung. Denn da Eber sahe, daß, die den Thurm Babel baueten, der Sprache irre wurden, und eben zu der Zeit den Sohn kriegte, gab er ihm davon den Namen zum Gedächtniß.

22. So sind die Kinder Sem auch beschrieben. Aber da ist kein Regiment noch weltliche Gewalt; sie haben noch den Vater Sem, den lassen sie sich regieren. Die andern richten weltliche Herrschaft auf, daß sie andere Leute unter sich zwingen und weit regieren.

23. Da ist nun in dem Capitel beschrieben, wo die Welt hergekommen ist nach der Sündfluth. Aber mit dem Sem hat die Schrift wenig zu schaffen, desgleichen auch mit Japheth; aber mit Ham, und was von ihm gekommen ist, hat sie am meisten zu thun. Denn es gehen fast alle Historien damit um, wie wir hören werden.
\quellseite{110}

\section{Das elfte Capitel}

V. 1—9. Es hatte aber alle Welt einerlei Zunge und Sprache. Da sie nun zogen gen Morgen, fanden sie einen Plan im Lande Sinear, und wohneten daselbst. Und sprachen unter einander: Wohlauf, lasset uns Ziegel streichen und brennen. Und nahmen Ziegel zu Stein, und Thon zu Kalk, und sprachen: Wohlauf, lasset uns eine Stadt und Thurm bauen, deß Spitze bis an den Himmel reiche, daß wir uns einen Namen machen; denn wir werden vielleicht zerstreuet in alle Länder. Da stieg der HErr hernieder, daß er sähe die Stadt und Thurm, den die Menschenkinder baueten. Und der HErr sprach: Siehe, es ist einerlei Volk und einerlei Sprache unter ihnen allen, und haben das angefangen zu thun, sie werden nicht ablassen von allem, das sie vorgenommen haben zu thun. Wohlauf, lasset uns hernieder steigen, und ihre Sprache daselbst verwirren, daß keiner des andern Sprache vernehme. Also zerstreuete sie der HErr von dannen in alle Länder, daß sie aufhöreten die Stadt zu bauen. Daher heißt ihr Name Babel, daß der HErr daselbst verwirret hatte aller Länder Sprache, und sie zerstreuet von dannen in alle Länder.

1. Hierinne beschließet Moses abermal ein sonderlich Stück göttliches Werkes, das er gethan hat, und noch auf den heutigen Tag bleibet, nämlich, daß er die Sprache verwirrt, verändert und gemehret hat. Es ist auch ein groß Wunderzeichen, daß alle samt einerlei Sprache gewesen ist, und sich so weit getheilet hat. Es haben auch von der Zertheilung der Sprachen viele geschrieben, und sagen gemeiniglich, daß zwei und siebenzig Sprachen überall sind in der Welt. Diejenige Zahl, halte ich, daß sie genommen sei von der Zahl der Kinder der drei Söhne Noahs, Ham, Japheth und Sem; derer sind bisher fast zwei und siebenzig erzählt.

2. Dies hat nun keinen Grund, wir können es auch nicht wissen, wie mancherlei Sprachen auf Erden sind. Denn die Söhne Canaan, welcher doch elf gewesen sind, davon elf Völker mit Unterschied der Regenten gekommen sind, alle fast einerlei Sprache geredet haben, ohne daß ein wenig verändert mag gewesen sein; wie bei uns die deutsche Sprache von andern anders geredet wird, daß sich’s nahe bei hundertmal verändert. Danach weiter, sind die arabische, syrische, medianische und chaldäische Sprachen einander fast nahe. Als, da Abraham von Chaldäa zog, hat er ja die Leute verstanden, wo er hin kam, und ist dennoch ein groß Land durchzogen. Das rede ich darum, daß man bei dem Texte bleibe und die Schrift rein lasse, nicht allemal eine Glosse gebe, gleich als müßten es gerade zwei und siebenzig Sprachen sein, weil man so viel Kinder von Noah zählt. Das wissen wir, daß ihrer viel sind; wie viel ihrer aber ist, können wir nicht sagen.

3. Item, von dem Thurm Babel, davon Moises so schreibt, daß sie wollten bauen eine Stadt und Thurm, deß Spitze bis an [den] Himmel reichete, leget man auch getrost, und hat Ursache genommen aus den Worten so zu dichten, daß bereits der Thurm so hoch aufgeführet sei worden, daß man konnte die Engel singen hören, daß er schon gegangen wäre über die Wolken. Item, so lügen sie weiter, es stehe noch heutiges Tages ein Stück davon; aber also, daß er drei Meilen in die Erde sei gesunken, und drei Meilen oben hinweg gewehet vom Winde, und stehe noch drei Meilen hoch davon. Das sind alles Lügenhebendinge.

4. Es ist eine wohlgebaute Stadt gewesen, wie eine andere, ohne daß sie sich wohl unterstanden haben, etwas Hohes und Großes aufzurichten, wo ihnen Gott nicht gewehret hätte; und hat es freilich der Nimrod angefangen, davon oben [Cap. 10, § 9] gesagt ist. So verstehe nun die Worte Mosis nicht anders, denn daß er so viel hat wollen sagen, daß sie haben wollen eine hohe und mächtige Stadt bauen. Denn der Schrift Weise ist, also zu reden, wie im fünften Buche Mosis [Cap. 1, 28.]: „Unsere Brüder haben gesagt, das Volk sei größer und höher denn wir, die Städte sind groß und bis in den Himmel vermauert.“ Item, abermal im selbigen Buche [Cap. 4, 11.] steht also: „Und ihr tratet herzu, und stundet unten an dem Berge, der Berg brannte aber bis mitten an den Himmel.“ Also heißt nach der Schrift alles Himmel,
\quellseite{111}
was über und unter uns ist, alles, das blau sieht, nahe und fern; als auch Christus im Matthäo [Cap. 6, 26.] redet: „Sehet die Vögel des Himmels an“ 2c. Darum ist es hier auch so viel geredet; sie haben die Spitze des Thurms wollen hoch aufführen, bis zum blauen Himmel. Denn daß sie in Himmel sollte gereicht haben, das war unmöglich.

5. Das sei darum gesagt, daß man nicht Ursache nehme, solche großen Lügen zu machen, wie man hier gethan hat, als auch das ist, daß sie den Thurm darum sollten gebaut haben, ob einmal wieder eine Sündfluth käme, daß sie davor bleiben könnten, so doch Moses hier klar anzeigt, wie sie gesagt haben: „Lasset uns einen Namen machen, denn wir werden vielleicht zerstreuet in alle Lande.“ Sie wollten ihr Datum auf das Gebäu setzen, daß man darnach davon in aller Welt zu sagen und singen müßte; wollten Ehre und Ruhm davon tragen, daß alle Leute sagen sollten: Dort ist eine Stadt und Gebäu, desgleichen nie gesehen noch gehört ist. Das wird hernach zur heimlichen Deutung dienen.

6. Nun lasset uns sehen, was er mit dieser Historie hat wollen anzeigen. Es ist ein Exempel, darinne man sehe, was die Welt ist. Droben [Cap. 6, § 3 f.] hat er beschrieben Gottes Kinder; hier beschreibt er die Menschenkinder. Denn das Buch ist darum geschrieben, daß man beiderlei Volk lerne kennen. So finden wir allemal Gottes Kinder so beschrieben, daß sie der Welt mit gebraucht haben, aber sich nicht darinne gewühlt, noch um sich gegriffen zu herrschen, als wollten sie ewig darinne bleiben; jene aber haben alle ihr Datum, Trotz und Trost darauf gesetzt. Diese haben nicht am zeitlichen Gut genug gehabt, sondern am geistlichen; jene haben allein auf das Zeitliche getrachtet, daß sie nur Gut und Ehre hätten.

7. So findet sich nun der Vorwitz, der viel Leuten anhängt, zu bauen; wiewohl das noch gering ist, daß man den Vorwitz büßt. Das ist aber hier gestraft, wie der Text spricht, daß sie solches vorgenommen hatten, und wollten nicht davon lassen; das ist, sie waren ersoffen im zeitlichen Gut, daß sie ihre Zuversicht und Ruhm setzten in solches Gebäue; suchten nicht Gott. Solchen Hochmuth und Trotz wollte er nicht leiden, darum stieg er herab, und stürzte alles.

8. So sagt nun Moses: „Da sie zogen gegen dem Morgen, fanden sie einen Plan im Lande Sinear, und wohneten daselbst“ 2c. Sinear ist gewesen ein Platz oder flaches Feld, darauf sie die Stadt bauten. Denn so sagen die Schreiber, daß sie liege in einem köstlichen Lande und freiem, ebenem Felde, daß man sie überall weit sehen kann; hat auch viel mächtige Könige gehabt, als erstlich den Nimrod, item, die Semiramis 2c., daß man so viel Wunders davon schreibt, daß ich nicht weiß, ob es alles zu glauben ist oder nicht. Es ist ein Anfang des weltlichen Reichs gewesen, und hat länger gestanden, denn keine Stadt. Es sagt auch Aristoteles, es sei nicht eine Stadt, sondern ein Land mit Mauern umringt. Da haben sie nun lernen Ziegel streichen, haben aber keinen Kalk gehabt, sondern nur Thon, wie man auch noch sagt, daß man in dem Lande nur mit Leim maure.

9. Was ist aber das, daß er sagt: „Da stieg der Herr hernieder, daß er sähe die Stadt und Thurm.“ Item, und der Herr sprach: „Wohlauf, lasset uns hernieder steigen, und ihre Sprache daselbst verwirren“, und redet dazu von sich, als von vielen. Das ist je nicht von Einer Person zu deuten, wiewohl es die Juden so deuten, Gott habe es zu den Engeln gesagt; ist aber nicht genugsam. Es ist freilich davon geredet, daß der Herr, davon Moses sagt, sei hernieder gestiegen; denn er zuvor von Einem Herrn geredet hatte. Derselbige Einige spricht nun: „Lasset uns hernieder steigen.“ Also hat er von Anbeginn schon angezeigt den hohen Artikel, wie auch vor [Cap. 1, §§ 13. 53] gehört ist, und hier und hernach mehr, daß Ein Gott sei, und doch mehr denn Eine Person.

10. Warum mußte aber der Herr herab steigen, hätte er doch wohl mögen droben bleiben? Sagt er doch durch den Propheten Jeremiam [Cap. 23, 23.]: Non Deus a longe, ich bin nicht ein Gott von ferne; item, V. 24.: Implebo coelum et arida, ich erfülle Himmel und Erde. Auch sagen wir im Glauben: Ich glaube an Gott, Schöpffer Himmels und Erden. Hat er Himmel und Erde geschaffen, so muß er je an allen Orten sein, da sein Geschöpf ist. Ja, spricht er im 99. Psalm, V. 1.: Qui sedet super Cherub, et intuetur abyssos, er sitzet auf dem Cherub, und siehet alle Tiefe, daß man bekennen muß, wo die Creatur ist, daß da auch Gott ist. Denn keiner das Wesen von ihm selbst hat, müssen alle erhalten werden von Gottes Gewalt
\quellseite{112}
. Wie geht es denn zu, so er schon hernieden ist, daß er sagt, er wolle herab steigen?

11. Antwort: Also haben wir droben auch gehört im 3. Capitel, V. 9., wie der HErr Adam suchte im Garten, und sprach: „Wo bist du?“ Gerade als wüßte er nicht, wo er wäre. Item, am vierten, V. 9., sprach er zu Kain: „Wo ist dein Bruder Abel?“ Auch kommt er herab zu Noah und redet mit ihm [Cap. 6, 13.]. Also redet die Schrift allenthalben, daß Gott herab vom Himmel steige, welches doch nicht sein kann. Ist aber eben das, so ich oben auch gesagt habe, Gott steigt also hernieder, wenn er sich offenbart, entweder inwendig im Herzen, oder äußerlich durch ein Werk; wie im Evangelio Matthäi, Cap. 3, 16. 17., der Heilige Geist in einer leiblichen Gestalt auf Christum stieg, als er getauft ward, und eine Stimme vom Himmel fiel, das ist, Gott der Vater; dazu auch der Heilige Geist offenbart sich durch ein äußerlich Werk, damit er sich zuvor nicht offenbart hatte, wiewohl er zuvor überall ist. Also ist er auch Mensch worden und auf Erden kommen, nicht daß er den Himmel verlassen habe, sondern hat sich da in einem sonderlichen Werke der Welt offenbart. So thut er noch, wenn er uns heimsucht mit Gutem, oder straft. Das heißt alles: zu uns kommen; ist doch nichts Anderes, denn daß er sich offenbart. Das ist eine Weise.

12. Die andere ist, daß der Mensch auch die Augen aufthue, und solches fühle. Denn Gott kommt zu vielen, die ihn doch nicht sehen, noch sein gewahr werden. Als, da sich Christus offenbarte, waren ihrer wenig, die ihn erkannt haben; die ihn aber erkannten, haben gesehen, daß Gott da herab kommen war, die andern aber meinten, er wäre ein Mensch, wie ein anderer. Dies Stück ist nun höher und tiefer denn das erste. Denn ich habe oft gesagt, wie sich das Gewissen gegen Gott hält, also ist er. Hältst du, daß er gnädig sei, so ist er gnädig; fürchtest du dich vor ihm, als vor einem schrecklichen Richter, so ist er es auch, richtet dich immerdar nach deinem Gewissen. Wie nun solch Gewissen sich ändert gegen Gott, so ändert sich auch die Sprache der Schrift, die redet davon, gerade wie man ihn fühlt. Also sagt David oft im Psalter: „HErr, stehe auf, warum schläfst du, warum zürnest du über uns?“ etc. [Ps. 7, 7. 68, 2. 44, 24.] Also gibt sie Gott zu: sehen, hören, gehen, stehen, reden, schweigen, schlafen und wachen, daß wir es also empfinden, und weil es sich so wandelt zwischen mir und ihm, wiewohl sonst Gottes Natur unverwandelt bleibt. [Es] ist aber alles um des Glaubens willen zu thun.

13. Also verstehe nun auch, daß Gott hier hernieder steigt durch ein äußerlich Werk, nämlich, daß er die Sprache zertrennt, daß keiner den andern vernehmen kann; nicht allein also, daß solche greuliche Plage und Werk äußerlich geschieht, sondern auch, daß ihr Gewissen fühlt, daß ihnen Gott ungnädig sei, und ihr Bauen, dadurch sie Ruhm suchen, verdammt. Wie sie nun fühlen, also ist es gewesen; darum sind sie auch durch solch Schrecken des Gewissens und äußerlich Verwirren alle zertrennt, und von einander getheilt in die Welt. Das Herz ward verzagt, und außen Irrung, darum konnten sie nicht bei einander bleiben. Denn die Einigkeit der Sprache ist darum gegeben, daß man sich freundlich zusammen halte. So ist es nun auch oben mit Adam und Kain zugegangen, wie gesagt. Darum ist das alles nicht mehr, denn daß Gott durch äußerliche Werke das Herz rührt; wie das fühlt, so heißt denn Gott schlafen, aufstehen, wirken, reden etc.

14. Wie aber, daß der HErr nicht frisch zufährt, ihre Sprache trennt, und sie irre macht, sondern nimmt einen Rath darüber, und beschließt es zuvor bei sich selbst? Was ist ihm solches vonnöthen, daß er sich erstlich stellt, als sähe er es nicht? Antwort: Der Spruch ist mit Fleiß gesetzt wider das frevel Urtheil. Daran sollten sich stoßen alle Afterredner und Ehrenbläser. Wiewohl Gott alle Dinge wußte und sah, dennoch stieg er herab, zu sehen, was sie machten. Darum soll man kein Urtheil fällen, wie gewiß man auch der Sache ist. Man soll zuvor darinne und davon handeln, sich aller Dinge erkundigen, daß das Urtheil nicht fehlen kann.

15. Dies ist ein nöthig Stück, und wohl zu merken. Denn das ist nun das dritte Mal angezeigt, erstlich in Adam, darnach mit Kain. Also werden wir hernach abermal hören von Sodom und Gomorra, da er [Cap. 18, 20. 21.] spricht: Das Geschrei Sodom und Gomorra ist viel worden. Ich will hernieder steigen, und sehen, ob sie nach dem Geschrei, das vor mich kommen ist, gethan haben oder nicht, auf daß ich es wisse.

16. Und ist darum nöthig zu merken, denn
\quellseite{113}
der Teufel hat immer die Leute besessen mit giftigen Zungen, daß einer dem andern nachredet, und doch so scheinbarlich vorträgt, wie denn ein jeglicher seine Sache schmückt und die beste macht, daß, die es hören, oft gefangen werden, und sagen, es ist vor Gott nicht recht. Solchem zu wehren, hat Gott dies Exempel lassen schreiben, daß man sich vorsehe, niemand verurtheile, sondern zuvor den andern Theil auch höre. Darum soll man keinem Menschen nimmer so bald glauben, wenn er von einem andern redet in seinem Abwesen, wie heilig er auch sein mag, also, daß man sich ja hüte vor plötzlichem Urtheilen. Willst du je glauben, magst du es thun; aber halte nur das Maul inne, und richte nicht; gehe nicht hin und sage, es sei also, wie du gehört hast; schleuß nicht ein Urtheil, du siehest es denn vor selbst.

17. Also soll man kein Ding ansehen auf eines andern Rede. Denn es gilt nicht, daß du dich auf einen andern berufest, und sprechest: Der oder dieser hat es gesagt. Thust du es aber, so fehlst du gewißlich, wirst bald verführt und zum Narren gemacht. O was Jammers und Unglücks möchte nachbleiben, wenn man solches nicht thäte! So gehen wir hin und bauen auf Menschen, so man doch weiß, daß er leugt, wie die Schrift [Ps. 116, 11.] sagt. Gott will selbst nicht urtheilen, er habe denn zuvor gefragt: „Wo bist du?“ Darum laß zuvor hören, was der Selbstschuldige sagt. Wenn er nun öffentlich überwunden ist, so magst du urtheilen. Wollte Gott, daß wir das in das Herz bildeten, und darnach richteten, daß wir die Zunge dämpften, die alles Herzeleid anrichtet, daß schwerlich immer zu rathen ist, auch durch wahrhaftige Zungen.

18. So haben wir nun gehört, was uns die Historie dieses Capitels von dem Thurm Babel lehrt. Ist nun auch zu sagen von der geistlichen Deutung, so darinne angezeigt ist. Ich habe also gesagt, daß in diesem Werke die Menschenkinder nichts so sehr versehen haben, als daß sie wollten einen Namen haben. Das wird sonderlich dazu gesetzt, und ihr Vornehmen eben darum getadelt und gestraft. Es ist aber damit angezeigt der große Haufe und das gemeine Uebel der Welt, das da heißt falsche Prediger. Denn da legt Gott am meisten an, daß sein heiliges Wort bleibe einig in der Christenheit, ohne allen Zusatz der Menschenlehre. Aber das Wort kann nicht bleiben, es sei denn, daß man Christum allein habe für den Bauherrn, und lasse ihn den Namen allein haben. Wo der nicht ist, da ist keine Einigkeit, und muß gewiß auch ein Babel daraus werden; nämlich also:

19. Wenn man von Christo predigt, so gibt man einem Menschen so viel, als dem andern, daß keiner mehr von Christo hat, denn der andere; St. Peter nicht mehr, denn ich, Maria nicht mehr, denn ein ander ehelich Weib; denn es ist nicht mehr, denn ein einiger Christus, der wird einem jeglichen ganz gegeben. Doch ist es auch wahr, daß ihn einer besser fassen kann, denn der andere; doch ist es Ein Sach. Als, wenn einer ein Kleinod fasset in Gold, der andere in ein schlecht Gefäß, das Gefäß mag besser und baß zugerichtet sein, aber das darein gefasset ist, wird nicht besser noch edeler: also mag einer sonderliche, höhere Gaben haben denn der andere, dennoch bleibt Christus gemein; und wie der Mann gemein bleibt, also auch alles das Gut, das er mit sich bringt. Also hat ein jung Kind gerade so viel als ein Alter, ein Gelehrter nicht mehr denn ein Laie, ein Herr nicht mehr denn ein Knecht. Diese Namen, Taufe, Evangelium und Geist, sind alle gemein.

20. So lange nun diese Lehre bleibt, so lange bleibt Einigkeit in der Welt. Denn da muß ein jeglicher sagen: Ich habe nicht mehr, denn der Allergeringste; wir sind alle gleich im Erbe; einer ist des andern Bruder; wir haben alle gleich Recht dazu. Wenn aber andere Bauleute kommen und aufstreten, machen sie so bald aus der Einigkeit eine Zerspaltung und Zertrennung. Wie es geht im päbstlichen, und auch gehen muß in allem weltlichen Regimente. Da regieret man also, daß ein Volk mehr hat, denn das andere. Ja, ein Fürst muß also regieren, daß es ungleich zugehe.

21. Also auch in dem Papstthum, welches man geistlich nennt, da ist immer ein Orden nach dem andern aufgekommen, allein darum, daß immer einer hat wollen besser sein denn der andere. Denn sie je öffentlich sagen: Ja, sollten wir nicht mehr haben denn andere, was richten wir denn für Wesen an? So hängen sie allewege einen Zusatz daran, daß es soll etwas Besonderes sein; so ist es schon geschieden und getrennt, und nimmer gleich; dazu kommt eitel Neid und Haß unter ihnen daraus.

22. Sollten sie aber solches nicht thun, und sagen: Ich bin nichts besser, denn der geringste
\quellseite{114}
Christ, wir haben alle gleich Einen Christum, Taufe, Evangelium und Glauben [Eph. 4, 5. 6.]: so müßten sie bekennen, daß ihr Ding Narrenwerk wäre, so würde es von sich selbst dahin fallen. Darum kann ihr Wesen nicht bestehen, es richte denn Zwiespältigkeit auf, und trete von der Einigkeit. So tritt immer einer nach dem andern auf, und spricht: Mein Stand ist der beste. [Es] nimmt ein jeglicher etwas vor sich, das etwas Köstlicheres sein solle, daher er die andern verachtet. Wo nun solche Lehre und Wesen geht und getrieben wird, so ist das ausgetilgt, daß wir alle gleich seien. Darum leidet sich das Evangelium nicht mit Menschenlehren und Secten. Darum ist solches alles überall in der Schrift gewehrt; noch hilft es nicht, man hat es lassen liegen, als ginge es uns nichts an, so es doch allein ein lebendig Buch ist, voll lebendiger Worte [Joh. 6, 68. Hebr. 4, 12.].

23. Also folgt es nun nach einander. Wenn man abtritt, und einen sonderlichen Stand aufrichtet, muß man bald auch einen andern Namen aufwerfen, daß der Name mannigfaltig werde, wie das Wesen mannigfaltig und getrennt ist. Darum sprechen sie: Ja, sollte ich nicht einen andern und bessern Namen haben, denn sie alle haben? Wir wollen heißen Benedicter, Francisc(er), Carthäuser etc., den wollen wir führen über Christi Namen. Das sind die, so hier bedeutet, die den Thurm bauen, die ihnen wollten Namen machen, daß sie es gethan haben. Es ist je eine große Blindheit und Vermessenheit, ja, Sünde und Schande, daß man solche Prediger leiden und hören soll, die unverschämt sagen: Wer St. Franciscus nachfolgt, der ist in einem bessern Leben und höhern Stande, denn sonst ein Christ, als sollte Franciscus heiliger sein denn Christus, so er doch Mühe und Arbeit gehabt hat, daß er Christo nachfolgte.

24. So haben wir die Deutung dieser Historie, daß die Bauleute sind Prediger. Die nun nicht im Glauben sind, die bauen und predigen wohl etwas Großes, Sonderliches vor andern, daß sie großen Namen haben; aber es wird die Plage und Strafe folgen, die hier angezeigt ist, daß sie uneinig, in ihrer Sprache verwirrt und zertrennt werden, so weit die Welt ist. Das wirst du so finden in allen Secten. Wo wahrhaftige Christen sind, die sind aller Dinge einträchtig, und kann sich keiner über den andern erheben noch sich zertrennen, da muß die Sprache auf eine Weise gehen, und gar übereinstimmen. Was du redest und predigst, das predige und rede ich auch; und wiederum. So bleibt man eins, inwendig im Herzen, und äußerlich im Leben.

25. Dort aber geht es so zu, daß sie alle haben gesagt: Man muß bleiben unter Einem Haupte, das ist der Pabst, das ist noch einig. Aber daraus fließen so viel Secten, als Klöster und Bisthümer sind, da ist keines mit dem andern eins, will jegliches das beste sein, und muß eines dem andern Feind sein, und es verachten. Dahin gehören auch die hohen Schulen mit ihren mancherlei Secten, die ihnen folgen. Da verstehet keiner des andern Sprache, keines Predigt gefällt dem andern; ist noch nie gehört, daß ein Predigerordens Mönch\footnote{So die Jenaer. Wittenberger: eines andern Ordens Mönch. Erlanger: ein ander Prediger.} aufstände, der da sollte den Francisc(er)orden hoch heben, sondern ein jeglicher Mönch preist seinen Orden.\footnote{Die eingeklammerten Worte fehlen in der Erlanger, stehen aber in der Wittenberger und in der Jenaer.} So hört auch keiner gerne, was der andere aufwirft und preist, oder wiederum.

26. Das ist die verwirrte und vermengte Sprache. Daraus folgt auch ein zertrenntes Wesen, daß sich keiner des andern annimmt, keiner dem andern Liebe erzeigt. [Es] ist so weit gekommen, daß ein Pfarrherr mit dem andern, ja, ein Altarist mit dem andern nicht eins ist. Darum auch kein ungeschichter Volk auf Erden ist denn diese Leute; ihrer keiner dient jemand, denn ihm selbst. Die Christenheit aber ist also gethan, daß sich ein jeglicher des andern Noth annimmt, und für ihn sorgt und ihm dient. Darum heißt auch solch uneiniges Wesen recht eine Babel, das ist, eine Vermengung oder Verwirrung. Es ist nichts mit der Babylonia gewesen, gegen diese. Hier treibt es Gott viel stärker, daß er sie trennt und zerstreut, so hetzt und treibt von einander, daß keiner dem andern Gutes gönnt. Das ist eine greuliche Strafe und Plage.

27. Es haben auch die Propheten in dies Capitel gesehen, und viel Prophezei daraus gezogen; als im 5. Psalm, V. 11.: „Schuldige sie, Gott, daß sie fallen von ihrem Vornehmen“ oder Ratschlägen, „stoße sie aus“ etc. Da rührt er diese Historie fein, als sollte er sagen: Sie hatten auch vorgeschlagen, einen Thurm zu bauen und hoch zu führen, was ist aber geschehen?
\quellseite{115}
Sie mußten abstehen und mit Schanden davon lassen. Item, also sagt er auch in einem andern Psalm: Herr, verschlinge und zertrenne ihre Zungen [Ps. 55, 10.].

28. Solche Gebete sollten mir auch thun wider unsere Babel, und sagen: Lieber Herr, gib Gnade wider ihre Rotten und Secten, daß sie davon fallen; trenne und theile sie, mache sie nur uneinig, daß sie zu Boden gehen, darum, daß sie so mannigfaltig dich mit ihren Secten erzürnen, und verderben dein einig Volk, Lehre und Namen. Solcher Prophezey sind viel mehr in der Schrift, die sich hieher ziehen, wer darauf Achtung hat.

29. Was beweget sie aber, daß sie so feindlich bauen an dem Ort, auf dem Platze Sinear, da sie ein fein fett Feld fanden? Es ist nimmer vergessen in der Schrift, daß sie ihnen immerdar Schuld gibt, daß sie Bauchdiener sind und geizige Wänste. Die zwei Stücke findet man gemeiniglich an ihnen, Geiz und Bauch. Das ist eben die Ursache, darum sie alles ihr Thun vornehmen. So sagt Paulus von ihnen Phil. 3, 19.: Quorum Deus venter est, Bauch ist ihr Gott; und Col. 3, 5.: Avaritia idolorum cultus, wer geizig ist, der ist abgöttisch. Das haben sie alle an sich, glauben an Gott nicht, darum sorgen und fürchten sie, sie können sich nicht ernähren. Darum brauchen sie Zungen, Hände und Werke dahin, daß sie die Leute schinden und schatzen, auf daß sie ihre Bäuche zuvor füllen und genug haben, ehe sie Gott dienen. Und das ist die Summa davon, daß sie den Thurm bauen auf den besten Platz, da eine feiste, volle Schmalzgrube ist, da wollen sie wohnen und bleiben.

\begin{quote}
V. 10—26. Dies sind die Geburt Sem.
\end{quote}

Sem war hundert Jahr alt, und zeugete Arphachsad zwei Jahr nach der Sündfluth, und lebte darnach fünfhundert Jahr, und zeugete Söhne und Töchter. Arphachsad war fünf und dreißig Jahr alt, und zeugete Salah, und lebte darnach vierhundert und drei Jahr, und zeugete Söhne und Töchter. Salah war dreißig Jahr alt, und zeugete Eber, und lebte darnach vierhundert und drei Jahr, und zeugete Söhne und Töchter. Eber war vier und dreißig Jahr alt, und zeugete Peleg, und lebte darnach vierhundert und dreißig Jahr, und zeugete Söhne und Töchter. Peleg war dreißig Jahr alt, und zeugete Regu, und lebte darnach zweihundert und neun Jahr, und zeugete Söhne und Töchter. Regu war zwei und dreißig Jahr alt, und zeugete Serug, und lebte darnach zweihundert und sieben Jahr, und zeugete Söhne und Töchter. Serug war dreißig Jahr alt, und zeugete Nahor, und lebte darnach zweihundert Jahr, und zeugete Söhne und Töchter. Nahor war neun und zwanzig Jahr alt, und zeugete Tharah, und lebte darnach hundert und neunzehn Jahr, und zeugete Söhne und Töchter. Tharah war siebenzig Jahr alt, und zeugete Abram, Nahor und Haran.

30. Da beschreibt Moses zehn Glied nach einander, von Sem bis auf Abraham, die Linie, davon Christus kommen sollte. Das sind eitel große, heilige Leute. Wie fein wird die Welt da gestanden sein, da jeglicher so lange gelebt hat, und so viel Kinder und Kindeskinder gesehen hat, dazu alles regiert ist durch väterliche Obrigkeit. Denn hier siehst du kein Weltregiment, darum ist es auf dieser Seite tödtlich gestanden. Der heilige Vater Noah hat noch den Abraham erlebt; Sem, sein Sohn, hat auch bei zehn Kindeskinder nach ihm gesehen, und unter ihnen regiert, daß wenn man die Jahre rechnet, findet sich’s, daß er bis an Jakob, ja, schier bis an Joseph, in das elfte Glied gelebt habe.

31. So haben die Väter alle an einander gehalten in einem feinen geistlichen Regiment. Moses beschreibt es aber darum, daß er die Linie Christi will nach einander führen; auch daß er den Ehestand bei Ehren behielte, denn er darf es auch wohl, sonderlich jetzund. Gott hat genug zu schaffen, daß er ihn bei Ehren erhalte wider unsere Geistlichen, die ihn nicht allein verachtet und gering gehalten, sondern auch gelästert haben. Ich halte, es sei noch ihrer keiner, die den Jungfrauenstand gehalten haben, so heilig als der Erzväter einer. Denn sie sind nicht so für schlechte Leute anzusehen, weil sie Moses durch den Heiligen Geist hervor zeucht. Hat es nun denen wohl angestanden, daß sie Weiber gehabt und Kinder gezeugt haben, so mögen wir es mit allen Ehren auch thun.

\begin{quote}
V. 27—32. Dies sind die Geburt Tharah.
\end{quote}

Tharah zeugete Abram, Nahor und Haran. Aber Haran zeugete Lot. Haran aber starb vor seinem Vater Tharah in seinem Vaterlande, zu Ur in Chaldäa. Da nahmen Abram und Nahor
Weiber\footnote{Wittenberger und Erlanger: oder.}.
\quellseite{116}

Abrams Weib heißt Sarai, und Nahors Weib Milca, Harans Tochter, der ein Vater war der Milca und der Jisca. Aber Sarai war unfruchtbar, und hatte kein Kind. Da nahm Tharah seinen Sohn Abram, und Lot, seines Sohnes Harans Sohn, und seine Schnur Sarai, seines Sohnes Abrams Weib, und führte sie von Ur aus Chaldäa, daß er in das Land Canaan zöge. Und sie kamen gen Haran, und wohneten daselbst. Und Tharah ward zweihundert und fünf Jahr alt, und starb in Haran.

32. Hier führt er nun das Geschlecht Tharahs, des Vaters Abrams, und zeucht an, daß er drei Söhne gehabt habe, welches man auch muß um folgender Historien willen wissen. Der Rechnung nach muß Abram der jüngste Sohn gewesen sein, wiewohl ihn Moses für den ersten angezählt. Das kann ich daher ermessen, daß er überlebt den Haran, seinen Bruder, und desselben zwei Töchter und einen Sohn. Denn Lot ist Abrams Bruder-Sohn gewesen; dessen zwei Schwestern waren Sarai und Milca. So haben die zwei Brüder, Abram und Nahor, zwei Schwestern genommen, ihres leiblichen Bruders Töchter.

33. Aus diesem kann man nun schließen, ist auch nie aufgehoben, und kann es noch niemand mehren, daß einer mit Gott und Ehren möge seines Bruders Töchter nehmen, weil es Abram und Nahor gethan haben, welche in der Schrift gepreiset werden, und wir sollen glauben, daß sie heiliger gewesen sind denn wir. Wäre es nicht recht noch göttlich gewesen, würde es weder Abram noch Sarai gethan haben. Doch ist dagegen verboten, daß einer seine Muhme nicht nehmen darf, wiewohl es ein gleichem Glied ist.

34. Von Abrams Bruder-Sohn, Lot, werden wir hernach hören; dazu auch von den beiden Schwestern. Denn Milca wird sich auch befreunden, und eine Mutter werden, 1 Mos. 22, 20. 21. Von den andern sagen die Juden, ob es wahr ist, weiß ich nicht gewiß, daß Jisca gewesen sei eben die Sarai, Abrams Weib, und zwei Namen gehabt habe. Wo dem also ist, so gibt die Rechnung, daß Abram viel jünger gewesen sei denn Haran. Denn man liest klar hernach also, daß er hundert Jahre alt gewesen ist, als ihm Isaak geboren ward, und Sarai, sein Weib, neunzig Jahr. Das ist so viel geredet, Abram ist zehn Jahre älter gewesen, denn sein Weib.

35. Soll nun Haran der Sarai Vater sein, und Abram älter sein, so müßte er die Tochter gezeuget haben im achten oder neunten Jahr, und seinem Bruder gegeben im zehnten Jahr. Darum sagen wir also, daß nicht gewiß ist, daß Sarai und Jisca Ein Weib sei. Ist es aber Ein Weib, so muß Abram der jüngste Sohn sein, und wird ihm doch die Ehre gethan, daß ihn Moses zum ersten nennt, darum, daß Gott etwas Sonderliches aus ihm machen wollte.

36. So sagt nun der Text: Tharah mit seinen drei Söhnen hat gewohnt in Chaldäa, das ist, im selben Lande, da sie den Thurm Babel bauten; denn Babylonia liegt im Lande Chaldäa. Da ist eine Stadt gewesen mit Namen Ur, darinne haben sie gewohnt. Aber der eine Bruder, Haran, ist darinne gestorben; Tharah aber wollte mit den andern zweien Söhnen in das Land Canaan ziehen, und gerieth, daß sie kamen und sich niedersetzten in Haran, das ist, im Lande Syria, welches liegt gegen dem gelobten Lande gegen Morgen, und oft angezogen wird in der Schrift.

37. Hier machen die Juden wunderlich Ding aus, [ich] weiß nicht, ob es alles zu glauben ist. Meinen also, daß Haran gemartert sein solle von den Chaldäern darum, daß er nicht wollte das Feuer anbeten; darum sei auch Tharah mit den andern Söhnen geflohen. Ich aber wollte gern den Abram zuvor zu einem Sünder machen, ehe ich ihn fromm machte; und das darum, denn also liest man in dem Buch Josua [Cap. 24, 2. 3.], daß er sagt zu dem Volk: „Jenseit des Wassers haben gewohnt eure Väter, Tharah, der Vater Abrahams und Nahors, und haben fremden Göttern gedienet, darum habe ich euren Vater Abraham genommen aus Mesopotamia, und gebracht in das Land Canaan“ etc., also, daß Abraham auch abgöttisch muß gewesen sein, auf daß man nicht den Patriarchen erhebe seiner Heiligkeit und freien Willens halben. Er ist eben so tief im Schlamm gelegen als wir; daß er aber zu einem solchen Manne geworden ist, hat er wohl gelernt, daß es aus lauterem Beruf, Gnade und Güte geschehen ist, dadurch ihn Gott heraus geführt hat in das Land Canaan. Derhalben kann ich nicht wohl glauben, daß sie darum sollten geflohen sein aus Chaldäa, daß sie nicht müßten Abgötter anbeten, es wäre denn ein anderer Abgott gewesen denn das Feuer.
\quellseite{117}

38. Das Wörtlein Ur habe ich lassen bleiben, wie es steht, und halte es mit denen, die da sagen, es sei ein Name einer Stadt in Chaldäa oder Babylonia. Etliche wollen, es heiße Feuer, darum, daß es die Babylonier für ihren Gott also geheißen und geehrt haben, davon etliche heidnische Bücher sagen. Also ist es noch nicht gar gewiß. Ur heißt zwar nicht eigentlich Feuer, sondern Licht, wiewohl es fast eines ist; daher, meine ich, haben sie der Stadt einen Namen gegeben. Desgleichen wir noch etliche bei uns Deutschen genannt haben, als Lichtenstein, Lichtenfels, Lichtenberg etc.

39. Ich halte es nicht dafür, daß die Babylonier daselbst so grobe Leute gewesen sind, daß sie nichts von Gott sollten gewußt haben, der alle Dinge geschaffen hat, oder, daß sie das Feuer für einen wahrhaftigen Gott gehalten haben; sondern das mag sein, daß sie Gott, der Himmel und Erde geschaffen hat, mit dem Namen genannt und geehrt haben. Wie auch die Juden Gott viel Namen gaben, als Baal, das ist, einen Mann, darum, daß sie ihn wollten so nahe sein als eine Braut oder ehelich Weib ihrem Manne. Irem, Moloch, das ist, ein König, und andere mehr; gingen immerdar damit um, daß sie den rechten Gott wollten also ehren.

40. Also haben wir auch gethan, und unserm Gott mancherlei Namen gemacht; als Augustinum, Benedictum, Franciscum etc. Denn wie wir von Gott halten, so ist er uns, so hat er auch einen Namen. Wenn ich ihn also ansehe, und von ihm denke, daß er mich durch meinen geistlichen Stand selig mache, so habe ich mir schon einen Abgott gemacht, und meine, der wahrhaftige Gott, der Himmel und Erde geschaffen hat, sei so ein Mann, wie ich ihn achte, und solle mich ansehen, daß ich einen solchen Gehorsam halte, bete und faste. So male ich mir ihn ab, und streiche ihm eine Farbe an, als habe er eine Carthäuser- oder Augustinerkappe an, und denke: Ei, wie fein Ding ist es um die Kappe. Also zeucht ihm ein jeglicher, der ohne Glauben ist, ein eigen Kleid an, macht einen Potzmann\footnote{Potzmann = Puppe, Popanz, Vogelscheuche.} aus ihm, gibt ihm im Herzen einen Namen und Gestalt, wie ihn gut dünkt.

41. Wenn aber der Glaube kommt, zeucht er ihm solch Kleid oder Farbe ab, und sagt: Der Gott sieht weder diesen noch jenen an, sondern allein den Herrn Jesum Christum; wer an das liebe Kind glaubt, der hat den rechten Gott, und sieht ihn, wie er ist.

42. Also haben nun diese auch gethan. Sind aber noch nicht so große Narren gewesen als wir, haben so klug und weise wollen sein, daß man nicht sollte sagen, sie beteten Holz und Steine an; dennoch nehmen sie Gott den rechten Namen und Ehre, ziehen ihn zu sich, und geben ihm einen andern; den muß er behalten in der Welt, so muß er sich noch das ganze Jahr in allen Predigten führen lassen in der Nummerei. So haben sie nun gelehrt, wo man das Licht sehe, da sollte man Gott anbeten; wie man uns gepredigt hat, wo ein Crucifix oder ander Bild stände, da sollte\footnote{In den Ausgaben: soll.} man Christum anbeten. So tolle Narren sind sie noch nicht, daß sie meinen, daß dies Licht oder Bild Gott sei; aber solche Narren sind sie mit uns, daß sie Gott da anbeten, und meinen, er habe Lust dazu und sehe dahin. Und lehren also: Wo du vor ein Licht kommst und anbetest, so triffst du Gott, das wird ihm gefallen. Es ist auch ein feiner, scheinbarlicher Gottesdienst gewesen, welchen niemand hätte können umstoßen, denn der Geist, wie auch zu unsern Zeiten. Wer wäre so klug gewesen, daß er hätte sagen dürfen, solches wäre unrecht, das wir getrieben haben mit unserm Gottesdienst, Messen und geistlichem Leben? Der Geist aber kommt und sagt: Gott hat es nicht befohlen, ist wider den Glauben, darum ist es unrecht und abgöttisch.

43. Also haben auch dazumal die heiligen Väter gethan, und darüber gepredigt: Nicht also, Gott hat es nicht geboten, ihr habt kein Zeugniß davon durch Gottes Wort, thut es aus eurem eigenen Kopf und Vermessenheit, daß ihr Gott also bindet an solche Ort- und euer Werk. Also sollt ihr aber thun: Gott hat uns verheißen [1 Mos. 3, 15.] einen Samen, der der Schlange Kopf zertreten soll;\footnote{Wittenberger und Jenser: sollt.} da sollt ihr an hangen und an ihn glauben. Aber solches haben sie nicht wollen hören, haben die Patriarchen für Ketzer gescholten und mit dem Tode gebräuet. Daher man auch sagt, daß sie die Leute mit Feuer gebrannt haben, wie jetzund, in der letzten
\quellseite{118}
Babylon, unsere Geistlichen auch thun. Denn dies Regiment, wie man schreibt, ist gern mit Feuer umgegangen.

44. So liest man von den Dreien, in den feurigen Ofen gelegt, im Propheten Daniel [Cap. 3, 22.]. Darum haben etliche geschrieben, wie gesagt, daß Haran, Abrams Bruder, in dem Lande gestorben sei um des Glaubens willen, und mit Feuer verbrannt; item, daß Abram darum geflohen sei, daß er dem Feuer entfliehe. Wie nun dem sei, haben wir keinen Text. So viel haben wir aber [Jos. 24, 3.], daß Abram auch abgöttisch gewesen sei, und Haran dazu, sammt dem dritten Bruder. Darnach ist GOtt gekommen, und hat den Abram durch sonderliche Gnade erleuchtet, auf daß er sähe, wie allein der Glaube vor GOtt gilt auf den künftigen Samen; da kam er in rechten Verstand, was der rechte Gottesdienst wäre.

45. Darum steht ein feiner Spruch in dem Propheten Jesaia [Cap. 50, 2. 40, 25. 46, 5. 9.]: Siehe, ich bin ein solcher GOtt, der das Meer trocken macht, und ist kein anderer, denn ich; wenn vergleichet ihr mich denn, dem ich sollte ähnlich sein? Als sollte er sagen: So große Gewalt habe ich, daß ich alle Dinge in meiner Hand habe, noch seid ihr so thöricht, daß ihr mich fassen wollt in ein klein Werk, und ich bin so groß, daß mich das ganze Meer nicht fassen kann. Was ist das anders, denn mit unsern Träumen umgehen und die anbeten, GOtt so enge spannen, wie wir gedenken, den man sollte heben über alle Creaturen, welches allein der Glaube thut?

46. Das habe ich darum geredet, daß man Mosen recht verstehe, der da beschreibt, wie Abram heraus gerissen ist von GOtt aus dem abgöttischen Wesen, aber noch nicht berufen, wie hernach davon im folgenden Capitel steht. Darum will ich ihn einen armen Sünder lassen bleiben, so lange bis GOtt kommt, und heißt ihn [1 Mos. 12, 1.] gehen aus seines Vaters Hause. Das geschieht alles uns zu Trost, daß wir sehen, wie GOtt wenig Leute gen Himmel hebt, die rein sind, sondern werden gemeiniglich alle aus dem Schlamm gehoben.

47. Aufs letzte hebt sich eine große Frage über dem Text, wenn man kommt in das zwölfte Capitel: wie sich zusammen reime, daß hier steht: Tharah ward zweihundert und fünfzig Jahr alt, und starb in Haran; und darauf [Cap. 12, 4.] folgt, wie Abram ist aus dem Lande gezogen; als habe er seines Vaters Tod erwartet, ehe er ist ausgegangen? Es ist aber leicht zu berichten, also, daß dies gesagt sei per anticipationem, und zuvor geschrieben wird, wie alt sein Vater geworden ist. Nun ist oben [B. 26.] gesagt, daß Tharah über siebenzig Jahre war, da er die drei Söhne zeugete; und folgt hernach [Cap. 12, 4.], daß Abram fünf und siebenzig Jahre alt war, da er in das Land Canaan kam; das sind noch nicht zweihundert Jahr; darum muß er je ein Jahr oder dreißig vor seines Vaters Tode ausgezogen sein.

\section{Das zwölfste Capitel}

\begin{quote}
V. 1—3. Und der HErr sprach zu Abram: Gehe aus deinem Vaterlande, und von deiner Freundschaft, und aus deines Vaters Hause, in ein Land, das ich dir zeigen will. Und ich will dich zum großen Volk machen, und will dich segnen, und dir einen großen Namen machen, und sollst ein Segen sein. Ich will segnen, die dich segnen, und verfluchen, die dich verfluchen. Und in dir sollen gesegnet werden alle Geschlechte auf Erden.
\end{quote}

1. Hier lasset uns sehen ein groß, mächtig, trefflich Exempel des Glaubens. Es ist auch fast der edelsten Capitel eines in diesem Buch; darum muß man es nicht schlecht ansehen, und überhin laufen, als gelte es jetzt nicht mehr. Ich habe dies Buch darum vorgenommen, daß nicht noch sei, Fabeln und Märlein zu predigen von der Heiligen Legenden, deren man alle Welt voll erdichtet hat, sondern daß man sehe rechtschaffene Exempel des Glaubens und der Liebe, auf daß man nicht allein mit Geboten und Gesetzen umgehe, sondern lebendige Worte und Weise vor Augen habe, wie es im Schwange geht, wenn man glaubt und liebt.
\quellseite{119}

2. Also haben wir auch oben gehört ein groß Exempel des Glaubens in dem heiligen Vater Noah; dagegen ein Exempel des Unglaubens von denen, die den Thurm Babel bauten, daß man allenthalben sehe, wie es um den Glauben zu thun ist. Wie wir nun jetzt gesehen haben, was des Unglaubens Art und Natur ist in den Ungläubigen, nämlich, wie sie ihre Sinne darauf haben, daß sie hier genug hätten, Gut, Ehre und Lust erlangeten, und ihre Herzen nicht richten auf das unsichtige, künftige Gut, darum sie auch sagten: Wir wollen uns einen Namen machen; was war das anders, denn daß sie gedachten zeitliche Ehre und Preis zu erlangen, daß sie nicht geringe und verachtete Leute wären? Wollten sich also einsehen, als sollten sie ewig hier bleiben. Also wollen wir hier wiederum sehen, was der Glaube schafft und für Kraft hat.

3. Ueber das ist daneben ein groß Stück, daß hier wieder eine Zuneigung gethan wird von Christo, daß er sollte von Abram kommen. Denn bisher haben sich die Väter alle mit der ersten Verheißung beholfen, die Gott Adam gethan hatte; das ist ihr Evangelium gewesen, das sie getrieben haben, bis es Gott hier wieder aufrichtet, heller und deutlicher macht, denn es nun schier verfallen war. Es sind kurze Worte, haben aber mächtig groß Ding hinter sich.

4. So ist nun das ein neu Evangelium, das Abram holt aus der Abgötterei; denn er muß vor fünf und siebenzig Jahre ein Ungchrist sein, und ein tapferer, starker Sünder, ehe er heilig wird, auf daß allenthalben die Schrift allein preise die Gnade und Barmherzigkeit. Da ist kein Verdienst; [es] hebt dürre an, wie Gott Abram ruft; schreibt nicht, daß er je ein gut Werk gethan habe, geopfert, oder Altar gebauet, wie hernach; sondern als er noch steckt in Sünden, und nie gedacht, daß er solche Gnade überkommen sollte, kommt es unversehens.

5. Darum habe ich gesagt: Soll jemand bekehrt, fromm, oder ein Christ werden, so heben wir nicht an; [es] dient kein Beten noch Fasten dazu, es muß vom Himmel und allein aus Gnaden kommen, daß Gott durch die Verheißung des Evangelii das Herz trifft, daß es fühlt, und müsse sagen, daß es vor nie bedacht oder in Sinn genommen habe, daß ihm solche Gnade sollte widerfahren. Wie hier: Ehe sich Abram hätte dürfen unterstehen zu bitten, ja, ehe er sich gedenkt zu bekehren, kommt Gott zuvor, nimmt ihn aus dem Irrthum, und setzt ihn in ein ander Wesen. Darum laß dir dies Exempel eben so neu sein, als wäre es gestern geschehen, denn es ist das Werk Gottes, das noch immerdar geht. Also muß es täglich gehen, wenn man fromm werden; was man anders lehrt, das gilt nichts. Wer da fromm will werden, der sage nur nicht: Ich will anheben und gute Werke thun, daß ich Gnade erlange; sondern also: Ich will warten, ob Gott durch sein Wort mir seine Gnade und Geist wollt geben. Das muß es thun, sonst ist es alles verloren; wie der Prophet in 85. Psalm, V. 9., sagt: Audiam, quid loquatur in me Dominus, ich will hören, was Gott der Herr redet.

6. So ist nun in diesen Worten beides verfasset, das Gesetz und Evangelium. Zum Gesetz gehört, daß er sagt: „Gehe aus deinem Vaterlande, und von deiner Freundschaft, und aus deines Vaters Hause, in ein Land, das ich dir zeigen will.“ Das ist so viel gesagt: Alle dein Thun und Wesen gefällt mir nicht, ich will dich nicht hier haben, du mußt heraus und alles verlassen. Das hat ihm auch wehe gethan, ist ein harter Ruff gewesen. Darum sind ihrer auch nicht viel mit ihm gegangen, denn sein Weib, und Lot, sein Vetter. Da hat er müssen sagen und bekennen vor Vater, Mutter, Bruder und jedermann: Dies Wesen taugt nicht; wäre etwas Gutes hier, so würde mich Gott nicht heißen heraus treten.

7. Nun sind diese Leute nicht Narren gewesen, sondern die besten, die man finden soll. Ist eben gewiß, als wenn man jetzt zu einem Carthäuser sagte: tritt aus deinem Wesen und Orden. So schwer es einem solchen würde, so schwer ist es Abram auch worden. Solches richtet niemand aus, denn die geistliche Predigt des Gesetzes, wenn es kommt und spricht: Was du thust, das ist nichts, du bist blind, siehst nicht, was gut ist, wie solltest du nun Gutes thun? So steckt es uns in die Sünde und Hölle hinein.

8. Darnach kommt die andere Predigt, das Evangelium, in diesen Worten: „Und ich will dich zum großen Volk machen und segnen“ etc., ja, sagt er, ich will dazu aus deinem Samen erwecken einen solchen Mann, der alle Welt wird segnen. Das\footnote{Erlanger: Da.} ist eitel Gnade, geistlich und weltlich, zeitlich und ewig. Das ist nun das
